% I :Produire un plan détaillé d’une dissertation — Français 1re Tech — Cours signé OnBuch+
% Programme officiel MINESEC. Compiler avec : tectonic N17-i-produire-plan-detaille-dissertation.tex
\newcommand{\DOCMATIERE}{Français}
\newcommand{\DOCNIVEAU}{1re Tech}
\newcommand{\DOCMODULE}{Français – classe de Première technique}
\newcommand{\DOCLECON}{N17}
\newcommand{\DOCTITRE}{I :Produire un plan détaillé d’une dissertation}
% OnBuch+ — préambule « cours signés » ENRICHI (compilé avec tectonic / XeLaTeX)
% Système visuel « Pop » d'OnBuch+ : fond crème (jamais blanc), bordures encre
% épaisses, ombres « dures » décalées, titres Archivo Black en capitales.
% Version MATHÉMATIQUES : graphes (pgfplots/tikz), boîtes pédagogiques
% (définition, propriété, méthode, attention, le savais-tu, exercice résolu,
% exercice, TP, bilan), page de garde et signature OnBuch+ ancrée en bas.
%
% Macros attendues dans main.tex AVANT \input{preamble} :
%   \DOCMATIERE  \DOCNIVEAU  \DOCMODULE  \DOCLECON (numéro)  \DOCTITRE
\documentclass[11pt]{article}

\usepackage{fontspec}
\usepackage[french]{babel}
\usepackage[a4paper,margin=2cm,top=3.4cm,bottom=2.8cm,headheight=1.4cm,footskip=1.3cm]{geometry}
\usepackage{amsmath,amssymb,amsfonts}
\usepackage{mathtools} % \xmapsto, \coloneqq... souvent utilisés par les modèles
\usepackage{cancel} % simplifications barrées (bilans de forces)
% \abs{z} (module, valeur absolue) et \norm{u} (norme), utilisés par les modèles
\providecommand{\abs}{}\let\abs\relax\DeclarePairedDelimiter{\abs}{\lvert}{\rvert}
\providecommand{\norm}{}\let\norm\relax\DeclarePairedDelimiter{\norm}{\lVert}{\rVert}
\usepackage[table]{xcolor} % [table] : fournit aussi \rowcolors (alternance de lignes)
\usepackage{pagecolor}
\usepackage{tikz}
\usetikzlibrary{arrows.meta,positioning,calc,shapes.geometric,shapes.misc,shapes.arrows,decorations.pathmorphing,decorations.pathreplacing,decorations.markings,patterns,shadows,mindmap,trees,fit,backgrounds,matrix,intersections,angles,quotes}
\usepackage{pgfplots}
\usepackage[version=4]{mhchem} % équations nucléaires \ce{^{235}_{92}U}
\usepackage[european,siunitx]{circuitikz} % schémas électriques
\pgfplotsset{compat=1.18}
\usepgfplotslibrary{fillbetween,groupplots} % aires entre courbes (calcul intégral)
\usepackage{enumitem}
\usepackage{fancyhdr}
% [most] charge toutes les bibliothèques tcolorbox (listings, minted, raster,
% documentation...) et gonfle inutilement la mémoire TeX sur un document long
% (~300 boîtes) : on ne charge que ce qui est réellement utilisé.
\usepackage[skins,breakable]{tcolorbox}
\usepackage{graphicx}
\usepackage{colortbl}
\usepackage{booktabs}
\usepackage{tabularx}
\usepackage{multirow}
\usepackage{diagbox} % case d'en-tête diagonale des tableaux à double entrée
% Colonnes extensibles centrée / gauche / droite (C, L, R), souvent utilisées par les modèles
\newcolumntype{C}{>{\centering\arraybackslash}X}
\newcolumntype{L}{>{\raggedright\arraybackslash}X}
\newcolumntype{R}{>{\raggedleft\arraybackslash}X}
\usepackage{array}
\usepackage{multicol}
\usepackage{siunitx}
% parse-numbers=false : accepte aussi \SI{2,5 \times 10^{-3}}{...}
\sisetup{locale=FR,output-decimal-marker={,},group-minimum-digits=4,parse-numbers=false}
\DeclareSIUnit\ohm{\ensuremath{\symup{\Omega}}} % U+2126 absent de la police
\DeclareSIUnit\division{div} % réglages d'oscilloscope (ms/div, V/div)
\DeclareSIUnit\tr{tr} % tours (vitesse de rotation, ex. \SI{1500}{\tr\per\minute})
\DeclareSIUnit\molar{M} % concentration molaire (ex. \SI{3}{\molar}, \SI{1}{\micro\molar})
\DeclareSIUnit\an{an} % année (le modèle confond parfois avec \year, un compteur TeX) — ex. \SI{5}{\kilo\meter\per\an}
\DeclareSIUnit\FCFA{FCFA} % franc CFA (ex. \SI{500}{\FCFA\per\kilo\gram})
\DeclareSIUnit\degreeBrix{\textdegree Brix} % teneur en sucre d'un sirop/jus (réfractométrie)
\DeclareSIUnit\month{mois} % le modèle utilise parfois \month, un compteur TeX, comme unité de durée
\DeclareSIUnit\microliter{\micro\liter} % le modèle écrit parfois \microliter en un seul token
\DeclareSIUnit\million{millions} % dénombrement (ex. \SI{15}{\million\per\mL} spermatozoïdes)
\usepackage{qrcode}
% \degree : commande fréquemment utilisée par les modèles pour un angle ou une
% température en dehors de \SI{}{} (non fournie par les paquets chargés).
\providecommand{\degree}{^{\circ}}
% \qed (fin de démonstration), utilisé par les modèles sans amsthm
\providecommand{\qed}{\hfill$\square$}
% \barre : double barre des valeurs interdites dans un tableau de variations (tkz-tab)
\providecommand{\barre}{\ensuremath{\|}}
% environnement proof (amsthm non chargé) utilisé par les modèles
\ifdefined\proof\else\newenvironment{proof}[1][Démonstration]{\par\noindent\textit{#1.}\ }{\qed\par}\fi
\usepackage{lastpage}
\usepackage{hyperref}
\hypersetup{hidelinks}

% ============================================================
% Palette « Pop » OnBuch+ — mêmes hex que la classe P (plus_ui.dart)
% ============================================================
\definecolor{poporange}{HTML}{F59321}
\definecolor{popdark}{HTML}{E07A0C}
\definecolor{popcream}{HTML}{FFF6E8}
\definecolor{popink}{HTML}{1C1714}
\definecolor{poppurple}{HTML}{7A5AE0}
\definecolor{popgreen}{HTML}{2E9E6A}
\definecolor{poppink}{HTML}{E0457B}
\definecolor{popblue}{HTML}{2F7DE1}
\definecolor{popgold}{HTML}{C98A12}
\definecolor{popmuted}{HTML}{8A7F76}
\definecolor{popline}{HTML}{EADFD2}
\definecolor{popgray}{HTML}{8A7F76}
\colorlet{popgrey}{popgray}
% Alias tolérants (le modèle nomme parfois les couleurs différemment)
\colorlet{popwhite}{white}
\colorlet{popblack}{popink}
\colorlet{popred}{poppink}
\colorlet{popyellow}{popgold}
% Teintes claires (fonds de tableaux/encadrés) — le modèle nomme parfois
% les couleurs avec un suffixe L (ex. popblueL) sans les avoir définies.
\colorlet{poporangeL}{poporange!20}
\colorlet{popdarkL}{popdark!20}
\colorlet{poppurpleL}{poppurple!20}
\colorlet{popgreenL}{popgreen!20}
\colorlet{poppinkL}{poppink!20}
\colorlet{popblueL}{popblue!20}
\colorlet{popgoldL}{popgold!20}
\colorlet{popredL}{popred!20}
\colorlet{popgrayL}{popgray!20}
% Teintes claires pour fonds de tableaux / zones
\colorlet{popblueL}{popblue!10!white}
\colorlet{poporangeL}{poporange!14!white}
\colorlet{popgreenL}{popgreen!12!white}
\colorlet{poppurpleL}{poppurple!12!white}
\colorlet{poppinkL}{poppink!12!white}
\colorlet{popgoldL}{popgold!14!white}

\pagecolor{popcream}

% ============================================================
% Typographie
% ============================================================
\defaultfontfeatures{Ligatures=TeX}
\setmainfont{PlusJakartaSans-Medium.ttf}[Path=../../../fonts/,BoldFont=PlusJakartaSans-Bold.ttf]
% Maths en sans-serif (Fira Math) avec chiffres et lettres droites de Plus
% Jakarta, pour que les formules se fondent dans le texte.
\usepackage[math-style=ISO,bold-style=ISO]{unicode-math}
\setmathfont{FiraMath-Regular.otf}[Path=../../../fonts/]
\setmathfont{PlusJakartaSans-Medium.ttf}[Path=../../../fonts/,range={up/num,up/{Latin,latin},\mathup}]
\usepackage{icomma} % virgule décimale sans espace en mode math
% Caractères mathématiques Unicode absents des polices de texte (Plus Jakarta,
% Archivo Black) : rendus par leur équivalent mathématique, en texte comme en
% formule — sinon ils s'affichent en carré vide (ex. « Arithmétique dans ℤ »).
\usepackage{newunicodechar}
\newunicodechar{ℝ}{\ensuremath{\mathbb{R}}}
\newunicodechar{ℤ}{\ensuremath{\mathbb{Z}}}
\newunicodechar{ℂ}{\ensuremath{\mathbb{C}}}
\newunicodechar{ℕ}{\ensuremath{\mathbb{N}}}
\newunicodechar{ℚ}{\ensuremath{\mathbb{Q}}}
\newunicodechar{𝔻}{\ensuremath{\mathbb{D}}}
\newunicodechar{α}{\ensuremath{\alpha}}
\newunicodechar{β}{\ensuremath{\beta}}
\newunicodechar{γ}{\ensuremath{\gamma}}
\newunicodechar{δ}{\ensuremath{\delta}}
\newunicodechar{ε}{\ensuremath{\varepsilon}}
\newunicodechar{θ}{\ensuremath{\theta}}
\newunicodechar{λ}{\ensuremath{\lambda}}
\newunicodechar{μ}{\ensuremath{\mu}}
\newunicodechar{σ}{\ensuremath{\sigma}}
\newunicodechar{τ}{\ensuremath{\tau}}
\newunicodechar{φ}{\ensuremath{\varphi}}
\newunicodechar{ψ}{\ensuremath{\psi}}
\newunicodechar{ω}{\ensuremath{\omega}}
\newunicodechar{Σ}{\ensuremath{\Sigma}}
% majuscules produites par \MakeUppercase dans les titres (α-aminés -> Α)
\newunicodechar{Α}{\ensuremath{\alpha}}
\newunicodechar{Β}{\ensuremath{\beta}}
\newunicodechar{Γ}{\ensuremath{\Gamma}}
\newunicodechar{Δ}{\ensuremath{\Delta}}
\newunicodechar{Ε}{\ensuremath{\varepsilon}}
\newunicodechar{Θ}{\ensuremath{\Theta}}
\newunicodechar{Λ}{\ensuremath{\Lambda}}
\newunicodechar{Μ}{\ensuremath{\mu}}
\newunicodechar{Τ}{\ensuremath{\tau}}
\newunicodechar{Φ}{\ensuremath{\Phi}}
\newunicodechar{Ψ}{\ensuremath{\Psi}}
\newunicodechar{Ω}{\ensuremath{\Omega}}
\newunicodechar{↦}{\ensuremath{\mapsto}}
\newunicodechar{∈}{\ensuremath{\in}}
\newunicodechar{∉}{\ensuremath{\notin}}
\newunicodechar{∧}{\ensuremath{\wedge}}
\newunicodechar{⇔}{\ensuremath{\Leftrightarrow}}
\newunicodechar{⟺}{\ensuremath{\Longleftrightarrow}}
\newunicodechar{⇒}{\ensuremath{\Rightarrow}}
\newunicodechar{⟹}{\ensuremath{\Longrightarrow}}
\newunicodechar{∘}{\ensuremath{\circ}}
\newunicodechar{∪}{\ensuremath{\cup}}
\newunicodechar{∩}{\ensuremath{\cap}}
\newunicodechar{≡}{\ensuremath{\equiv}}
\newunicodechar{⊂}{\ensuremath{\subset}}
\newunicodechar{⊥}{\ensuremath{\perp}}
\newunicodechar{∃}{\ensuremath{\exists}}
\newunicodechar{∀}{\ensuremath{\forall}}
\newunicodechar{∛}{\ensuremath{\sqrt[3]{\,}}}
\newunicodechar{∜}{\ensuremath{\sqrt[4]{\,}}}
\newunicodechar{⃗}{\ensuremath{{}^{\to}}}
\newunicodechar{ⁿ}{\ifmmode^{n}\else\textsuperscript{\ensuremath{n}}\fi}
\newunicodechar{ˣ}{\ifmmode^{x}\else\textsuperscript{\ensuremath{x}}\fi}
\newunicodechar{ᵇ}{\ifmmode^{b}\else\textsuperscript{\ensuremath{b}}\fi}
\newunicodechar{ʳ}{\ifmmode^{r}\else\textsuperscript{\ensuremath{r}}\fi}
\newunicodechar{ᵘ}{\ifmmode^{u}\else\textsuperscript{\ensuremath{u}}\fi}
\newunicodechar{ʸ}{\ifmmode^{y}\else\textsuperscript{\ensuremath{y}}\fi}
\newunicodechar{ᵃ}{\ifmmode^{a}\else\textsuperscript{\ensuremath{a}}\fi}
\newunicodechar{ᵉ}{\ifmmode^{e}\else\textsuperscript{\ensuremath{e}}\fi}
\newunicodechar{ᵏ}{\ifmmode^{k}\else\textsuperscript{\ensuremath{k}}\fi}
\newunicodechar{ᵖ}{\ifmmode^{p}\else\textsuperscript{\ensuremath{p}}\fi}
\newunicodechar{ᴺ}{\ifmmode^{N}\else\textsuperscript{\ensuremath{N}}\fi}
\newunicodechar{ᶜ}{\ifmmode^{c}\else\textsuperscript{\ensuremath{c}}\fi}
\newunicodechar{ʰ}{\ifmmode^{h}\else\textsuperscript{\ensuremath{h}}\fi}
\newunicodechar{ᵠ}{\ifmmode^{q}\else\textsuperscript{\ensuremath{q}}\fi}
\newunicodechar{ₙ}{\ifmmode_{n}\else\textsubscript{\ensuremath{n}}\fi}
\newunicodechar{ₐ}{\ifmmode_{a}\else\textsubscript{\ensuremath{a}}\fi}
\newunicodechar{ᵢ}{\ifmmode_{i}\else\textsubscript{\ensuremath{i}}\fi}
\newunicodechar{ᵣ}{\ifmmode_{r}\else\textsubscript{\ensuremath{r}}\fi}
\newunicodechar{ₖ}{\ifmmode_{k}\else\textsubscript{\ensuremath{k}}\fi}
\newunicodechar{ᵦ}{\ifmmode_{\beta}\else\textsubscript{\ensuremath{\beta}}\fi}
\newunicodechar{ₓ}{\ifmmode_{x}\else\textsubscript{\ensuremath{x}}\fi}
\newfontfamily\popdisplay{ArchivoBlack-Regular.ttf}[Path=../../../fonts/]
\newfontfamily\poplogo{SpaceGrotesk-Bold.ttf}[Path=../../../fonts/]
\newfontfamily\popsemi{PlusJakartaSans-SemiBold.ttf}[Path=../../../fonts/]
\newfontfamily\popxbold{PlusJakartaSans-ExtraBold.ttf}[Path=../../../fonts/]
\newfontfamily\popxboldls{PlusJakartaSans-ExtraBold.ttf}[Path=../../../fonts/,LetterSpace=3.0]

\setlist[enumerate]{leftmargin=1.7em,itemsep=5pt,topsep=4pt}
\setlist[itemize]{leftmargin=1.7em,itemsep=4pt,topsep=4pt,label={\color{poporange}\textbullet}}
\setlength{\parindent}{0pt}
\setlength{\parskip}{5pt}
\renewcommand{\arraystretch}{1.25}

% pgfplots : style Pop par défaut
\pgfplotsset{
  every axis/.append style={axis line style={popink,line width=1pt},tick style={popink},
    grid style={popline,line width=0.5pt},label style={font=\small},tick label style={font=\footnotesize}},
  popcurve/.style={poporange,line width=1.8pt,no markers,smooth},
}
% Même style disponible dans un \draw TikZ simple (hors pgfplots)
\tikzset{popcurve/.style={poporange,line width=1.8pt}}
\tikzset{popgrid/.style={popline,very thin}}
\colorlet{popcurve}{poporange} % le nom du style est parfois utilisé comme couleur

% ============================================================
% Pill / squiggle
% ============================================================
\newcommand{\poppill}[3]{%
  \tikz[baseline=-0.55ex]\node[draw=popink,line width=1.2pt,rounded corners=2.6pt,%
    fill=#2,inner xsep=6.5pt,inner ysep=3pt]{%
    \popxboldls\fontsize{7}{7}\selectfont\color{#3}\MakeUppercase{#1}};%
}
\newcommand{\popsquiggle}[1]{%
  \tikz[baseline=-0.4ex]{
    \draw[#1,line width=2.6pt,line cap=round]
      plot [smooth] coordinates {(0,0.09) (0.34,0.01) (0.68,0.21) (1.02,0.01) (1.36,0.10)};
  }%
}
% Mot-clé surligné (marqueur orange sous le texte)
\newcommand{\cle}[1]{{\popxbold\color{popdark}#1}}

% ============================================================
% En-tête / pied
% ============================================================
\pagestyle{fancy}
\fancyhf{}
\renewcommand{\headrulewidth}{0pt}
\renewcommand{\footrulewidth}{0pt}
\lhead{\raisebox{-0.15ex}{\poplogo\fontsize{13}{13}\selectfont\color{popink}OnBuch\color{poporange}+}}
\rhead{\poppill{\DOCMATIERE\ \textbullet\ \DOCNIVEAU\ \textbullet\ Leçon \DOCLECON}{white}{popink}}
\lfoot{\popxbold\fontsize{7.6}{7.6}\selectfont\color{popmuted}\MakeUppercase{Cours signé OnBuch+}\ \textbullet\ {\popsemi\DOCTITRE}}
\rfoot{\tikz[baseline=-0.6ex]\node[rounded corners=8.5pt,fill=popink,minimum height=17pt,minimum width=17pt,inner xsep=5pt,inner sep=0]{\popxbold\fontsize{8}{8}\selectfont\color{popcream}\thepage\,/\,\pageref*{LastPage}};}

% ============================================================
% Titres
% ============================================================
\newcommand{\chaptitre}[2]{%
  \begin{tcolorbox}[enhanced,colback=poporange,colframe=popink,boxrule=2.6pt,arc=11pt,
      drop shadow=popink,left=15pt,right=15pt,top=12pt,bottom=11pt,before skip=3pt,after skip=18pt]
    \poppill{#2}{popink}{popcream}\par\vspace{9pt}
    {\raggedright\hyphenpenalty=10000\exhyphenpenalty=10000\popdisplay\fontsize{22}{24}\selectfont\color{popcream}\MakeUppercase{#1}\par}
  \end{tcolorbox}
}
% Section numérotée : pastille orange avec le numéro + titre Archivo
\newcounter{popsec}
\newcommand{\coursec}[1]{%
  \stepcounter{popsec}\setcounter{popsub}{0}%
  \par\vspace{16pt}\noindent
  \begin{minipage}{\linewidth}
  \tikz[baseline=-0.7ex]\node[draw=popink,line width=1.6pt,fill=poporange,rounded corners=4pt,minimum size=22pt,inner sep=0]{\popdisplay\fontsize{12}{12}\selectfont\color{popcream}\thepopsec};%
  \hspace{8pt}{\popdisplay\fontsize{15}{17}\selectfont\color{popink}\MakeUppercase{#1}}\par
  \vspace{4pt}\noindent{\color{poporange}\rule{36pt}{3.6pt}}
  \end{minipage}\par\nopagebreak\vspace{8pt}
}
\newcounter{popsub}
\newcommand{\courssub}[1]{%
  \stepcounter{popsub}%
  \par\vspace{9pt}\noindent{\popxbold\fontsize{11.5}{13}\selectfont\color{popink}{\color{poporange}\thepopsec.\thepopsub}\hspace{6pt}#1}\par\nopagebreak\vspace{4pt}
}

% ============================================================
% Boîtes pédagogiques — même recette : fond blanc, bordure épaisse colorée,
% ombre dure encre, badge plein. Titre optionnel [#1].
% ============================================================
\tcbset{popbase/.style={breakable,enhanced,colback=white,boxrule=2.3pt,arc=9pt,
  shadow={3pt}{-3pt}{0pt}{popink},left=14pt,right=14pt,top=11pt,bottom=11pt,before skip=11pt,after skip=15pt,
  fontupper=\popsemi\fontsize{10.6}{14.5}\selectfont\color{popink},
  fontlower=\popsemi\fontsize{10.6}{14.5}\selectfont\color{popink}}}
% Coche dessinée (le glyphe ✓ n'existe pas dans les polices du document)
\newcommand{\popcheck}[1]{\tikz[baseline=-0.5ex]\draw[#1,line width=1.6pt,line cap=round,line join=round](-0.13,0.0)--(-0.04,-0.09)--(0.13,0.1);}
\providecommand{\coche}{\popcheck{popgreen}}
\providecommand{\resultat}[1]{\par\smallskip\noindent\fcolorbox{popgreen}{popgreenL}{\parbox{\dimexpr\linewidth-2\fboxsep-2\fboxrule\relax}{\textbf{Résultat.} #1}}\par\smallskip}
\newcommand{\popboxhead}[3]{\poppill{#1}{#2}{white}\ifx\relax#3\relax\else\hspace{6pt}{\popxbold\fontsize{10.6}{12}\selectfont\color{popink}#3}\fi\par\vspace{7pt}}

\newtcolorbox{aretenir}[1][]{popbase,colframe=popgreen,before upper={\popboxhead{À retenir}{popgreen}{#1}}}
\newtcolorbox{exemplebox}[1][]{popbase,colframe=poporange,before upper={\popboxhead{Exemple}{poporange}{#1}}}
\newtcolorbox{definition}[1][]{popbase,colframe=popblue,before upper={\popboxhead{Définition}{popblue}{#1}}}
\newtcolorbox{propriete}[1][]{popbase,colframe=poppurple,before upper={\popboxhead{Propriété}{poppurple}{#1}}}
\newtcolorbox{methode}[1][]{popbase,colframe=popgold,before upper={\popboxhead{Méthode}{popgold}{#1}}}
\newtcolorbox{attention}[1][]{popbase,colframe=poppink,before upper={\popboxhead{Attention}{poppink}{#1}}}
\newtcolorbox{experience}[1][]{popbase,colframe=popblue,colback=popblueL,before upper={\popboxhead{Expérience}{popblue}{#1}}}
% « Le savais-tu ? » : carte violette pleine (contexte, culture, Cameroun)
\newtcolorbox{savaistu}[1][]{popbase,colback=poppurple,colframe=popink,
  fontupper=\popsemi\fontsize{10.4}{14.2}\selectfont\color{white},
  before upper={\poppill{Le savais-tu\,?}{white}{poppurple}\ifx\relax#1\relax\else\hspace{6pt}{\popxbold\color{white}#1}\fi\par\vspace{7pt}}}
% Exercice résolu : énoncé en haut, \tcblower puis la solution sur fond crème
\newcounter{popexo}\newcounter{popexr}
\newtcolorbox{exoresolu}[1][]{popbase,colframe=poporange,colbacklower=popcream,
  segmentation style={popink,line width=1.2pt,dashed},
  before upper={\stepcounter{popexr}\popboxhead{Exercice résolu \thepopexr}{poporange}{#1}},
  before lower={\poppill{Solution}{popink}{popcream}\par\vspace{6pt}}}
% Exercice (non résolu en place) — la correction est dans la partie Corrigés
\newtcolorbox{exercice}[1][]{popbase,colframe=popink,
  before upper={\stepcounter{popexo}\popboxhead{Exercice \thepopexo}{popink}{#1}}}
% Corrigé
\newtcolorbox{corrige}[1][]{popbase,colframe=popgreen,colback=popgreenL,
  before upper={\popboxhead{Corrigé}{popgreen}{#1}}}
% Objectifs / prérequis / bilan
\newtcolorbox{objectifs}{popbase,colframe=popink,colback=poporangeL,
  before upper={\popboxhead{Objectifs de la leçon}{popink}{}}}
\newtcolorbox{prerequis}{popbase,colframe=popink,colback=popblueL,
  before upper={\popboxhead{Prérequis}{popblue}{}}}
\newtcolorbox{bilan}{popbase,colframe=popink,colback=white,boxrule=2.8pt,
  before upper={\popboxhead{Fiche bilan}{poporange}{L'essentiel à connaître pour l'examen}}}
% Figure encadrée avec légende
\newtcolorbox{popfigure}[1][]{enhanced,breakable=false,colback=white,colframe=popline,boxrule=1.4pt,arc=8pt,
  left=10pt,right=10pt,top=10pt,bottom=8pt,before skip=10pt,after skip=12pt,halign=center,
  lower separated=false,halign lower=center,
  fontlower=\popsemi\fontsize{9}{11.5}\selectfont\color{popmuted},
  #1}
\newcommand{\legende}[1]{\par\vspace{6pt}{\popsemi\fontsize{9}{11.5}\selectfont\color{popmuted}#1}}

% ============================================================
% Signature OnBuch+
% ============================================================
\newcommand{\poplogoinline}[1]{{\poplogo\fontsize{#1}{#1}\selectfont\color{popink}OnBuch\color{poporange}+}}
\newcommand{\signatureonbuch}{%
  \begin{tcolorbox}[enhanced,colback=popink,colframe=popink,boxrule=2.2pt,arc=10pt,
      drop shadow=poporange,left=14pt,right=14pt,top=10pt,bottom=10pt,before skip=0pt,after skip=0pt]
    \tikz[baseline=-0.7ex]\node[circle,fill=popgreen,draw=popcream,line width=1.3pt,minimum size=22pt,inner sep=0]{\popcheck{white}};%
    \hspace{9pt}\begin{minipage}[c]{0.62\linewidth}
      {\popxbold\fontsize{12}{13}\selectfont\color{popcream}Signé\ \textbullet\ {\poplogo OnBuch\color{poporange}+}}\par\vspace{2pt}
      {\raggedright\popsemi\fontsize{8}{10.5}\selectfont\color{popline}\MakeUppercase{Équipe pédagogique OnBuch+}\\\MakeUppercase{Conforme au programme officiel MINESEC}\par}
    \end{minipage}\hfill
    {\poplogo\fontsize{22}{22}\selectfont\color{popcream}OnBuch\color{poporange}+}
  \end{tcolorbox}%
}

% ============================================================
% Page de garde
% #1 titre · #2 matière · #3 niveau · #4 module · #5 n° leçon · #6 durée ·
% #7 description / objectifs (texte ou itemize)
% ============================================================
\newcommand{\pagegarde}[7]{%
  \thispagestyle{empty}
  \newgeometry{a4paper,margin=1.7cm,top=1.6cm,bottom=1.4cm}%
  \begin{tikzpicture}[remember picture,overlay]
    \fill[poporange!18] ([xshift=-3.2cm,yshift=-3.6cm]current page.north east) circle (4.6cm);
    \fill[poppurple!14] ([xshift=2.4cm,yshift=7.6cm]current page.south west) circle (3.8cm);
  \end{tikzpicture}%
  {\poplogo\fontsize{30}{30}\selectfont\color{popink}OnBuch\color{poporange}+}\hfill
  \raisebox{0.6ex}{\poppill{Cours signé \textbullet\ Premium}{popink}{popcream}}\par
  \vspace{1pt}\popsquiggle{poppurple}\par
  \vspace{8pt}
  \begin{tcolorbox}[enhanced,colback=poporange,colframe=popink,boxrule=2.8pt,arc=13pt,
      drop shadow=popink,left=18pt,right=18pt,top=12pt,bottom=13pt,after skip=10pt]
    \poppill{#2 \textbullet\ #3}{popink}{popcream}\hspace{5pt}\poppill{#4}{white}{popink}\par\vspace{8pt}
    {\popdisplay\fontsize{14}{14}\selectfont\color{popink}LEÇON #5}\par\vspace{4pt}
    {\raggedright\hyphenpenalty=10000\exhyphenpenalty=10000\popdisplay\fontsize{25}{27}\selectfont\color{popcream}\MakeUppercase{#1}\par}
    \vspace{8pt}
    {\popxbold\fontsize{9.5}{9.5}\selectfont\color{popink}\MakeUppercase{Durée conseillée : #6}}
  \end{tcolorbox}
  \begin{tcolorbox}[enhanced,colback=white,colframe=popink,boxrule=2.2pt,arc=10pt,
      drop shadow=popink,left=16pt,right=16pt,top=10pt,bottom=10pt,after skip=8pt]
    \poppill{Dans cette leçon}{poppurple}{white}\par\vspace{5pt}
    \popsemi\fontsize{9.3}{12.6}\selectfont\color{popink}#7
  \end{tcolorbox}
  \noindent\begin{minipage}[t]{0.487\linewidth}
    \begin{tcolorbox}[enhanced,colback=popgreen,colframe=popink,boxrule=2.2pt,arc=10pt,
        drop shadow=popink,left=11pt,right=11pt,top=10pt,bottom=10pt,height=3.1cm,valign=center]
      \tcbox[colback=white,colframe=popink,boxrule=1.3pt,arc=5pt,boxsep=0pt,nobeforeafter,
        left=3pt,right=3pt,top=3pt,bottom=3pt,valign=center]{\qrcode[height=1.9cm]{https://play.google.com/store/apps/details?id=com.luvvix.onbuch&hl=fr}}%
      \hspace{8pt}\begin{minipage}[c]{0.5\linewidth}
        {\popdisplay\fontsize{11}{12.5}\selectfont\color{white}\MakeUppercase{Télécharge OnBuch}}\par\vspace{4pt}
        {\raggedright\popsemi\fontsize{8.4}{11}\selectfont\color{white}Gratuit \textbullet\ Google Play\\Tuteur IA, annales, concours, résultats\par}
      \end{minipage}
    \end{tcolorbox}
  \end{minipage}\hfill
  \begin{minipage}[t]{0.487\linewidth}
    \begin{tcolorbox}[enhanced,colback=poppurple,colframe=popink,boxrule=2.2pt,arc=10pt,
        drop shadow=popink,left=11pt,right=11pt,top=10pt,bottom=10pt,height=3.1cm,valign=center]
      \tikz[baseline=-0.7ex]\node[circle,fill=white,draw=popink,line width=1.3pt,minimum size=1.6cm,inner sep=0]{\popdisplay\fontsize{24}{24}\selectfont\color{poppurple}+};%
      \hspace{8pt}\begin{minipage}[c]{0.6\linewidth}
        {\popdisplay\fontsize{11}{12.5}\selectfont\color{white}\MakeUppercase{Passe à OnBuch+}}\par\vspace{4pt}
        {\raggedright\popsemi\fontsize{8.4}{11}\selectfont\color{white}Tous les cours signés, Tuteur IA illimité, fiches PDF — onglet OnBuch+ de l'app\par}
      \end{minipage}
    \end{tcolorbox}
  \end{minipage}
  \vspace{8pt}
  \popcontenu
  \vfill
  \signatureonbuch
  \restoregeometry
  \newpage
}

% Rangée « Ce cours contient » (page de garde)
\newcommand{\poptile}[3]{% #1 couleur #2 gros texte #3 légende
  \begin{tcolorbox}[enhanced,colback=#1,colframe=popink,boxrule=2pt,arc=9pt,shadow={2.5pt}{-2.5pt}{0pt}{popink},
      left=6pt,right=6pt,top=9pt,bottom=9pt,halign=center,valign=center,height=1.75cm,width=\linewidth,nobeforeafter]
    {\popdisplay\fontsize{12.5}{13}\selectfont\color{white}\MakeUppercase{#2}}\par\vspace{3pt}
    {\centering\popsemi\fontsize{7.8}{9.5}\selectfont\color{white}#3\par}
  \end{tcolorbox}}
\newcommand{\popcontenu}{%
  \par\noindent{\popxboldls\fontsize{7.5}{7.5}\selectfont\color{popmuted}CE COURS SIGNÉ CONTIENT}\par\vspace{7pt}
  \noindent\begin{minipage}[t]{0.235\linewidth}\poptile{popblue}{Cours}{complet, illustré, pas à pas}\end{minipage}\hfill
  \begin{minipage}[t]{0.235\linewidth}\poptile{poporange}{Méthodes}{et exercices résolus}\end{minipage}\hfill
  \begin{minipage}[t]{0.235\linewidth}\poptile{poppink}{Exercices}{type examen + corrigés}\end{minipage}\hfill
  \begin{minipage}[t]{0.235\linewidth}\poptile{popgreen}{Fiche bilan}{l'essentiel à retenir}\end{minipage}\par}

% Sommaire
\newtcolorbox{sommaire}{enhanced,colback=white,colframe=popink,boxrule=2.2pt,arc=10pt,
  drop shadow=popink,left=18pt,right=18pt,top=13pt,bottom=13pt,before skip=6pt,after skip=20pt,
  before upper={\poppill{Sommaire}{poppurple}{white}\par\vspace{9pt}\popsemi\fontsize{10.6}{14.5}\selectfont\color{popink}}}

% Signature de fin de cours — poussée tout en bas de la dernière page
\newcommand{\findecours}{%
  \par\vfill\nopagebreak
  \begin{center}\popsquiggle{poporange}\end{center}\vspace{4pt}
  \signatureonbuch
}
\AtBeginDocument{\def\llbracket{\mathopen{[\mkern-3.5mu[}}\def\rrbracket{\mathclose{]\mkern-3.5mu]}}} % intervalles d'entiers (glyphe ⟦ absent de la police)
% Glyphes absents de Fira Math : redéfinis à partir de glyphes disponibles
\AtBeginDocument{%
  \def\setminus{\mathbin{\backslash}}\let\smallsetminus\setminus
  \def\vdots{\mathord{\vbox{\baselineskip3.2pt\lineskiplimit0pt\kern4pt\hbox{.}\hbox{.}\hbox{.}}}}%
  \def\ddots{\mathinner{\mkern1mu\raise6.5pt\hbox{.}\mkern2mu\raise3.6pt\hbox{.}\mkern2mu\raise0.7pt\hbox{.}\mkern1mu}}%
  \def\triangle{\mathord{\scalebox{0.9}{$\symup{\Delta}$}}}%
  \def\nabla{\mathord{\raisebox{\depth}{\scalebox{1}[-1]{$\symup{\Delta}$}}}}%
  \def\checkmark{\popcheck{popdark}}%
  \def\star{\mathbin{\ast}}\def\bigstar{\mathord{\ast}}%
  \def\diamond{\mathbin{\scalebox{0.6}{$\Diamond$}}}%
  \def\bigcup{\mathop{\mathchoice{\raisebox{-0.3em}{\scalebox{1.7}{$\cup$}}}{\scalebox{1.25}{$\cup$}}{\cup}{\cup}}\displaylimits}%
  \def\bigcap{\mathop{\mathchoice{\raisebox{-0.3em}{\scalebox{1.7}{$\cap$}}}{\scalebox{1.25}{$\cap$}}{\cap}{\cap}}\displaylimits}%
  \def\lesseqqgtr{\mathrel{\lessgtr}}\def\bigoplus{\mathop{\oplus}\displaylimits}%
}
\providecommand{\ssi}{si et seulement si} % abréviation fréquente
\AtBeginDocument{\providecommand{\wideparen}{\overparen}} % arc de cercle

\begin{document}

\pagegarde{I :Produire un plan détaillé d’une dissertation}{Français}{1re Tech}{Français – classe de Première technique}{N17}{10 séances (fiche de progression harmonisée)}{Cette leçon outille l'élève de Première technique pour construire, étape par étape, le plan détaillé d'une dissertation : de l'analyse du sujet à la problématique, de la recherche des idées à l'organisation du plan. Elle s'appuie sur l'étude du Lion et la Perle de Wole Soyinka et mobilise les notions de communication et de ponctuation forte.
\vspace{4pt}
\begin{multicols}{2}\small
\begin{itemize}
\item À la fin de cette leçon, je sais identifier et exploiter les paratextes d'une œuvre littéraire (titre, couverture, quatrième de couverture, biographie de l'auteur).
\item À la fin de cette leçon, je sais analyser un sujet de dissertation (dégager le thème, les mots-clés, la consigne).
\item À la fin de cette leçon, je sais formuler une problématique pertinente à partir d'un sujet donné.
\item À la fin de cette leçon, je sais mobiliser les techniques de recherche d'idées (brainstorming, mobilisation des connaissances, exploitation de documents).
\item À la fin de cette leçon, je sais distinguer et construire les types de plans (thématique, dialectique, analytique).
\item À la fin de cette leçon, je sais élaborer un plan détaillé complet (introduction, développement en parties et sous-parties, conclusion).
\end{itemize}\end{multicols}}

\begin{sommaire}
\begin{multicols}{2}
\begin{enumerate}
\item Réception du texte : paratextes et lecture ouvroir du Lion et la Perle
\item Analyse du sujet et formulation de la problématique
\item La recherche des idées
\item Les types de plans et l'élaboration du plan détaillé
\item Langue : conditions et situations de communication
\item Langue : la ponctuation forte et lecture méthodique du Lion et la Perle
\item Activité d'intégration
\item Méthodes et exercices résolus
\item Exercices
\item Corrigés des exercices
\item Fiche bilan
\end{enumerate}
\end{multicols}
\end{sommaire}

\begin{objectifs}
\begin{itemize}
\item À la fin de cette leçon, je sais identifier et exploiter les paratextes d'une œuvre littéraire (titre, couverture, quatrième de couverture, biographie de l'auteur).
\item À la fin de cette leçon, je sais analyser un sujet de dissertation (dégager le thème, les mots-clés, la consigne).
\item À la fin de cette leçon, je sais formuler une problématique pertinente à partir d'un sujet donné.
\item À la fin de cette leçon, je sais mobiliser les techniques de recherche d'idées (brainstorming, mobilisation des connaissances, exploitation de documents).
\item À la fin de cette leçon, je sais distinguer et construire les types de plans (thématique, dialectique, analytique).
\item À la fin de cette leçon, je sais élaborer un plan détaillé complet (introduction, développement en parties et sous-parties, conclusion).
\item À la fin de cette leçon, je sais identifier les conditions et les situations de communication dans un échange réel ou littéraire.
\item À la fin de cette leçon, je sais employer correctement la ponctuation forte (point, point d'interrogation, point d'exclamation, points de suspension) dans mes productions écrites.
\item À la fin de cette leçon, je sais rendre compte d'une lecture et justifier mes choix de plan par écrit.
\end{itemize}
\end{objectifs}

\begin{prerequis}
\begin{itemize}
\item Notions de base sur le texte argumentatif acquises en classe de Seconde
\item La structure générale d'une dissertation (introduction, développement, conclusion)
\item Les types de textes et les genres littéraires (théâtre notamment)
\item La phrase simple et la phrase complexe (morphosyntaxe de Seconde)
\item La lecture cursive d'une œuvre intégrale
\end{itemize}
\end{prerequis}

\newpage

\coursec{Réception du texte : paratextes et lecture ouvroir du Lion et la Perle}

Lors du concours d'entrée à l'ENAM, un candidat de Yaoundé rédige une dissertation riche en idées mais désorganisée : il obtient 8/20. Son voisin, avec les mêmes idées mais un plan clair en deux parties, obtient 15/20. Au marché Mokolo comme dans la salle d'examen, convaincre exige d'abord d'organiser. Comment transformer une problématique et des idées éparses en un plan détaillé cohérent, capable de guider la rédaction d'une dissertation réussie ? C'est tout l'enjeu de cette leçon, que nous mènerons en parallèle de notre lecture du \emph{Lion et la Perle} de Wole Soyinka.

\courssub{Présentation de Wole Soyinka et contexte de l'œuvre}

\begin{popfigure}
\begin{tikzpicture}[
  node distance=1.2cm and 1.5cm,
  every node/.style={font=\small},
  title/.style={rectangle, draw=popblue, thick, fill=popblueL, rounded corners, minimum width=3.5cm, text centered},
  cat/.style={rectangle, draw=poporange, fill=poporangeL, rounded corners, minimum width=3cm, text centered},
  detail/.style={rectangle, draw=popgreen, fill=popgreenL, rounded corners, minimum width=4cm, text centered, text width=3.8cm, align=center}
]
\node[title] (oeuvre) {\textbf{Le Lion et la Perle}};
\node[cat, below left=of oeuvre] (auteur) {Auteur};
\node[detail, left=of auteur, align=center] (auteurd){Wole Soyinka\\Dramaturge nigérian\\Prix Nobel 1986};
\node[cat, below=of oeuvre] (genre) {Genre};
\node[detail, below=of genre, align=center] (genred){Comédie satirique\\Théâtre en trois actes\\« Morning, Noon, Night »};
\node[cat, below right=of oeuvre] (date) {Date \& Lieu};
\node[detail, right=of date, align=center] (dated){Écrit vers 1959\\Publié 1963\\Nigeria (Yoruba land)};
\node[cat, above left=of oeuvre] (contexte) {Contexte};
\node[detail, above=of contexte, align=center] (contexted){Nigeria pré-indépendance\\Années 1950-1960\\Tradition vs Modernité};
\node[cat, above right=of oeuvre] (themes) {Thèmes majeurs};
\node[detail, above=of themes, align=center] (themesd){Pouvoir \& séduction\\Éducation \& culture\\Identité africaine};
\draw[popblue, thick, ->] (oeuvre.south west) -- (auteur.north east);
\draw[popblue, thick, ->] (oeuvre.south) -- (genre.north);
\draw[popblue, thick, ->] (oeuvre.south east) -- (date.north west);
\draw[popblue, thick, ->] (oeuvre.north west) -- (contexte.south east);
\draw[popblue, thick, ->] (oeuvre.north east) -- (themes.south west);
\end{tikzpicture}
\legende{Carte d'identité de l'œuvre : \emph{Le Lion et la Perle} de Wole Soyinka.}
\end{popfigure}

\textbf{Wole Soyinka} (né en 1934 à Abeokuta, Nigeria) est une figure majeure de la littérature africaine et mondiale. Dramaturge, poète, essayiste et militant politique, il devient en \textbf{1986 le premier Africain à recevoir le prix Nobel de littérature}, salué pour son théâtre qui « dans une perspective culturelle large et avec des accents poétiques, façonne le drame de l'existence ». Son œuvre mêle traditions orales yoruba, mythologie, satire sociale et engagement politique. Il a connu la prison sous la dictature militaire nigériane (1967-1969) pour avoir tenté de médier la guerre du Biafra.

\begin{savaistu}[Wole Soyinka, pionnier nobélisé]
Wole Soyinka est le \textbf{premier Africain} (et le premier auteur noir africain) à recevoir le prix Nobel de littérature en 1986. Avant lui, le prix était allé à des auteurs européens, américains ou asiatiques. Son discours de réception, « This Past Must Address Its Present », est une charge contre l'apartheid et les dictatures africaines. Au Cameroun, son théâtre inspire les troupes universitaires et le Festival de théâtre de Yaoundé.
\end{savaistu}

\emph{Le Lion et la Perle} (titre original : \emph{The Lion and the Jewel}) est écrit vers 1959, à la veille de l'indépendance du Nigeria (1960). La pièce se déroule dans le village fictif d'Ilujinle, au cœur du pays yoruba. Elle met en scène trois personnages principaux : \textbf{Baroka}, le Bale (chef traditionnel), surnommé « le Lion » ; \textbf{Sidi}, la belle villageoise, « la Perle » ; et \textbf{Lakunle}, l'instituteur « moderne », partisan de la civilisation occidentale. La pièce est une comédie satirique qui oppose deux visions du monde : la tradition africaine (incarnée par Baroka) et la modernité importée (incarnée par Lakunle), avec Sidi comme enjeu de séduction et symbole de l'avenir du village.

\courssub{Le paratexte : définition, composantes et fonctions}

\begin{definition}[Paratexte (Gérard Genette)]
Le \cle{paratexte} est l'ensemble des éléments qui entourent le texte principal et en accompagnent la présentation, la diffusion et la réception. Selon Gérard Genette (\emph{Seuils}, 1987), il constitue un « seuil » qui « offre à tous le monde la possibilité d'entrer dans le texte ou de s'en détourner ». Il se divise en deux grandes catégories :
\begin{itemize}
\item Le \textbf{péritexte} : éléments matériellement attachés au livre (titre, couverture, préface, dédicace, notes, épilogue, etc.).
\item L'\textbf{épitexte} : éléments extérieurs au livre mais qui en parlent (interviews, critiques, correspondances de l'auteur, publicités, adaptations).
\end{itemize}
\end{definition}

\begin{popfigure}
\begin{tikzpicture}[
  node distance=1.5cm and 2cm,
  every node/.style={font=\small},
  root/.style={rectangle, draw=poppurple, thick, fill=poppurpleL, rounded corners, minimum width=3cm, text centered},
  branch/.style={rectangle, draw=popblue, fill=popblueL, rounded corners, minimum width=2.8cm, text centered},
  leaf/.style={rectangle, draw=popgreen, fill=popgreenL, rounded corners, minimum width=2.5cm, text centered, text width=2.2cm, align=center}
]
\node[root, align=center] (paratexte){PARATEXTE\\ (Seuil du texte)};
\node[branch, below left=of paratexte, align=center] (peritexte){PÉRITEXTE\\ (Dans le livre)};
\node[branch, below right=of paratexte, align=center] (epitexte){ÉPITEXTE\\ (Hors du livre)};
% Peritexte children
\node[leaf, below left=of peritexte] (titre) {Titre / Sous-titre};
\node[leaf, below=of peritexte] (couverture) {Couverture / Illustration};
\node[leaf, below right=of peritexte] (quatrieme) {Quatrième de couverture};
\node[leaf, left=of titre] (auteur) {Nom de l'auteur};
\node[leaf, right=of quatrieme] (preface) {Préface / Avant-propos};
\node[leaf, below=of titre] (dedicace) {Dédicace / Épigraphe};
\node[leaf, below=of couverture] (notes) {Notes / Postface};
% Epitexte children
\node[leaf, below left=of epitexte] (interviews) {Interviews de l'auteur};
\node[leaf, below=of epitexte] (critiques) {Critiques / Recensions};
\node[leaf, below right=of epitexte] (pub) {Publicité / Bandeau};
\node[leaf, right=of pub] (adapt) {Adaptations (film, théâtre)};
\draw[poppurple, thick] (paratexte) -- (peritexte);
\draw[poppurple, thick] (paratexte) -- (epitexte);
\draw[popblue, thick] (peritexte) -- (titre);
\draw[popblue, thick] (peritexte) -- (couverture);
\draw[popblue, thick] (peritexte) -- (quatrieme);
\draw[popblue, thick] (peritexte) -- (auteur);
\draw[popblue, thick] (peritexte) -- (preface);
\draw[popblue, thick] (peritexte) -- (dedicace);
\draw[popblue, thick] (peritexte) -- (notes);
\draw[poporange, thick] (epitexte) -- (interviews);
\draw[poporange, thick] (epitexte) -- (critiques);
\draw[poporange, thick] (epitexte) -- (pub);
\draw[poporange, thick] (epitexte) -- (adapt);
\end{tikzpicture}
\legende{Arbre des composantes du paratexte selon Genette : péritexte (interne) et épitexte (externe).}
\end{popfigure}

\textbf{Fonctions du paratexte} (trois fonctions principales) :
\begin{enumerate}
\item \textbf{Informer} : donner l'identité de l'œuvre (auteur, titre, genre, date, éditeur, collection), résumer l'intrigue (quatrième de couverture), situer le contexte.
\item \textbf{Orienter la lecture} : la préface propose des clés d'interprétation ; le titre suggère des pistes ; la couverture crée une atmosphère visuelle.
\item \textbf{Susciter l'attente du lecteur} : le paratexte éveille la curiosité, crée un « horizon d'attente » (Jauss) qui guide la réception. Un titre intriguant, une couverture attractive, une préface d'un auteur célèbre incitent à lire.
\end{enumerate}

\begin{attention}[Erreur fréquente : résumé vs analyse]
Ne confondez \textbf{pas} le résumé de la quatrième de couverture (paratexte à fonction informative et publicitaire) avec l'analyse littéraire de l'œuvre. La quatrième de couverture est écrite par l'éditeur pour \emph{vendre} le livre ; elle simplifie, parfois trahit la complexité du texte. L'analyse, elle, s'appuie sur la lecture intégrale de l'œuvre, cite le texte, examine les procédés d'écriture. Au Probatoire, on vous demande une analyse, pas une paraphrase du paratexte.
\end{attention}

\courssub{Analyse du titre : \emph{Le Lion et la Perle}}

Le titre \emph{Le Lion et la Perle} (traduction de \emph{The Lion and the Jewel}) est hautement symbolique. Il nomme les deux protagonistes principaux et condense le conflit central.

\begin{center}
\begin{tabularx}{\linewidth}{|l|X|X|}\hline
\textbf{Terme} & \textbf{Personnage} & \textbf{Symbolique} \\ \hline
\cle{Le Lion} & Baroka, le Bale (chef traditionnel) & Force, pouvoir, ruse, virilité, autorité ancestrale, royauté yoruba. Le lion est le roi de la brousse, respecté et craint. Baroka incarne la tradition vivante, capable de s'adapter (il conduit une voiture, veut le téléphone). \\ \hline
\cle{La Perle} & Sidi, la belle villageoise & Beauté, pureté, valeur, fragilité, objet de convoitise. La perle naît dans l'huître (le village) ; elle est rare, précieuse. Sidi est la « perle » que veulent posséder les deux hommes : Baroka par la ruse traditionnelle, Lakunle par le « mariage moderne » (dot refusée). \\ \hline
\cle{Et} (conjonction) & — & Opposition / confrontation / alternative. Le titre ne dit pas « Le Lion \emph{ou} la Perle » mais « \emph{et} » : les deux sont liés, le sort de l'un décide de l'autre. \\ \hline
\end{tabularx}
\end{center}

\begin{exemplebox}[Analyse du titre en situation d'examen]
\textbf{Sujet :} « Que vous suggère le titre \emph{Le Lion et la Perle} sur les enjeux de la pièce ? »
\textbf{Réponse type :} Le titre met en présence deux symboles forts de la culture yoruba. Le Lion (Baroka) représente le pouvoir traditionnel, la sagesse des anciens, la séduction par l'expérience. La Perle (Sidi) représente la jeunesse, la beauté, l'avenir du village, mais aussi l'enjeu d'une modernité qui veut la « civiliser » (Lakunle). La conjonction « et » annonce une rencontre, un affrontement, un marché : qui possédera la perle ? Le titre résume ainsi le conflit tradition/modernité et la place de la femme comme objet d'échange symbolique.
\end{exemplebox}

\courssub{La lecture ouvroir : démarche active et carnet de bord}

\begin{methode}[Mener une lecture ouvroir (lecture active avec carnet de bord)]
La \cle{lecture ouvroir} (ou lecture-atelier) transforme le lecteur passif en lecteur-acteur. Elle se pratique en trois temps :
\begin{enumerate}
\item \textbf{Avant la lecture (paratexte)} : Observer la couverture, lire le titre, l'auteur, la quatrième de couverture, la préface. Noter ses \emph{attentes}, ses \emph{hypothèses} de lecture, ses \emph{questions}.
\item \textbf{Pendant la lecture (texte)} : Lire avec un crayon. \textbf{Annoter} dans les marges : souligner les passages frappants, noter des \textbf{?} (incompréhension), des \textbf{!} (surprise), des \textbf{$\leftrightarrow$} (liens avec d'autres textes, le contexte). Rédiger des \emph{réactions immédiates} en marge ou sur un carnet à part.
\item \textbf{Après chaque acte / scène (bilan)} : Consigner dans le \textbf{carnet de bord} : résumé synthétique, personnages nouveaux, thèmes émergents, procédés de langue (proverbes, chants, ironie), hypothèses révisées, questions pour la suite.
\end{enumerate}
Le carnet de bord devient la matière première pour le compte rendu, le commentaire et la dissertation.
\end{methode}

\begin{exoresolu}[Première page de carnet de bord pour \emph{Le Lion et la Perle}]
\textbf{Énoncé :} Rédigez la première page de votre carnet de bord après la lecture des paratextes et de la première scène (Morning).
\tcblower
\textbf{Carnet de bord — \emph{Le Lion et la Perle} — 15/09/2025}
\textbf{Paratexte :} Couverture : photo d'un chef yoruba coiffé de perles, fond ocre. Titre en rouge. Quatrième : « Comédie sur le choc des cultures en pays yoruba ». Attente : humour, conflit générationnel, femme objet ?
\textbf{Hypothèses :} 1) Baroka va tromper Lakunle. 2) Sidi choisira la tradition. 3) La fin sera ironique, pas tragique.
\textbf{Lecture Acte 1 (Morning) :}
- Lakunle : parle « anglais de collège », refuse la dot (« coutume barbare »), veut transformer Sidi en « femme moderne ». \textbf{!} Caricature ? Satire de l'assimilation ?
- Sidi : fière de sa beauté (photos du magazine), refuse Lakunle sans dot. \textbf{$\leftrightarrow$} Vanité, mais aussi dignité.
- Baroka (évoqué) : « le Renard », 62 ans, veut Sidi. \textbf{?} Comment va-t-il s'y prendre ?
\textbf{Thèmes naissants :} Tradition vs Modernité ; Condition féminine ; Pouvoir de l'image (photos) ; Langue (pidgin vs anglais standard).
\textbf{Questions :} Quel rôle pour Sadiku (épouse aînée) ? La danse du « voyageur perdu » : symbole ?
\end{exoresolu}

\courssub{Premières hypothèses de lecture : tradition contre modernité}

À l'issue de la lecture ouvroir des paratextes et de la première partie (Morning), trois grandes hypothèses structurent notre entrée dans l'œuvre :

\begin{propriete}[Hypothèses de lecture fondées sur le paratexte et l'amorce]
\begin{enumerate}
\item \textbf{Conflit central : Tradition vs Modernité.} Non pas une opposition binaire (le bien vs le mal), mais une \emph{négociation} : Baroka (tradition) utilise la modernité (timbre-poste, machine à coudre, voiture) pour asseoir son pouvoir ; Lakunle (modernité) est ridicule dans son mimétisme occidental. L'œuvre suggère une \textbf{synthèse africaine} possible.
\item \textbf{La femme comme enjeu symbolique.} Sidi (la Perle) n'est pas un personnage passif. Sa beauté (reconnue par le magazine \emph{Man}) lui donne un pouvoir médiatique. Elle négocie, refuse, choisit. L'hypothèse : Sidi incarne l'\textbf{avenir du village}, disputé entre deux modèles de masculinité.
\item \textbf{Le théâtre comme espace de parole yoruba.} Soyinka intègre proverbes, chants, danses, mime (le voyageur perdu). Hypothèse : la forme théâtrale \emph{est} le message : la tradition n'est pas un musée, c'est une scène vivante où se joue l'identité.
\end{enumerate}
\end{propriete}

\courssub{Contexte historique et parallèle camerounais}

\begin{center}
\begin{tabularx}{\linewidth}{|X|X|}\hline
\textbf{Nigeria (années 1950-1960)} & \textbf{Parallèle camerounais (même époque)} \\ \hline
Colonie britannique, indépendance 1960. & Cameroun français (indépendance 1960) et britannique (réunification 1961). \\
Société yoruba : chefferies traditionnelles (Oba, Bale) coexistent avec administration coloniale. & Chefferies traditionnelles (Lamidats au Nord, chefferies Bamiléké, Bassa, Douala) encadrées par l'administration coloniale puis l'État. \\
École missionnaire/coloniale forme une élite « assimilée » (Lakunle). & Écoles missionnaires (catholiques, protestantes, laïques) forment les « évolués », fonctionnaires, instituteurs. \\
Urbanisation rapide (Lagos, Ibadan) : migration jeunes, rupture village/ville. & Urbanisation (Douala, Yaoundé) : exode rural, bidonvilles, nouvelles identités. \\
Débat : « African personality » vs « Westernization » (Soyinka, Achebe, Senghor). & Débat : « Authenticité » (Ahmadou Ahidjo, puis Biya) vs modernisation ; littérature (Mongo Beti, Ferdinand Oyono) dénonce l'aliénation. \\ \hline
\end{tabularx}
\end{center}

Ce parallèle n'est pas anecdotique : il ancre la lecture de Soyinka dans l'expérience vécue des élèves camerounais. Le village d'Ilujinle ressemble à bien des villages du Cameroun des années 1960 (et d'aujourd'hui) : le chef traditionnel (fon, lamido, chef de groupement) qui négocie avec l'administration, l'instituteur qui « parle blanc » et refuse la dot, la jeune fille courtisée pour sa beauté et son statut symbolique.

\begin{aretenir}[Points clés pour la suite]
\begin{itemize}
\item Le paratexte est un \textbf{seuil herméneutique} : il guide l'entrée dans le texte, ne le remplace pas.
\item La lecture ouvroir (carnet de bord, annotations, hypothèses) est la \textbf{méthode professionnelle} pour passer de la lecture impressionniste à l'analyse argumentée.
\item \emph{Le Lion et la Perle} met en scène, sous forme comique, le \textbf{grand débat de l'Afrique des indépendances} : comment être moderne sans s'aliéner ? La réponse de Soyinka : par le rire, la ruse, et la réappropriation critique de la tradition.
\item Le parallèle Nigeria-Cameroun permet d'ancrer l'universel de l'œuvre dans le \textbf{concret local}, atout majeur pour la dissertation et l'oral au Probatoire.
\end{itemize}
\end{aretenir}

\begin{exercice}[Entraînement : paratexte et hypothèses]
Vous découvrez un roman camerounais dont voici les éléments paratextuels :
\begin{itemize}
\item Titre : \emph{Le Cri de la terre}
\item Auteur : Inconnu de vous, né à Dschang en 1948.
\item Couverture : photo d'un champ de caféiers en fleurs, au premier plan une main noire tenant une poignée de grains rouges.
\item Quatrième de couverture : « Dans les années 1950, au cœur du pays bamiléké, un jeune cultivateur refuse de vendre sa récolte à la coopérative coloniale. Il rêve d'une Afrique qui maîtrise son économie. Roman historique et initiatique. »
\end{itemize}
\textbf{Consignes :}
\begin{enumerate}
\item Identifiez les éléments de péritexte et d'épitexte (imaginables) présents ici.
\item Formulez trois hypothèses de lecture (thèmes, conflit, ton) en vous appuyant \emph{uniquement} sur ce paratexte.
\item Quelle fonction remplit la mention « Roman historique et initiatique » ?
\end{enumerate}
\end{exercice}

\begin{corrige}[Exercice n°1]
\begin{enumerate}
\item \textbf{Péritexte :} titre, nom de l'auteur, couverture (image), quatrième de couverture. \textbf{Épitexte imaginable :} préface d'un historien, interview de l'auteur à la CRTV, bandeau « Prix littéraire », adaptation théâtrale annoncée.
\item \textbf{Hypothèses :} 1) Conflit économico-politique : paysan vs administration coloniale (coopérative) → résistance par la terre. 2) Initiation du héros : passage de l'adolescence à l'âge adulte conscient, prise de conscience nationale. 3) Ton : réaliste, engagé, peut-être lyrique (image de la main/terre).
\item La mention « Roman historique et initiatique » remplit une \textbf{fonction générique} (informer sur le genre) et \textbf{horizon d'attente} : le lecteur attend une double trame (Histoire collective + parcours individuel), une documentation précise, un héros qui apprend et grandit.
\end{enumerate}
\end{corrige}

\coursec{Analyse du sujet et formulation de la problématique}

Dans la section précédente, nous avons appris à recevoir un texte : grâce aux paratextes (titre, couverture, quatrième de couverture, biographie de l'auteur) et à la lecture ouvroir du \emph{Lion et la Perle} de Wole Soyinka, nous avons dégagé les premières hypothèses de lecture. Recevoir un texte, c'est déjà analyser : observer, définir, questionner. Nous allons maintenant utiliser exactement la même démarche, mais appliquée non plus à une œuvre, mais à un \cle{sujet de dissertation}. Car avant d'écrire la moindre phrase d'un devoir, il faut comprendre précisément ce que le sujet demande. C'est toute la différence entre un élève qui « parle du sujet » et un élève qui \emph{répond au sujet}.

\courssub{Qu'est-ce qu'un sujet de dissertation ?}

\begin{definition}[Le sujet de dissertation]
Le \cle{sujet de dissertation} est un énoncé bref — une \textbf{citation}, une \textbf{affirmation} ou une \textbf{question} — qui invite le candidat à mener une réflexion argumentée sur un problème. Il ne demande jamais de réciter des connaissances : il demande de \emph{penser} un sujet, de prendre position et de justifier cette position.
\end{definition}

Au Probatoire et au Baccalauréat technique, le sujet se présente sous trois formes principales :

\begin{itemize}
\item \textbf{La citation} : « "La tradition est la mémoire des peuples", écrit un auteur africain. Partagez-vous ce point de vue ? »
\item \textbf{L'affirmation} : « La tradition est un obstacle au développement. Discutez. »
\item \textbf{La question directe} : « Faut-il moderniser nos villages ? »
\end{itemize}

Dans les trois cas, l'examinateur attend la même chose : une réflexion personnelle, organisée, appuyée sur des arguments et des exemples. Tout commence donc par une étape que beaucoup de candidats négligent : l'\cle{analyse du sujet}.

\begin{attention}[Le piège de la précipitation]
L'erreur la plus grave en dissertation consiste à commencer à écrire après une seule lecture du sujet. Un sujet mal compris entraîne le \cle{hors-sujet}, sanctionné très lourdement à l'examen, même si la copie est bien écrite. Consacrez toujours un quart d'heure environ à l'analyse du sujet : c'est un temps gagné, jamais perdu.
\end{attention}

\courssub{Les quatre étapes de l'analyse du sujet}

Analyser un sujet, c'est le démonter méthodiquement, comme un mécanicien démonte un moteur pour comprendre comment il fonctionne. Cette analyse se fait en quatre étapes successives, qu'on ne doit jamais sauter.

\begin{methode}[Analyser un sujet de dissertation en quatre étapes]
\begin{enumerate}
\item \textbf{Lire et souligner} : lire le sujet au moins trois fois et souligner les mots-clés (les mots qui portent le sens).
\item \textbf{Définir et délimiter} : définir chaque mot-clé avec ses propres mots, puis délimiter le thème (de quoi parle exactement le sujet ? de quoi ne parle-t-il pas ?).
\item \textbf{Repérer les présupposés et la tension} : identifier ce que le sujet affirme implicitement, puis dégager le paradoxe ou l'opposition qui fonde la réflexion.
\item \textbf{Interroger} : transformer cette tension en une question centrale, la problématique.
\end{enumerate}
\end{methode}

\begin{popfigure}
\begin{center}
\begin{tikzpicture}[
  node distance=1.2cm and 2.5cm,
  etape/.style={draw, rounded corners=3pt, minimum width=4.8cm, minimum height=1.2cm, align=center, font=\small, thick},
  fleche/.style={-{Stealth[length=3mm]}, thick}
]
\node[etape, draw=popblue, fill=popblueL, align=center] (e1){\textbf{Étape 1 : les mots-clés}\\ lire, souligner};
\node[etape, draw=poporange, fill=poporangeL, below right=of e1, align=center] (e2){\textbf{Étape 2 : le thème}\\ définir, délimiter};
\node[etape, draw=popgreen, fill=popgreenL, below left=of e2, align=center] (e3){\textbf{Étape 3 : la tension}\\ présupposés, paradoxe};
\node[etape, draw=poppurple, fill=poppurpleL, below right=of e3, align=center] (e4){\textbf{Étape 4 : la problématique}\\ interroger};
\draw[fleche, poporange] (e1) -- (e2);
\draw[fleche, popgreen] (e2) -- (e3);
\draw[fleche, poppurple] (e3) -- (e4);
\end{tikzpicture}
\end{center}
\legende{Les quatre étapes de l'analyse du sujet, en escalier : chaque étape descend d'un niveau de réflexion, des mots du sujet jusqu'à la question centrale qui guidera tout le devoir.}
\end{popfigure}

\courssub{Étape 1 : lire le sujet et souligner les mots-clés}

La première lecture se fait à voix haute ou en murmurant, pour bien sentir le poids de chaque mot. Ensuite, on souligne les \cle{mots-clés}, c'est-à-dire les mots qui portent l'essentiel du sens. On néglige les mots-outils (articles, prépositions, conjonctions) pour se concentrer sur les noms, les verbes et les adjectifs porteurs de la charge sémantique.

\begin{exemplebox}[Repérage des mots-clés]
Sujet : « La \textbf{tradition} est un \textbf{obstacle} au \textbf{développement}. »

Trois mots-clés : \emph{tradition}, \emph{obstacle}, \emph{développement}. Ce sont eux qu'il faudra définir et mettre en relation.
\end{exemplebox}

\courssub{Étape 2 : définir les mots-clés et délimiter le thème}

Chaque mot-clé doit être défini avec ses propres mots, sans dictionnaire, comme on le ferait pour un camarade. Cette définition personnelle est la preuve qu'on a compris le sujet et qu'on s'approprie le vocabulaire du débat.

\begin{center}
\begin{tabularx}{\linewidth}{|l|X|}
\hline
\rowcolor{popblueL} \textbf{Mot-clé} & \textbf{Définition personnelle} \\
\hline
Tradition & Ensemble des coutumes, croyances, langues, rites et savoir-faire transmis de génération en génération dans une communauté. \\
\hline
Obstacle & Ce qui barre la route, ce qui empêche d'avancer ou de progresser ; un frein. \\
\hline
Développement & Progrès économique, social et technique d'une société (infrastructures, éducation, santé, industrie, bien-être des populations). \\
\hline
\end{tabularx}
\end{center}

Définir permet ensuite de \cle{délimiter le thème} : le sujet porte sur le rapport entre la tradition (héritage culturel) et le développement (progrès matériel et social). Il ne porte \emph{pas} sur la politique politicienne, ni sur la religion en général, ni sur l'histoire ancienne pour elle-même. Savoir ce dont le sujet ne parle pas est aussi important que savoir ce dont il parle : c'est la garde-fou contre le hors-sujet.

\courssub{Étape 3 : repérer les présupposés et dégager la tension}

Un sujet de dissertation n'est jamais neutre. Il contient des \cle{présupposés} : ce qu'il affirme implicitement, sans le dire ouvertement, et qu'il faut mettre au jour pour pouvoir le discuter. C'est l'étape heuristique par excellence.

Dans notre exemple, le sujet « La tradition est un obstacle au développement » présuppose :
\begin{itemize}
\item que la tradition et le développement seraient \textbf{opposés}, comme deux forces ennemies (vision binaire) ;
\item que la tradition appartiendrait exclusivement au \textbf{passé} et le développement à l'\textbf{avenir} (vision linéaire du temps) ;
\item que le développement serait forcément un \textbf{bien} absolu auquel il faut tout sacrifier (vision utilitariste).
\end{itemize}

Ces présupposés font apparaître la \cle{tension} (ou paradoxe) qui fonde la réflexion : d'un côté, certaines traditions semblent effectivement freiner le progrès (refus de la scolarisation des filles, pratiques rétrogrades, dépenses excessives des funérailles) ; de l'autre, la tradition est aussi une richesse, une identité, voire un moteur de développement (artisanat, tourisme culturel, solidarité communautaire type tontine, valeurs morales). Faut-il donc détruire la tradition pour se développer, ou peut-on se développer \emph{avec} elle ? C'est ce conflit d'idées qui rend le sujet intéressant — et discutable.

\begin{aretenir}[L'essentiel de l'analyse]
Un bon sujet de dissertation contient toujours une \textbf{tension} entre deux idées opposées (thèse et antithèse potentielles). Le travail du candidat consiste à la repérer : c'est elle qui donnera naissance à la problématique, puis au plan dialectique du devoir.
\end{aretenir}

\courssub{Étape 4 : formuler la problématique}

\begin{definition}[La problématique]
La \cle{problématique} est la question principale qui naît de l'analyse du sujet et que le développement devra résoudre. Elle se formule en \textbf{une seule phrase interrogative}, claire et précise, et elle oriente l'ensemble du devoir : chaque partie, chaque argument devra contribuer à y répondre.
\end{definition}

La problématique est comme la question qu'un enquêteur se pose avant de commencer son enquête : tant qu'elle n'est pas formulée, il ne sait pas ce qu'il cherche. Dans la copie, elle apparaît généralement à la fin de l'introduction, juste avant l'annonce du plan.

\begin{propriete}[Critères d'une bonne problématique]
Une bonne problématique est :
\begin{itemize}
\item \textbf{pertinente} : elle découle directement du sujet et de ses mots-clés, sans s'en écarter ;
\item \textbf{claire} : elle est formulée en une phrase simple, compréhensible immédiatement ;
\item \textbf{ouverte} : elle appelle un débat, c'est-à-dire qu'on peut y répondre de manière nuancée (elle n'attend pas un simple « oui » ou « non » définitif).
\end{itemize}
\end{propriete}

\begin{attention}[Erreur fréquente : la fausse problématique]
Beaucoup de candidats se contentent de \textbf{répéter le sujet} sous forme de question : « La tradition est-elle un obstacle au développement ? » Cette formulation ne fait que recopier l'énoncé sans l'interroger : elle ne montre aucun travail d'analyse et n'ouvre aucun débat réel. Une vraie problématique va \emph{au-delà} du sujet : elle met en relation les termes, elle soulève un enjeu, elle interroge les présupposés.
\end{attention}

\courssub{Exemple guidé : du sujet à la problématique}

Appliquons la méthode complète à deux sujets, dont l'un est entièrement résolu.

\begin{exoresolu}[Analyse complète du sujet « La tradition est un obstacle au développement »]
Énoncé : Analysez le sujet « La tradition est un obstacle au développement » et formulez une problématique.
\tcblower
Solution détaillée :

\textbf{Étape 1 — Mots-clés} : tradition ; obstacle ; développement.

\textbf{Étape 2 — Définitions et thème} : la tradition est l'héritage culturel transmis entre générations ; un obstacle est ce qui empêche d'avancer ; le développement est le progrès matériel et social d'une société. Thème délimité : le rapport entre héritage culturel et progrès, notamment dans les sociétés africaines en mutation.

\textbf{Étape 3 — Présupposés et tension} : le sujet suppose que tradition et développement s'opposent. Or, si certaines coutumes freinent le progrès (refus de l'école pour les filles, dépenses excessives des funérailles), d'autres le servent (solidarité de la tontine, artisanat, tourisme culturel, valeurs morales). La tension est donc : la tradition freine-t-elle le progrès, ou peut-elle au contraire l'accompagner ?

\textbf{Étape 4 — Problématique} : « Dans quelle mesure la tradition constitue-t-elle un frein au développement, et ne peut-elle pas devenir, au contraire, un levier de progrès pour nos sociétés ? »

Cette problématique est pertinente (elle reprend les trois mots-clés), claire (une seule phrase) et ouverte (elle invite à discuter les deux faces du problème). Elle permettra de construire un plan dialectique : thèse (la tradition comme obstacle), antithèse (la tradition comme richesse), synthèse (un développement enraciné dans la culture).
\end{exoresolu}

\begin{exemplebox}[Un second exemple]
Sujet : « Faut-il moderniser nos villages ? »

Analyse rapide : les mots-clés sont \emph{moderniser} (transformer selon les techniques et les modes de vie actuels) et \emph{villages} (cadre de vie traditionnel, communauté). La tension : la modernisation apporte le confort (électricité, eau potable, routes, numérique), mais elle risque de détruire ce qui fait l'âme du village (liens familiaux, coutumes, autorité traditionnelle, cohésion sociale).

Problématique : « La modernisation des villages camerounais doit-elle sacrifier les traditions qui fondent leur identité, ou peut-elle s'enraciner dans le patrimoine local ? »
\end{exemplebox}

On remarque que cette problématique ne répète pas le sujet : elle le déplace vers un véritable enjeu (modernité \emph{contre} ou \emph{avec} identité), ce qui annonce déjà un débat en deux ou trois parties.

\courssub{Entraînement}

\begin{exercice}[À vous de jouer]
Pour chacun des sujets suivants, soulignez les mots-clés, définissez-les, dégagez la tension, puis formulez une problématique en une phrase interrogative :
\begin{enumerate}
\item « L'école est la clé de la réussite. Discutez. »
\item « Faut-il abandonner nos langues locales au profit du français et de l'anglais ? »
\end{enumerate}
\end{exercice}

\begin{corrige}[Exercice]
\textbf{Sujet 1} : mots-clés : \emph{école}, \emph{clé}, \emph{réussite}. Tension : l'école donne des diplômes et des savoirs, mais on connaît des réussites sans diplômes (commerçants, artisans, artistes, entrepreneurs) et des diplômés sans emploi. Problématique possible : « L'instruction scolaire est-elle la seule voie qui mène à la réussite sociale, ou d'autres chemins peuvent-ils y conduire ? »

\textbf{Sujet 2} : mots-clés : \emph{langues locales}, \emph{français}, \emph{anglais}. Tension : les langues officielles ouvrent sur le monde, l'administration et l'emploi formel, mais les langues locales portent la culture, l'identité et la pensée profonde. Problématique possible : « La promotion des langues officielles doit-elle conduire à l'abandon des langues locales, ou les unes et les autres peuvent-elles cohabiter harmonieusement dans la formation du citoyen camerounais ? »
\end{corrige}

\begin{aretenir}[Bilan de la section]
Analyser un sujet de dissertation, c'est suivre quatre étapes : \textbf{souligner} les mots-clés, \textbf{définir} et \textbf{délimiter} le thème, dégager les \textbf{présupposés} et la \textbf{tension}, puis \textbf{interroger} pour formuler la problématique. Cette question centrale, pertinente, claire et ouverte, est la boussole de tout le devoir : c'est à partir d'elle que nous pourrons, dans la prochaine étape, rechercher les idées et les exemples qui nourriront le développement.
\end{aretenir}

\coursec{La recherche des idées}

Après avoir analysé le sujet et formulé la problématique — étape que nous venons de maîtriser — il faut maintenant nourrir la démonstration. Une dissertation sans idées pertinentes, c'est une coquille vide : belle structure, mais rien à l'intérieur. Cette section vous apprend à transformer votre sujet en une réserve d'arguments solides, illustrés d'exemples précis, pour convaincre l'examinateur du Probatoire.

\courssub{Le brainstorming (remue-méninges) : libérer la pensée sans la juger}

\begin{definition}[Brainstorming]
Le \cle{brainstorming} (ou remue-méninges) est une technique de créativité collective ou individuelle qui consiste à produire le maximum d'idées en un temps limité, \emph{sans aucune censure ni critique}, pour explorer toutes les pistes possibles autour d'un sujet.
\end{definition}

L'objectif n'est pas d'écrire la dissertation, mais de vider son esprit. Toute idée, même farfelue, peut déclencher une piste utile. On note tout : mots-clés, noms propres, dates, titres d'œuvres, proverbes, faits d'actualité, souvenirs de cours.

\begin{methode}[Conduire un brainstorming efficace en 5 minutes]
\begin{enumerate}
\item \textbf{Préparez le terrain} : une feuille blanche (ou un tableau), un stylo, un minuteur réglé sur 5 minutes.
\item \textbf{Écrivez le sujet au centre} : par exemple « Tradition et modernité au Cameroun ».
\item \textbf{Laissez jaillir les idées} : notez \emph{tout} ce qui vous vient — pas de phrases complètes, des mots, des flèches, des dessins. Aucune idée n'est « bête ».
\item \textbf{Ne vous arrêtez pas} : si vous bloquez, relisez le sujet, changez d'angle (économique, culturel, politique, scolaire, familial).
\item \textbf{À la fin, ne triez pas encore} : passez directement à l'étape suivante (mobilisation des connaissances) pour enrichir la liste.
\end{enumerate}
\end{methode}

\begin{exemplebox}[Brainstorming individuel — 5 minutes sur « Tradition et modernité »]
\textbf{Idées brutes notées :} chefferies, code de la famille, mariages coutumiers, dot, téléphones portables dans les villages, TikTok, école formelle vs initiation, exode rural, startups à Douala, langues nationales vs français/anglais, médecine traditionnelle vs hôpital, fêtes traditionnelles (Ngondo, Medumba), urbanisation, perte des valeurs, génération Z, respect des aînés, individualisme, diaspora, transferts d'argent, agriculture moderne, changement climatique, chefferies 2.0 (pages Facebook), réconciliation nationale.
\end{exemplebox}

\begin{attention}[Piège fréquent : l'auto-censure]
Ne dites pas « ça ne vaut rien » ou « c'est hors-sujet » pendant le brainstorming. Une idée sur « les influenceurs TikTok au village » peut sembler anecdotique, mais elle devient un excellent exemple pour illustrer la modernité qui s'invite dans la tradition. \textbf{Tout se note, le tri vient après.}
\end{attention}

\courssub{La mobilisation des connaissances : faire le lien avec vos acquis}

Le brainstorming a vidé votre mémoire vive ; maintenant, il faut aller chercher dans la mémoire longue : cours de français, d'histoire, de géographie, d'éducation civique, de philosophie, lectures personnelles (dont \emph{Le Lion et la Perle}), actualité camerounaise et africaine.

\begin{propriete}[Réservoirs de connaissances mobilisables]
Pour un sujet de dissertation en Première technique, les réservoirs principaux sont :
\begin{itemize}
\item \textbf{Le programme de français} : textes étudiés (\emph{Le Lion et la Perle}, poèmes, articles de presse), notions de style, registres, genres littéraires.
\item \textbf{Les autres disciplines} : histoire (colonisation, indépendances), géographie (urbanisation, ressources), éducation civique (citoyenneté, droits de l'homme), économie (secteur informel, entrepreneuriat).
\item \textbf{L'actualité camerounaise et africaine} : CAN, élections, réformes éducatives, projets structurants (barrage de Nachtigal, autoroute Yaoundé-Douala), diaspora, conflits, innovations technologiques (fintech, agritech).
\item \textbf{La culture générale africaine} : écrivains (Mongo Beti, Calixthe Beyala, Léonora Miano), penseurs (Achille Mbembe, Felwine Sarr), mouvements artistiques, proverbes, contes.
\item \textbf{Votre expérience vécue} : stage en entreprise, vie associative, famille, quartier, observations de terrain.
\end{itemize}
\end{propriete}

\begin{methode}[Mobiliser ses connaissances par thème]
\begin{enumerate}
\item Reprenez chaque grosse idée issue du brainstorming.
\item Demandez-vous : « Quel cours, quel texte, quel fait précis illustre cette idée ? »
\item Notez la référence exacte : « Cours d'histoire, chapitre 3 : résistance nationale » ; « \emph{Le Lion et la Perle}, scène du marché : Baroka vs Lakunle » ; « Article \emph{Jeune Afrique} mars 2024 : startup camerounaise AgriTech ».
\item Visez \textbf{au moins 2 arguments et 2 exemples précis par partie} de votre futur plan.
\end{enumerate}
\end{methode}

\begin{exemplebox}[Mobilisation pour « Tradition et modernité » — Partie « Conflits et tensions »]
\textbf{Argument :} La modernité scolaire dévalorise les savoirs traditionnels.
\textbf{Exemple 1 (texte étudié) :} Dans \emph{Le Lion et la Perle}, Lakunle (instituteur « moderne ») méprise les coutumes (refus de payer la dot, moquerie du culte d'Ogun) et veut imposer un mariage « occidental » à Sidi.
\textbf{Exemple 2 (actualité) :} Rapport MINEDUB 2023 : 68\% des élèves du secondaire ne parlent pas leur langue maternelle ; programmes scolaires peu intégrateurs des savoirs endogènes.
\textbf{Exemple 3 (expérience) :} Lors de mon stage en zone rurale, des jeunes formés à l'école refusent d'apprendre la pharmacopée locale auprès des aînés.
\end{exemplebox}

\begin{aretenir}[Repère chiffré]
Un bon devoir de Première technique mobilise en \textbf{moyenne 2 arguments et 2 exemples précis par partie} (soit 6 à 8 arguments-exemples au total pour un plan en 3 parties). Moins de 2 arguments par partie affaiblit la démonstration ; plus de 3 risque de noyer le fil conducteur.
\end{aretenir}

\courssub{L'exploitation de documents : transformer des sources en arguments}

Au Probatoire, le sujet de dissertation peut s'accompagner d'un corpus (textes, images, graphiques, tableaux). Même sans corpus fourni, vous devez agir \emph{comme si} vous en aviez un : votre culture documentaire est votre corpus interne.

\begin{definition}[Exploitation documentaire]
L'\cle{exploitation de documents} consiste à extraire d'une source (texte, image, donnée chiffrée, carte, témoignage) une information pertinente, à la contextualiser et à l'articuler à votre argumentation, en citant la source.
\end{definition}

\begin{methode}[Exploiter un document en 4 étapes]
\begin{enumerate}
\item \textbf{Identifier} : nature (texte journalistique, roman, statistique, photo, carte), auteur, date, source, intention.
\item \textbf{Comprendre} : idée principale, données clés, ton, thèse défendue.
\item \textbf{Sélectionner} : quelle citation, quel chiffre, quel détail visuel sert \emph{votre} argument ? (Pas plus de 2 lignes de citation).
\item \textbf{Intégrer} : introduire le document (« Selon le rapport X… », « Comme le montre l'image Y… »), citer/préciser, puis \emph{commenter} : en quoi cela prouve votre point ?
\end{enumerate}
\end{methode}

\begin{exemplebox}[Exploitation d'un chiffre — « Tradition et modernité »]
\textbf{Document :} « En 2022, 42\% des Camerounais de 15–35 ans utilisent Internet régulièrement (INS, EESI 5). »
\textbf{Intégration :} « Selon l'INS (EESI 5, 2022), 42\% des jeunes Camerounais sont connectés. Ce chiffre illustre la pénétration rapide du numérique (modernité) dans une société où l'oralité et la transmission intergénérationnelle (tradition) restaient hégémoniques il y a vingt ans. »
\end{exemplebox}

\begin{attention}[Erreur fatale : l'exemple vague]
\emph{« Quelqu'un m'a dit que… »}, \emph{« On dit que… »}, \emph{« Dans mon village, y'a un gars qui… »} — \textbf{zéro point}. L'examinateur attend des faits \textbf{vérifiables, datés, sourcés}. Remplacez « on dit » par « Selon le rapport INS 2022 » ou « Comme l'illustre la scène 3 de \emph{Le Lion et la Perle} ».
\end{attention}

\courssub{Le questionnement systématique (QQOQCP) : ne rien laisser dans l'ombre}

La méthode \cle{QQOQCP} (Qui ? Quoi ? Où ? Quand ? Comment ? Pourquoi ? / Combien ?) est un filet de sécurité : elle force à explorer toutes les dimensions du sujet.

\begin{methode}[Appliquer le QQOQCP à votre problématique]
Pour chaque grande partie de votre plan, posez les 6 questions et notez les réponses courtes.
\begin{enumerate}
\item \textbf{Qui ?} Acteurs : jeunes, aînés, chefs traditionnels, État, ONG, entreprises, diaspora, écoles, médias.
\item \textbf{Quoi ?} Objets : langues, rites, technologies, lois, valeurs, savoirs, espaces, temps.
\item \textbf{Où ?} Espaces : village, ville, école, marché, chefferie, hôpital, église/mosquée, réseaux sociaux, diaspora.
\item \textbf{Quand ?} Temporalités : précolonial, colonial, indépendances, années 90 (libéralisation), ère numérique, prospective 2035.
\item \textbf{Comment ?} Mécanismes : transmission, imitation, résistance, hybridation, appropriation, conflit, négociation, législation.
\item \textbf{Pourquoi ? / Combien ?} Causes \& ampleur : démographie, économie, idéologie, pouvoir, identité, survie, fierté, rentabilité.
\end{enumerate}
\end{methode}

\begin{exemplebox}[QQOQCP appliqué à « Tradition et modernité : hybridation »]
\begin{center}
\begin{tabularx}{\linewidth}{|l|X|}\hline
\textbf{Question} & \textbf{Réponses exploitées} \\ \hline
Qui ? & Jeunes entrepreneurs (ex. \emph{KmerPad}, \emph{GiftedMom}), chefs 2.0, ONG culturelles \\ \hline
Quoi ? & Applications en langues nationales, e-commerce produits artisanaux, festivals modernisés \\ \hline
Où ? & Yaoundé, Douala, Bafoussam, mais aussi villages connectés (ex. \emph{Internet Society Cameroon}) \\ \hline
Quand ? & Depuis 2010 (essor 3G/4G), accélération post-Covid \\ \hline
Comment ? & Crowdfunding pour rites, streaming de danses, formation digitale d'artisans \\ \hline
Pourquoi ? & Valoriser l'héritage \textbf{et} gagner sa vie ; réponse au chômage, fierté identitaire \\ \hline
\end{tabularx}
\end{center}
\end{exemplebox}

\courssub{Le tri des idées : éliminer, regrouper, nommer}

Vous avez maintenant une masse d'idées brutes, de connaissances mobilisées, de documents exploités, de réponses QQOQCP. Il faut faire le ménage.

\begin{methode}[Tri en 3 passes]
\begin{enumerate}
\item \textbf{Passage 1 — Élimination (hors-sujet, doublons, trop vagues)} : Rayez ce qui ne répond pas à la problématique, ce qui est répété, ce qui reste flou (« la culture », « les gens »).
\item \textbf{Passage 2 — Regroupement (familles d'arguments)} : Collez des post-it ou tracez des cercles de couleur autour des idées qui vont ensemble. Donnez un nom à chaque famille (ce sera vos \cle{sous-parties}).
\item \textbf{Passage 3 — Vérification d'équilibre} : Chaque grande partie du plan doit contenir 2 à 3 familles d'arguments. Si une partie n'en a qu'une, elle est trop faible ; si elle en a 5, elle est trop lourde — scindez-la.
\end{enumerate}
\end{methode}

\begin{exemplebox}[Tri pour « Tradition et modernité » — Résultat]
\textbf{Familles identifiées (couleurs) :}
\begin{itemize}
\item \textcolor{popblue}{\textbf{Conflits}} : dot vs mariage civil, médecine trad. vs hôpital, école vs initiation, langue française vs langues nationales.
\item \textcolor{popgreen}{\textbf{Hybridations}} : chefferies 2.0, startups agro/artisanat, festivals modernisés, télé-médecine trad., influenceurs culturels.
\item \textcolor{poporange}{\textbf{Enjeux de pouvoir}} : législation (code de la famille), éducation (programmes), médias (contenu local), économie (brevetage savoirs trad.).
\end{itemize}
\textbf{Équilibre :} 3 familles → 3 parties parfaites. Chaque famille contient 3–4 arguments précis.
\end{exemplebox}

\courssub{La hiérarchisation : ordonner du plus faible au plus fort}

Dans chaque partie, les arguments ne se valent pas. La progression classique (et attendue au Bac) est \textbf{croissante} : on commence par l'argument le plus évident/accessible, on monte vers le plus fin/structurant, on termine par le plus percutant/ouvrant.

\begin{propriete}[Hiérarchisation des arguments dans une partie]
\begin{itemize}
\item \textbf{Niveau 1 (entrée) :} Constat factuel, exemple simple, illustration concrète. \emph{« Au village, les jeunes utilisent WhatsApp pour organiser les réunions de chefferie. »}
\item \textbf{Niveau 2 (analyse) :} Mise en relation, mécanisme, cause/effet. \emph{« Cet usage transforme la parole publique : l'aîné ne monopolise plus la parole, le débat s'écrit, s'archive, s'exporte. »}
\item \textbf{Niveau 3 (portée) :} Portée générale, enjeu systémique, ouverture. \emph{« Ainsi, le numérique ne se contente pas de moderniser la tradition : il en redistribue l'autorité, préfigurant une gouvernance locale plus horizontale. »}
\end{itemize}
\end{propriete}

\begin{attention}[Piège : l'argument « choc » en premier]
Si vous placez votre meilleure idée en ouverture, la suite paraîtra fade. \textbf{Gardez le meilleur pour la fin de chaque partie} — c'est ce que l'examinateur retiendra en corrigeant.
\end{attention}

\courssub{Le choix des exemples : précision, vérifiabilité, adaptation au destinataire}

Un argument sans exemple est une assertion ; un exemple mal choisi (vague, faux, hors-contexte) décrédibilise tout le paragraphe.

\begin{definition}[Exemple dissertationnable]
Un \cle{bon exemple} pour le Probatoire est :
\begin{itemize}
\item \textbf{Précis} : nom propre, date, lieu, chiffre, citation courte (≤ 2 lignes).
\item \textbf{Vérifiable} : source identifiable (ouvrage, article, rapport, texte au programme, observation datée).
\item \textbf{Adapté au destinataire} : l'examinateur camerounais connaît le contexte national ; pas besoin d'expliquer ce qu'est le Ngondo, mais citez l'édition 2023. Pour un fait international, contextualisez en une phrase.
\item \textbf{Varié} : alternez littéraire (\emph{Le Lion et la Perle}), statistique (INS, Banque Mondiale), historique (indépendance 1960), actualité (CAN 2021, loi 2024), expérience de stage.
\end{itemize}
\end{definition}

\begin{exoresolu}[Choisir le bon exemple — Sujet : « Tradition et modernité »]
\textbf{Argument :} La modernité juridique menace les droits coutumiers des femmes.
\tcblower
\begin{itemize}
\item \textcolor{poppink}{\textbf{Mauvais exemple :}} « Les femmes sont maltraitées par la tradition. » (Vague, pas d'exemple, jugement moral non argumenté)
\item \textcolor{poporange}{\textbf{Exemple faible :}} « Dans mon quartier, une veuve a été chassée. » (Anecdote non sourcée, non datée, non généralisable)
\item \textcolor{popgreen}{\textbf{Bon exemple :}} « L'article 732 du Code civil camerounais (héritage) entre en tension avec la coutume bamiléké (exclusion des filles) : affaire \emph{Mme X c. Famille Y}, Cour d'appel du Littoral, arrêt 2021 — la cour a reconnu le droit successoral de la fille sous réserve de la coutume. » (Précis, juridique, daté, illustre le conflit norme moderne / norme coutumière)
\item \textcolor{popblue}{\textbf{Excellent exemple (couplé) :}} « Ce conflit juridique s'incarne littérairement dans \emph{Le Lion et la Perle} : Sidi, bien que « moderne » (école, photos), subit la loi de la dot imposée par Baroka — la tradition juridico-symbolique prime sur son statut individuel. » (Lien texte programme + droit positif = double ancrage)
\end{itemize}
\end{exoresolu}

\courssub{Constituer sa banque d'arguments et d'exemples : l'outil de rédaction}

Tout ce travail (brainstorming → mobilisation → documents → QQOQCP → tri → hiérarchisation → choix exemples) aboutit à un \textbf{document de travail unique} : votre \cle{banque d'arguments}. C'est la « matière première » de votre plan détaillé.

\begin{methode}[Construire sa banque d'arguments (fiche unique par sujet)]
\begin{enumerate}
\item Sur une feuille A4 (ou fichier), tracez un tableau 4 colonnes : \textbf{Partie} | \textbf{Sous-partie (famille)} | \textbf{Argument (phrase nominale)} | \textbf{Exemple(s) précis + source}.
\item Remplissez pour les 3 parties, 2–3 sous-parties chacune, 1–2 arguments par sous-partie.
\item Cochez chaque exemple utilisé → vous verrez si vous êtes déséquilibré (tout littéraire, pas de chiffre ; tout Cameroun, pas d'Afrique ; tout passé, pas de présent).
\item Gardez cette fiche \emph{sous les yeux} pendant la rédaction du plan détaillé et de la dissertation.
\end{enumerate}
\end{methode}

\begin{exemplebox}[Extrait de banque d'arguments — « Tradition et modernité »]
\begin{center}
\begin{tabularx}{\linewidth}{|l|l|X|X|}\hline
\textbf{Partie} & \textbf{Sous-partie} & \textbf{Argument} & \textbf{Exemple(s) + source} \\ \hline
I. Conflits & A. Juridique & Dot vs mariage civil & Code civil art. 144 (mariage civil) vs coutume dot obligatoire ; \emph{Le Lion et la Perle} sc. 3 (Baroka/Lakunle) \\ \hline
I. Conflits & B. Médicale & Médecine trad. vs hôpital & Enquête MINSANTÉ 2022 : 70\% population recourt d'abord trad. ; conflit pharmacopée / brevet (affaire \emph{Prunus africana}) \\ \hline
II. Hybridations & A. Numérique & Chefferies 2.0 & Chefferie Bafut : page FB, WhatsApp conseil, streaming danses (reportage CRTV 2023) \\ \hline
II. Hybridations & B. Économique & Startups valorisant savoirs & \emph{KmerPad} (serviettes locales), \emph{GiftedMom} (suivi grossesse SMS langues nat.) — Prix AFRICA CEO 2021 \\ \hline
III. Enjeux & A. Éducatif & École vs initiation & Rapport MINEDUB 2023 : 68\% élèves ne parlent pas langue maternelle ; réforme programmes 2024 (introduction cultures nat.) \\ \hline
III. Enjeux & B. Identitaire & Diaspora & Transferts diaspora 2023 : 1 200 Mds FCFA (Banque Mondiale) ; associations \emph{Cameroun Diaspora USA} financent chefferies, écoles \\ \hline
\end{tabularx}
\end{center}
\end{exemplebox}

\begin{savaistu}[Au Probatoire (et au Baccalauréat technique)]
L'épreuve de dissertation est notée sur \textbf{la pertinence des idées autant que sur la langue} (grille officielle : 10 pts idées / 10 pts expression). Une copie avec une langue parfaite mais des arguments creux (généralités, pas d'exemples) plafonne à 8–9/20. Une copie avec quelques fautes mais des arguments solides, sourcés, hiérarchisés, peut obtenir 14–15/20. \textbf{La banque d'arguments est votre assurance-vie.}
\end{savaistu}

\begin{popfigure}
\begin{tikzpicture}[
  mindmap,
  concept color=popblue,
  level 1/.style={level distance=4.5cm, sibling angle=90},
  level 2/.style={level distance=3.2cm, sibling angle=45},
  level 3/.style={level distance=2.5cm, sibling angle=30},
  every node/.style={font=\small},
  grow cyclic
]
\node[concept, text width=2.8cm, align=center] {Tradition \&\\ Modernité}
  child[concept color=poporange] {
    node[concept, text width=2.5cm, align=center] {Conflits\\ \& Tensions}
    child[concept color=poporange!70] { node[concept, text width=2.2cm, align=center] {Juridique\\ (dot, héritage)} }
    child[concept color=poporange!70] { node[concept, text width=2.2cm, align=center] {Médicale\\ (trad. vs hôpital)} }
    child[concept color=poporange!70] { node[concept, text width=2.2cm, align=center] {Éducative\\ (école vs initiation)} }
    child[concept color=poporange!70] { node[concept, text width=2.2cm, align=center] {Linguistique\\ (français vs langues nat.)} }
  }
  child[concept color=popgreen] {
    node[concept, text width=2.5cm, align=center] {Hybridations\\ \& Innovations}
    child[concept color=popgreen!70] { node[concept, text width=2.2cm, align=center] {Chefferies 2.0\\ (FB, WhatsApp)} }
    child[concept color=popgreen!70] { node[concept, text width=2.2cm, align=center] {Startups\\ (AgriTech, FinTech)} }
    child[concept color=popgreen!70] { node[concept, text width=2.2cm, align=center] {Festivals\\ modernisés} }
    child[concept color=popgreen!70] { node[concept, text width=2.2cm, align=center] {Influenceurs\\ culturels} }
  }
  child[concept color=poppurple] {
    node[concept, text width=2.5cm, align=center] {Enjeux\\ de Pouvoir}
    child[concept color=poppurple!70] { node[concept, text width=2.2cm, align=center] {Législation\\ (Code famille)} }
    child[concept color=poppurple!70] { node[concept, text width=2.2cm, align=center] {Médias\\ (contenu local)} }
    child[concept color=poppurple!70] { node[concept, text width=2.2cm, align=center] {Économie\\ (brevetage)} }
    child[concept color=poppurple!70] { node[concept, text width=2.2cm, align=center] {Diaspora\\ (transferts)} }
  };
\end{tikzpicture}
\legende{Carte mentale issue du brainstorming + tri + hiérarchisation sur « Tradition et modernité ». Trois grandes familles (couleurs) = trois parties du plan. Chaque nœud niveau 2 = sous-partie ; niveau 3 = arguments précis avec exemples sourcés.}
\end{popfigure}

\begin{exercice}[Entraînement : votre banque d'arguments]
Choisissez un sujet type Probatoire : \emph{« Le jeune Camerounais face à l'emploi : entre attente de l'État et esprit d'entreprise »}.
\begin{enumerate}
\item Faites un brainstorming de 5 minutes (chronomètre).
\item Mobilisez 3 connaissances précises (cours, actualité, \emph{Le Lion et la Perle}).
\item Appliquez le QQOQCP à la problématique « Comment le jeune Camerounais construit-il son insertion professionnelle ? ».
\item Triez, regroupez en 3 familles, hiérarchisez.
\item Remplissez une fiche « banque d'arguments » (tableau 4 colonnes) avec au moins 2 exemples sourcés par sous-partie.
\end{enumerate}
\end{exercice}

\begin{corrige}[Exercice n°1 — Pistes de correction]
\textbf{Brainstorming (échantillon) :} chômage jeunes (INS 2022 : 13\% 15–35 ans), programmes PIAASI, PAJER-U, GIC, startup \emph{KmerPad}, \emph{Wascal}, diaspora, formation duale, secteur informel (90\% emplois), migration, formation professionnelle, incubateurs (JFN Center, ActivSpaces), agriculture jeune, foncier, financement (FNE, ANPME), mindset « fonctionnaire », valorisation métiers manuels, \emph{Le Lion et la Perle} : Lakunle (instituteur salarié) vs Baroka (chef entrepreneur terroir).
\textbf{QQOQCP :} Qui ? Jeunes, État, secteur privé, ONG, diaspora. Quoi ? Emploi, formation, financement, terre, mindset. Où ? Villes, zones rurales, incubateurs, champs, cybercafés. Quand ? 1990 (dévaluation), 2010 (PIAASI), 2020 (Covid), 2024 (PND30). Comment ? Concours, auto-emploi, GIC, startup, migration, formation duale. Pourquoi ? Démographie, inadéquation formation/emploi, culture salariat, manque capital.
\textbf{Familles :} I. Contraintes structurelles (État, formation, foncier, financement). II. Initiatives endogènes (GIC, agro-jeunes, artisanat, diaspora). III. Nouveaux modèles (incubateurs, numérique, formation duale, mindset entrepreneurial).
\textbf{Banque d'arguments :} voir structure tableau ci-dessus — adapter avec exemples précis (chiffres INS, lois 2023, noms startups, arrêts jurisprudence, scènes \emph{Le Lion et la Perle}).
\end{corrige}

\begin{aretenir}[Synthèse : les 7 clés d'une recherche d'idées réussie]
\begin{enumerate}
\item \textbf{Brainstorming} 5 min sans censure.
\item \textbf{Mobilisation} croisée : cours + \emph{Le Lion et la Perle} + actualité Cameroun/Afrique + expérience.
\item \textbf{Documents} : exploitez chiffres, citations, images comme des preuves.
\item \textbf{QQOQCP} : check-list pour ne rien oublier.
\item \textbf{Tri} : éliminez → regroupez (familles) → équilibrez (2–3 sous-parties/partie).
\item \textbf{Hiérarchisation} : faible → fort dans chaque partie.
\item \textbf{Banque d'arguments} : fiche unique, 4 colonnes, exemples sourcés, variés, précis.
\end{enumerate}
\end{aretenir}

La recherche d'idées est maintenant complète. Vous possédez la matière brute, triée, hiérarchisée, sourcée. L'étape suivante — \emph{l'élaboration du plan détaillé} — consistera à architecturer cette matière en une démonstration logique, fluide, prête à être rédigée.

\coursec{Les types de plans et l'élaboration du plan détaillé}

Dans la section précédente, nous avons appris à \cle{rechercher les idées} : à partir de la problématique, nous avons accumulé des arguments, des exemples, des références littéraires et des faits de société. Mais un tas de briques ne fait pas une maison. De même, une liste d'idées, aussi riche soit-elle, ne fait pas une dissertation. Il faut maintenant \textbf{organiser} ces idées, les \textbf{hiérarchiser} et les \textbf{ordonner} : c'est le travail du plan. Cette section vous donne les trois grands modèles de plans exigibles au Probatoire et au Baccalauréat technique, puis la méthode complète pour rédiger un \cle{plan détaillé}, véritable canevas de votre devoir.

\courssub{Les trois types de plans}

\begin{definition}[Plan]
Le \cle{plan} est l'organisation générale d'un devoir : il répartit les idées en grandes parties (généralement deux ou trois), elles-mêmes subdivisées en sous-parties, de manière à répondre de façon ordonnée et progressive à la problématique.
\end{definition}

Il existe trois types de plans classiques. Le choix de l'un ou de l'autre dépend de la nature du sujet.

\courssub{Le plan thématique}

Le \cle{plan thématique} étudie successivement les \textbf{différents aspects} d'une même question, sans les opposer. On passe d'un aspect à l'autre selon un ordre logique : du plus simple au plus complexe, du plus connu au moins connu, ou selon le schéma \emph{causes / manifestations / conséquences}.

\begin{exemplebox}[Exemple de plan thématique]
Sujet : \emph{Le paludisme au Cameroun}.

\textbf{I.} Les causes du paludisme (le moustique anophèle, les eaux stagnantes, les conditions d'hygiène).

\textbf{II.} Les manifestations de la maladie (fièvres, frissons, anémie, formes graves).

\textbf{III.} Les conséquences et les moyens de lutte (absentéisme scolaire, coût des soins, moustiquaires imprégnées, assainissement).
\end{exemplebox}

Ce plan convient aux sujets \textbf{descriptifs ou explicatifs}, c'est-à-dire aux sujets qui demandent d'\emph{exposer} et d'\emph{expliquer} un phénomène, sans débat.

\courssub{Le plan dialectique}

Le \cle{plan dialectique} est le plan du \textbf{débat}. Il convient aux sujets \textbf{polémiques}, c'est-à-dire aux sujets sur lesquels les opinions s'opposent. Il se construit en trois temps :

\begin{enumerate}
\item la \cle{thèse} : on expose d'abord le point de vue contenu dans le sujet (ou le point de vue le plus répandu), avec ses arguments ;
\item l'\cle{antithèse} : on expose ensuite le point de vue opposé, qui nuance ou contredit le premier ;
\item la \cle{synthèse} : on dépasse l'opposition en proposant une position équilibrée, personnelle et argumentée.
\end{enumerate}

\begin{attention}[La synthèse n'est pas un compromis mou]
La synthèse ne consiste pas à dire simplement « les deux ont raison ». Elle doit \emph{dépasser} l'opposition : montrer dans quelles conditions chaque point de vue est valable, ou proposer une troisième voie. Une synthèse faible ruine tout le devoir.
\end{attention}

\begin{exemplebox}[Exemple de plan dialectique]
Sujet : \emph{La tradition est un obstacle au développement. Discutez.}

\textbf{I. Thèse} : la tradition freine le développement (refus des techniques nouvelles, poids des coutumes sur les choix individuels, résistance au changement — exemple : certains rites qui retardent la scolarisation des filles).

\textbf{II. Antithèse} : la tradition peut aussi soutenir le développement (solidarité familiale et villageoise, tontines qui financent des projets, savoir-faire artisanaux et agricoles, valeurs morales qui structurent la société).

\textbf{III. Synthèse} : il faut distinguer les traditions vivantes et utiles des pratiques archaïques ; le vrai développement s'enracine dans la culture tout en s'ouvrant au progrès — comme le montre le conflit entre Baroka et Lakunle dans \emph{Le Lion et la Perle} de Wole Soyinka.
\end{exemplebox}

\courssub{Le plan analytique}

Le \cle{plan analytique} convient aux \textbf{sujets de société} qui posent un problème concret à résoudre. Il suit la démarche du médecin : observer, diagnostiquer, soigner.

\begin{enumerate}
\item le \cle{constat} : on décrit le problème, on montre qu'il existe et qu'il est grave ;
\item les \cle{causes} : on analyse les origines du problème ;
\item les \cle{solutions} : on propose des remèdes réalistes et argumentés.
\end{enumerate}

\begin{exemplebox}[Exemple de plan analytique]
Sujet : \emph{L'insécurité routière au Cameroun}.

\textbf{I.} Constat : le nombre élevé d'accidents, les victimes, le coût humain et économique.

\textbf{II.} Causes : état des routes, excès de vitesse, alcool au volant, surcharge des véhicules, laxisme des contrôles.

\textbf{III.} Solutions : entretien du réseau, répression des infractions, sensibilisation, formation des conducteurs.
\end{exemplebox}

\begin{popfigure}
\begin{center}
\begin{tikzpicture}[font=\small]
\node[draw=popblue, fill=popblueL, rounded corners, minimum width=3.4cm, minimum height=2.6cm, align=center] (thema) at (0,0) {\textbf{Plan thématique}\\[2pt] Aspects d'une\\ question\\ \emph{causes / manifestations / conséquences}};
\node[draw=poporange, fill=poporangeL, rounded corners, minimum width=3.4cm, minimum height=2.6cm, align=center] (dial) at (5.2,0) {\textbf{Plan dialectique}\\[2pt] Sujet polémique\\ \emph{thèse / antithèse / synthèse}};
\node[draw=popgreen, fill=popgreenL, rounded corners, minimum width=3.4cm, minimum height=2.6cm, align=center] (ana) at (10.4,0) {\textbf{Plan analytique}\\[2pt] Problème de société\\ \emph{constat / causes / solutions}};
\node[below=0.15cm of thema, align=center, text width=3.6cm] {\footnotesize Ex. : \emph{Le paludisme au Cameroun}};
\node[below=0.15cm of dial, align=center, text width=3.6cm] {\footnotesize Ex. : \emph{La tradition est un obstacle au développement}};
\node[below=0.15cm of ana, align=center, text width=3.6cm] {\footnotesize Ex. : \emph{L'insécurité routière}};
\end{tikzpicture}
\end{center}
\legende{Tableau comparatif des trois types de plans : structure, usage et exemple de sujet.}
\end{popfigure}

\courssub{Choisir le bon plan}

\begin{methode}[Comment choisir son plan]
\begin{enumerate}
\item \textbf{Relire la consigne} et repérer le verbe directeur : \emph{discutez}, \emph{que pensez-vous} annoncent un débat ; \emph{expliquez}, \emph{présentez} annoncent un exposé ; \emph{comment résoudre}, \emph{quelles solutions} annoncent un problème à traiter.
\item \textbf{Identifier la nature du sujet} : sujet polémique (opinions opposées) $\rightarrow$ plan \cle{dialectique} ; sujet descriptif ou explicatif $\rightarrow$ plan \cle{thématique} ; problème à résoudre $\rightarrow$ plan \cle{analytique}.
\item \textbf{Vérifier la matière disponible} : le plan choisi doit pouvoir accueillir les idées rassemblées lors de la recherche des idées. Si vos idées se répartissent naturellement en « pour » et « contre », le dialectique s'impose.
\item \textbf{Tester le plan} : formuler chaque grande partie en une phrase et vérifier que l'ensemble répond exactement à la problématique.
\end{enumerate}
\end{methode}

\courssub{L'architecture du plan détaillé}

\begin{definition}[Plan détaillé]
Le \cle{plan détaillé} est le canevas complet du devoir : il indique, pour chaque partie et chaque sous-partie, les \textbf{arguments}, les \textbf{exemples} qui les illustrent et les \textbf{transitions} qui relient les parties entre elles. C'est la dernière étape avant la rédaction.
\end{definition}

Le plan détaillé respecte l'architecture classique de la dissertation : une introduction, un développement en deux ou trois parties, une conclusion.

\textbf{1. L'introduction} comporte quatre éléments, dans l'ordre :
\begin{itemize}
\item l'\cle{amorce} (ou accroche) : une phrase générale qui attire l'attention du lecteur et situe le thème ;
\item la \cle{présentation du sujet} : on définit les termes importants et on situe la question ;
\item la \cle{problématique} : la ou les questions directrices, souvent formulées par un interrogatif ;
\item l'\cle{annonce du plan} : on présente brièvement les grandes parties (« Nous verrons d'abord que..., puis que... »).
\end{itemize}

\textbf{2. Le développement} comporte deux ou trois grandes parties. Chaque partie contient :
\begin{itemize}
\item deux ou trois \textbf{sous-parties}, chacune portant un argument ;
\item pour chaque argument, au moins un \textbf{exemple} précis (fait, statistique, référence littéraire comme \emph{Le Lion et la Perle}, expérience de la vie courante) ;
\item une \textbf{transition} à la fin de chaque partie, qui fait le bilan de ce qui vient d'être dit et annonce la partie suivante.
\end{itemize}

\textbf{3. La conclusion} comporte deux éléments :
\begin{itemize}
\item le \cle{bilan} : on répond clairement à la problématique en résumant la position adoptée, sans répéter mot à mot le développement ;
\item l'\cle{ouverture} : on élargit la réflexion vers une question voisine, sans traiter un nouveau sujet.
\end{itemize}

\begin{propriete}[Équilibre du plan]
Un plan équilibré comporte des parties de \textbf{longueur comparable} (aucune partie ne doit être trois fois plus longue qu'une autre), et chaque partie est subdivisée en \textbf{deux ou trois sous-parties}. Une partie à une seule sous-partie est mal construite : si l'on ne peut pas la diviser, c'est qu'elle est trop pauvre ou mal délimitée.
\end{propriete}

\begin{popfigure}
\begin{center}
\begin{tikzpicture}[font=\small]
% Pyramide
\fill[popgoldL] (0,4.6) -- (-1.1,3.4) -- (1.1,3.4) -- cycle;
\fill[popblueL] (-1.1,3.3) -- (1.1,3.3) -- (1.9,1.9) -- (-1.9,1.9) -- cycle;
\fill[poporangeL] (-2.0,1.8) -- (2.0,1.8) -- (2.8,0.4) -- (-2.8,0.4) -- cycle;
\fill[popgreenL] (-2.9,0.3) -- (2.9,0.3) -- (3.7,-1.1) -- (-3.7,-1.1) -- cycle;
\draw[popdark, thick] (0,4.6) -- (-3.7,-1.1) (0,4.6) -- (3.7,-1.1);
\draw[popdark] (-1.1,3.4) -- (1.1,3.4) (-1.95,1.85) -- (1.95,1.85) (-2.85,0.35) -- (2.85,0.35);
\node at (0,3.9) {\footnotesize \textbf{Introduction}};
\node at (0,2.6) {\footnotesize \textbf{Partie I}};
\node at (0,1.1) {\footnotesize \textbf{Partie II}};
\node at (0,-0.4) {\footnotesize \textbf{Conclusion}};
% Annotations
\node[right, align=left, text width=5.2cm] at (4.0,3.9) {\footnotesize amorce, présentation du sujet, problématique, annonce du plan};
\node[right, align=left, text width=5.2cm] at (4.0,2.6) {\footnotesize 2 ou 3 sous-parties : argument + exemple ; transition};
\node[right, align=left, text width=5.2cm] at (4.0,1.1) {\footnotesize 2 ou 3 sous-parties : argument + exemple ; transition};
\node[right, align=left, text width=5.2cm] at (4.0,-0.4) {\footnotesize bilan (réponse à la problématique) + ouverture};
\draw[->, popmuted] (3.6,3.9) -- (1.2,3.9);
\draw[->, popmuted] (3.6,2.6) -- (1.6,2.6);
\draw[->, popmuted] (3.6,1.1) -- (2.2,1.1);
\draw[->, popmuted] (3.6,-0.4) -- (3.0,-0.4);
\end{tikzpicture}
\end{center}
\legende{Architecture du plan détaillé : l'introduction est brève, les parties du développement forment le corps du devoir, la conclusion referme l'ensemble.}
\end{popfigure}

\courssub{Connecteurs logiques et transitions}

Les \cle{connecteurs logiques} sont les articulations du devoir : ils montrent au correcteur que la pensée progresse. Il faut les choisir selon la relation logique :

\begin{center}
\begin{tabularx}{\linewidth}{|l|X|}
\hline
\rowcolor{popblueL} \textbf{Relation} & \textbf{Connecteurs utiles} \\
\hline
Addition & d'abord, ensuite, de plus, en outre, enfin \\
\hline
Opposition & mais, cependant, toutefois, en revanche, néanmoins \\
\hline
Cause & car, parce que, en effet, puisque \\
\hline
Conséquence & donc, ainsi, c'est pourquoi, par conséquent \\
\hline
Illustration & par exemple, ainsi, c'est le cas de \\
\hline
Conclusion & en définitive, en somme, pour conclure \\
\hline
\end{tabularx}
\end{center}

La \cle{transition} est une phrase (ou deux) placée entre deux grandes parties. Elle comporte deux temps : un \textbf{bilan} de la partie qui se termine, puis une \textbf{annonce} de la partie suivante. Exemple : « Si la tradition apparaît donc comme un frein au changement, il serait pourtant injuste de n'y voir qu'un obstacle : elle peut aussi servir de moteur au développement, comme nous allons le montrer. »

\begin{attention}[Erreurs fréquentes de construction]
\begin{itemize}
\item Construire deux parties qui \textbf{se chevauchent} : les mêmes arguments reviennent dans les deux parties, ce qui prouve que le découpage est mauvais. Chaque partie doit avoir son domaine propre.
\item Construire des parties qui \textbf{ne répondent pas à la problématique} : on parle du sujet, mais on oublie la question posée. Chaque partie doit apporter un élément de réponse.
\item Annoncer un plan à deux parties et en rédiger trois, ou l'inverse : l'annonce du plan engage le correcteur, il faut la respecter.
\end{itemize}
\end{attention}

\courssub{Vérification finale du plan}

Avant de rédiger, relisez votre plan détaillé en vous posant quatre questions :

\begin{methode}[La grille de vérification]
\begin{enumerate}
\item \textbf{Cohérence} : les parties s'enchaînent-elles logiquement ? Les transitions fonctionnent-elles ?
\item \textbf{Équilibre} : les parties sont-elles de longueur comparable, chacune avec deux ou trois sous-parties ?
\item \textbf{Fidélité à la problématique} : chaque partie répond-elle à la question posée ? Le sujet est-il traité entièrement, et lui seul ?
\item \textbf{Richesse} : chaque argument est-il accompagné d'au moins un exemple précis ?
\end{enumerate}
\end{methode}

\begin{exoresolu}[Élaborer un plan détaillé complet]
Énoncé : à partir du sujet « \emph{Le téléphone portable : un danger pour la jeunesse ?} », indiquez le type de plan adapté, puis rédigez un plan détaillé (introduction, grandes parties avec sous-parties et exemples, conclusion).
\tcblower
\textbf{Solution détaillée.}

\emph{Étape 1 — Choix du plan.} Le sujet est formulé comme une question polémique (« un danger ? ») : on peut répondre oui ou non. Le plan adapté est donc le plan \textbf{dialectique}.

\emph{Étape 2 — Plan détaillé.}

\textbf{Introduction.} Amorce : le téléphone portable est devenu l'objet le plus répandu chez les jeunes Camerounais. Présentation : outil de communication, d'information et de divertissement, il suscite pourtant des inquiétudes. Problématique : le téléphone portable représente-t-il un danger pour la jeunesse, ou peut-il être un instrument de progrès ? Annonce du plan : nous montrerons d'abord les dangers réels du portable, puis ses apports, avant de dégager les conditions d'un usage maîtrisé.

\textbf{I. Thèse : le portable expose la jeunesse à de réels dangers.}
\begin{itemize}
\item a) Danger pour les études : distraction en classe, baisse de concentration (exemple : élèves qui consultent les réseaux sociaux pendant les révisions).
\item b) Danger pour la santé et les mœurs : sommeil réduit, exposition à des contenus violents ou immoraux, cyberharcèlement.
\item c) Danger financier et social : coût des forfaits, vols, dépendance (exemple : jeunes qui dépensent leur argent de poche en connexion).
\end{itemize}
Transition : ces dangers sont indéniables ; faut-il pour autant condamner le téléphone portable ?

\textbf{II. Antithèse : le portable est aussi un outil précieux pour la jeunesse.}
\begin{itemize}
\item a) Outil d'apprentissage : accès aux cours en ligne, dictionnaires, groupes de travail scolaire (exemple : élèves qui révisent le Probatoire grâce à des applications éducatives).
\item b) Outil de communication et de sécurité : lien avec la famille, appels d'urgence.
\item c) Outil d'insertion économique : mobile money, petits commerces, information sur les offres d'emploi.
\end{itemize}
Transition : le portable n'est donc ni bon ni mauvais en soi ; tout dépend de l'usage qu'on en fait.

\textbf{III. Synthèse : pour un usage maîtrisé du téléphone portable.}
\begin{itemize}
\item a) Responsabiliser les jeunes : éducation aux médias, règles d'usage à la maison et au lycée.
\item b) Encadrer sans interdire : l'école et la famille doivent fixer des limites tout en valorisant les usages positifs.
\end{itemize}

\textbf{Conclusion.} Bilan : le portable n'est un danger que lorsqu'il est mal utilisé ; bien encadré, il devient un instrument de réussite. Ouverture : cette question renvoie à celle, plus large, de la place des technologies numériques dans l'éducation en Afrique.
\end{exoresolu}

\begin{exercice}[À vous de jouer]
Pour chacun des sujets suivants, indiquez le type de plan le plus adapté et rédigez les intitulés des grandes parties : 1) \emph{Les causes et les conséquences de l'exode rural au Cameroun} ; 2) \emph{Faut-il interdire les réseaux sociaux aux élèves ?} ; 3) \emph{Comment lutter contre la déforestation ?}
\end{exercice}

\begin{corrige}[Exercice]
1) Sujet explicatif $\rightarrow$ plan \textbf{thématique} : I. les causes de l'exode rural ; II. ses manifestations ; III. ses conséquences sur les campagnes et les villes. 2) Sujet polémique $\rightarrow$ plan \textbf{dialectique} : I. les raisons d'interdire ; II. les limites d'une interdiction ; III. synthèse : un usage encadré plutôt qu'une interdiction. 3) Problème à résoudre $\rightarrow$ plan \textbf{analytique} : I. constat de la déforestation ; II. causes (exploitation, agriculture, bois de chauffe) ; III. solutions (reboisement, énergies alternatives, législation).
\end{corrige}

\begin{aretenir}[L'essentiel]
Trois types de plans : \cle{thématique} (aspects d'une question), \cle{dialectique} (thèse / antithèse / synthèse, pour les sujets polémiques), \cle{analytique} (constat / causes / solutions, pour les problèmes de société). Le \cle{plan détaillé} est le canevas complet du devoir : introduction (amorce, présentation, problématique, annonce), développement (parties équilibrées, sous-parties avec arguments et exemples, transitions), conclusion (bilan et ouverture). Avant de rédiger, on vérifie : cohérence, équilibre, fidélité à la problématique.
\end{aretenir}

\coursec{Langue : conditions et situations de communication}

Après avoir maîtrisé l'élaboration du plan détaillé, pierre angulaire de la dissertation, nous devons porter notre attention sur l'outil même de notre expression : la langue. Rédiger une dissertation, c'est communiquer avec un destinataire précis — le correcteur — dans une situation codifiée — l'examen. Comprendre les mécanismes de la communication, ses conditions de réussite et les normes qui régissent ses différentes situations, c'est s'assurer que notre message (la dissertation) sera reçu, compris et évalué conformément à nos intentions. C'est l'objet de cette section.

\courssub{Qu'est-ce que communiquer ? Définition et modèle théorique}

Communiquer, ce n'est pas seulement parler ou écrire : c'est \cle{établir une relation} entre un sujet qui émet et un sujet qui reçoit, au moyen d'un système de signes partagé.

\begin{definition}[La communication]
La \textbf{communication} est un processus dynamique d'échange d'informations par lequel un \textbf{émetteur} transmet un \textbf{message} à un \textbf{récepteur} (ou destinataire), en utilisant un \textbf{code} commun, à travers un \textbf{canal} matériel, en référence à un \textbf{référent} (ou contexte) partagé.
\end{definition}

Pour analyser finement ce processus, le linguiste Roman Jakobson a proposé un modèle schématique incontournable, qui met en évidence six facteurs indispensables, chacun correspondant à une fonction du langage.

\begin{popfigure}
\begin{tikzpicture}[
    node distance=1.8cm and 2.5cm,
    box/.style={draw=popblue, thick, rounded corners, fill=popblueL, minimum width=2.8cm, minimum height=1.2cm, align=center, font=\sffamily\small},
    arrow/.style={-Stealth, thick, popdark},
    labelarrow/.style={midway, fill=white, inner sep=2pt, font=\scriptsize\sffamily, text=popdark}
]
% Nodes
\node[box, align=center] (emetteur){\textbf{ÉMETTEUR}\\(Locuteur / Écrivain)};
\node[box, right=of emetteur, align=center] (message){\textbf{MESSAGE}\\(Énoncé)};
\node[box, right=of message, align=center] (recepteur){\textbf{RÉCEPTEUR}\\(Destinataire / Lecteur)};
\node[box, below=of message, align=center] (code){\textbf{CODE}\\(Langue, dialecte, jargon)};
\node[box, below left=of code, align=center] (canal){\textbf{CANAL}\\(Air, papier, écran, onde)};
\node[box, below right=of code, align=center] (referent){\textbf{RÉFÉRENT}\\(Contexte, objet du discours)};

% Arrows
\draw[arrow] (emetteur) -- node[labelarrow] {encode} (message);
\draw[arrow] (message) -- node[labelarrow] {décode} (recepteur);
\draw[arrow] (message) -- node[labelarrow, left] {utilise} (code);
\draw[arrow] (message) -- node[labelarrow, right] {passe par} (canal);
\draw[arrow] (message) -- node[labelarrow, above] {porte sur} (referent);

% Functions labels
\node[align=center, font=\scriptsize\sffamily\itshape, text=poppurple, above=0.3cm of emetteur] {Fonction \textbf{Émotive / Expressive}};
\node[align=center, font=\scriptsize\sffamily\itshape, text=poppurple, above=0.3cm of message] {Fonction \textbf{Poétique}};
\node[align=center, font=\scriptsize\sffamily\itshape, text=poppurple, above=0.3cm of recepteur] {Fonction \textbf{Conative / Appel}};
\node[align=center, font=\scriptsize\sffamily\itshape, text=poppurple, below=0.3cm of code] {Fonction \textbf{Métalinguistique}};
\node[align=center, font=\scriptsize\sffamily\itshape, text=poppurple, below=0.3cm of canal] {Fonction \textbf{Phatique}};
\node[align=center, font=\scriptsize\sffamily\itshape, text=poppurple, below=0.3cm of referent] {Fonction \textbf{Référentielle}};
\end{tikzpicture}
\legende{Schéma de la communication selon Jakobson : les six facteurs et les six fonctions du langage.}
\end{popfigure}

\begin{propriete}[Les six fonctions du langage (Jakobson)]
Chaque facteur de la communication oriente le message vers une fonction dominante :
\begin{itemize}
    \item \textbf{Référentielle} (ou informative) : centrée sur le \textbf{référent}. Informe sur un fait, une réalité. \emph{Ex : « Le Probatoire débute le 3 juin. »}
    \item \textbf{Émotive (ou expressive)} : centrée sur l'\textbf{émetteur}. Exprime des sentiments, des émotions. \emph{Ex : « Je suis ravi de ce résultat ! »}
    \item \textbf{Conative (ou appel)} : centrée sur le \textbf{récepteur}. Vise à provoquer une action, une réaction. \emph{Ex : « Fermez la porte, s'il vous plaît. »}
    \item \textbf{Phatique} : centrée sur le \textbf{canal}. Vérifie, ouvre ou maintient le contact. \emph{Ex : « Allô ? Vous m'entendez ? »}
    \item \textbf{Métalinguistique} : centrée sur le \textbf{code}. Parle de la langue elle-même. \emph{Ex : « Le mot "dissertation" vient du latin \emph{dissertare}. »}
    \item \textbf{Poétique} : centrée sur le \textbf{message} pour lui-même. Recherche esthétique, jeu sur la forme. \emph{Ex : « Un corbeau, sur un arbre perché... »}
\end{itemize}
\end{propriete}

\begin{attention}[Piège fréquent : Code vs Canal]
Ne confondez jamais le \textbf{code} (le système de signes : français, anglais, bassa, morse, langage SMS) et le \textbf{canal} (le support physique : air/voix, papier, écran, fil téléphonique, fibre optique).
\begin{itemize}
    \item \emph{Erreur :} « Le canal de ce SMS est le français. »
    \item \emph{Correct :} « Le \textbf{code} de ce SMS est le français (avec abréviations) ; son \textbf{canal} est l'écran du téléphone (support numérique). »
\end{itemize}
Au Probatoire, une question sur « le support du message » attend le \textbf{canal} ; une question sur « la langue employée » attend le \textbf{code}.
\end{attention}

\courssub{Les conditions de la communication : les six facteurs en action}

Pour qu'un acte de communication \textbf{réussisse} (que le message soit compris comme l'émetteur l'entendait), six conditions doivent être réunies simultanément. L'absence ou le dysfonctionnement de l'une d'elles entraîne une \textbf{rupture de communication}.

\begin{aretenir}[Check-list des conditions de réussite]
\begin{enumerate}
    \item \textbf{Un Émetteur} : source de l'intention communicative. Il doit avoir une intention claire.
    \item \textbf{Un Récepteur} : cible du message. Il doit être en état de recevoir (attention, capacité cognitive, disponibilité).
    \item \textbf{Un Message} : contenu informatif structuré. Il doit être cohérent et pertinent.
    \item \textbf{Un Code commun} : système de règles partagé (langue, terminologie technique, conventions). \textbf{Condition sine qua non :} émetteur et récepteur doivent maîtriser le même code.
    \item \textbf{Un Canal fonctionnel} : support physique intact (voix audible, écrit lisible, connexion stable).
    \item \textbf{Un Référent (Contexte) partagé} : connaissance du monde, du sujet, de la situation. Sans contexte commun, le message est ambigu.
\end{enumerate}
\end{aretenir}

\begin{exemplebox}[Analyse d'une rupture de communication]
\textbf{Situation :} Un enseignant dicte une consigne complexe en français technique à un élève de Première qui n'a pas le vocabulaire spécialisé, dans une cour de récréation bruyante.
\begin{itemize}
    \item \textbf{Code :} Partiellement partagé (français commun oui, jargon technique non). $\rightarrow$ \emph{Échec du décodage.}
    \item \textbf{Canal :} Perturbé par le bruit (bruit parasite). $\rightarrow$ \emph{Perte du signal.}
    \item \textbf{Récepteur :} En situation de distraction. $\rightarrow$ \emph{Attention défaillante.}
\end{itemize}
\textbf{Conclusion :} La communication échoue par accumulation de dysfonctionnements sur trois facteurs.
\end{exemplebox}

\courssub{Les situations de communication : typologie et normes}

Une \textbf{situation de communication} est l'ensemble des circonstances concrètes (spatio-temporelles, sociales, psychologiques) qui entourent l'acte de communiquer. Elle détermine les \textbf{normes} (règles implicites ou explicites) que les interlocuteurs doivent respecter.

\begin{definition}[Situation de communication]
C'est le \textbf{contexte global} — lieu, moment, identité et relation entre interlocuteurs, sujet traité — qui donne du sens à l'échange et impose des contraintes de forme (registre, politesse, structure) et de fond.
\end{definition}

On classe habituellement les situations selon quatre axes majeurs qui se combinent :

\begin{center}
\begin{tabularx}{\linewidth}{|l|X|X|}\hline
\rowcolor{popblueL}
\textbf{Axe} & \textbf{Pôle 1} & \textbf{Pôle 2} \\ \hline
\textbf{Degré de formalité} & \textbf{Formelle} : cadres officiels, institutions, examen, administration. Normes strictes, langue soutenue, politesse codifiée. & \textbf{Informelle} : vie privée, amis, famille. Normes souples, langue familière/courante, tolérance à l'implicite. \\ \hline
\textbf{Public / Privé} & \textbf{Publique} : espace ouvert, auditoire large ou indéterminé (conférence, tribune, réseaux sociaux publics). Responsabilité engagée. & \textbf{Privée} : espace clos, interlocuteurs identifiés et limités (lettre personnelle, conversation intime). Confidentialité présumée. \\ \hline
\textbf{Mode} & \textbf{Orale} : présence physique (généralement), instantanéité, prosodie (intonation, geste), éphémère, interaction immédiate. & \textbf{Écrite} : distance spatio-temporelle, différé, révisable, trace durable, monologue apparent (pas de feed-back immédiat). \\ \hline
\textbf{Relation interpersonnelle} & \textbf{Symétrique} : pairs (élève-élève, collègues). Réciprocité des droits linguistiques. & \textbf{Asymétrique} : hiérarchie (prof-élève, admin-usager, parent-enfant). Marques de respect et de déférence obligatoires. \\ \hline
\end{tabularx}
\end{center}

\begin{exemplebox}[Deux situations, deux normes : Le palabre vs L'exposé]
\begin{center}
\begin{tabularx}{\linewidth}{|l|X|X|}\hline
\rowcolor{poporangeL}
\textbf{Critère} & \textbf{Palabre villageois (Bafoussam)} & \textbf{Exposé en classe de Première} \\ \hline
\textbf{Lieu / Moment} & Place publique / case à palabres ; temps long, cyclique. & Salle de classe ; temps limité (10-15 min), horaire fixe. \\ \hline
\textbf{Interlocuteurs} & Notables, chefs, population (asymétrie forte, respect de l'âge/statut). & Professeur (évaluateur) + pairs (public). Asymétrie prof/élève. \\ \hline
\textbf{Sujet} & Affaires de la communauté, conflits, coutumes. & Sujet d'étude programmé (littéraire, scientifique, technique). \\ \hline
\textbf{Code / Registre} & Langue locale (Ghomala', Bamiléké) ou français ; proverbes, paraboles, style oratoire. & Français standard / soutenu ; vocabulaire précis, structure logique (intro/développement/concl). \\ \hline
\textbf{Canal} & Voix (orale), gestuelle, silence signifiant. & Voix + support visuel (tableau, diaporama) ; écrit (notes). \\ \hline
\textbf{Normes majeures} & Prise de parole codifiée (droit d'aînesse), écoute respectueuse, consensus. & Levée de main, argumentation rationnelle, respect du temps, objectivité. \\ \hline
\end{tabularx}
\end{center}
\end{exemplebox}

\begin{savaistu}[Le Cameroun : champion du « code-switching » situationnel]
Au Cameroun, un même locuteur navigue dans une journée entre \textbf{trois, quatre, voire cinq codes linguistiques} selon la situation :
\begin{itemize}
    \item \textbf{Famille / Quartier :} Langue maternelle (Bassa, Douala, Bamiléké, Fulfulde...) ou \textbf{Camfranglais} (argot urbain mixte).
    \item \textbf{Marché / Transport :} \textbf{Pidgin English} (zone anglophone) ou \textbf{Français populaire} / Camfranglais.
    \item \textbf{Administration / École / Travail formel :} \textbf{Français standard} (zone francophone) ou \textbf{English standard} (zone anglophone) — \emph{code officiel, normé}.
    \item \textbf{Église / Mosquée / Cérémonie traditionnelle :} Latin (parfois), Arabe, ou langue vernaculaire sacralisée.
    \item \textbf{Médias / Réseaux sociaux :} Français standard, Anglais, Pidgin, Camfranglais, émojis (code visuel).
\end{itemize}
Cette compétence \textbf{plurilingue situationnelle} est une richesse cognitive : elle apprend à adapter son registre, son vocabulaire et sa structure à l'interlocuteur et au lieu. C'est exactement la compétence exigée en dissertation : \textbf{basculer en « mode examen » (code soutenu, structure rigide, destinataire correcteur)}.
\end{savaistu}

\courssub{Application : Analyser la communication dans \emph{Le Lion et la Perle}}

Wole Soyinka, dans \emph{Le Lion et la Perle}, met en scène des affrontements communicationnels où le choc des situations (tradition vs modernité, oral vs écrit, pouvoir vs séduction) structure l'intrigue. Analysons deux dyades majeures.

\begin{exoresolu}[Analyse de la situation de communication : Baroka vs Sidi (La scène de la chambre à coucher)]
\textbf{Énoncé :} Identifiez les facteurs de la communication et la fonction dominante lors de la séduction de Sidi par Baroka (Acte 2, scène « La chambre de Baroka »).

\tcblower

\textbf{Solution détaillée :}
\begin{enumerate}
    \item \textbf{Émetteur :} \textbf{Baroka} (le Bale, chef traditionnel, 62 ans). Statut : pouvoir politique, autorité coutumière, expérience rhétorique.
    \item \textbf{Récepteur :} \textbf{Sidi} (la « Perle », jeune, belle, instruite à l'école occidentale, vaniteuse). Statut : sujette, mais fière de sa beauté (capital symbolique).
    \item \textbf{Relation :} \textbf{Asymétrique forte} (Chef / Sujette ; Homme âgé / Jeune femme ; Tradition / Modernité). Baroka détient l'autorité institutionnelle ; Sidi détient le pouvoir de la jeunesse et de l'image (la photo).
    \item \textbf{Code :} \textbf{Yoruba (langue vernaculaire) + Proverbes + Paraboles}. Baroka utilise le code de la tradition orale, métaphorique, imagée. Sidi comprend le code mais le rejette idéologiquement (elle préfère le « moderne »).
    \item \textbf{Canal :} \textbf{Oral direct} (face-à-face), renforcé par le \textbf{non-verbal} (gestes, regards, posture, offrande de boisson, manipulation de la machine à timbres).
    \item \textbf{Référent / Contexte :} La chambre privée du Bale (espace intime, pouvoir), le soir. Sujet : le mariage / la possession de Sidi / la modernité (machine à timbres).
    \item \textbf{Message :} « Je ne suis pas un vieillard fini ; je m'adapte (machine) ; tu seras ma favorite, reine du palais. »
    \item \textbf{Fonction dominante :} \textbf{CONATIVE (Appel)}. Baroka \textbf{veut obtenir un « oui »} (le mariage, la soumission consentie). Il use de rhétorique persuasive (arguments logiques + séduction affective + promesse de pouvoir).
    \item \textbf{Fonction secondaire :} \textbf{ÉMOTIVE} (il feint la faiblesse, la jalousie pour attendrir) et \textbf{RÉFÉRENTIELLE} (il donne des infos sur la machine, le timbre).
\end{enumerate}
\textbf{Analyse critique :} Baroka \textbf{maîtrise la situation}. Il adapte son code (proverbes) à son récepteur (Sidi connaît les proverbes), utilise le canal oral (sa force), contrôle le référent (il introduit la modernité pour la neutraliser). Il gagne la partie communicationnelle.
\end{exoresolu}

\begin{exoresolu}[Analyse de la situation : Lakunle vs Sidi (La scène de l'école / Le matin)]
\textbf{Énoncé :} Comparez la stratégie communicationnelle de Lakunle à celle de Baroka face à Sidi.

\tcblower

\textbf{Solution détaillée :}
\begin{itemize}
    \item \textbf{Émetteur :} \textbf{Lakunle} (instituteur, « moderne », semi-instruit, dogmatique).
    \item \textbf{Code :} \textbf{Anglais standard (dans le texte original) + Vocabulaire livresque, citations mal digérées, refus du proverbe}. Code « scolaire », déconnecté du vécu villageois.
    \item \textbf{Canal :} Oral (souvent en public, place du village ou devant l'école), gestuelle caricaturale (mime le « civilisé »).
    \item \textbf{Relation :} \textbf{Asymétrique revendiquée par Lakunle} (Il se croit supérieur : « civilisé » vs « païenne »). Mais Sidi \textbf{renverse l'asymétrie} par le mépris et le rire.
    \item \textbf{Référent :} La « dot » (coutume), la modernité (monogamie, mariage civil), le livre (Bible/Code civil).
    \item \textbf{Fonction dominante :} \textbf{RÉFÉRENTIELLE (désir d'informer/convaincre par la logique)} + \textbf{MÉTALINGUISTIQUE (corrige le langage de Sidi)}.
    \item \textbf{Échec :} Lakunle \textbf{ignore la situation}. Il parle \emph{sur} Sidi, pas \emph{à} Sidi. Il viole les normes locales (refus de la dot = insulte à la famille), utilise un code inadapté (jargon incompréhensible ou arrogant), et son canal public humilie le récepteur.
\end{itemize}
\textbf{Conclusion :} Baroka \textbf{s'adapte à la situation} (contexte, récepteur, code, canal) ; Lakunle \textbf{s'impose à la situation} et échoue. \textbf{Leçon pour la dissertation :} Comme Baroka, le candidat doit s'adapter à la situation d'examen (correcteur, consignes, code soutenu, structure attendue), non comme Lakunle qui impose sa propre logique inadaptée.
\end{exoresolu}

\courssub{Lien direct avec la dissertation : La situation d'examen}

Rédiger une dissertation au Probatoire, c'est un acte de communication \textbf{hautement codifié}. Le candidat (émetteur) qui réussit est celui qui modélise parfaitement la situation de réception (le correcteur / récepteur).

\begin{methode}[Adapter son écrit à la situation d'examen (Modélisation du correcteur)]
\begin{enumerate}
    \item \textbf{Identifier le Récepteur :} C'est un \textbf{correcteur professionnel}, fatigué, pressé, qui applique une \textbf{grille de critères} (hors-sujet, analyse, plan, connaissances, langue). Il ne vous connaît pas. \textbf{Objectif :} Faciliter sa lecture, ne pas l'obliger à deviner.
    \item \textbf{Choisir le Code :} \textbf{Français standard soutenu}. Zéro familier, zéro SMS, zéro anglicisme non intégré, zéro approximation syntaxique. Vocabulaire précis (connecteurs logiques, termes d'analyse).
    \item \textbf{Maîtriser le Canal :} \textbf{Copie manuscrite}. Lisibilité impérative (calligraphie, marges, saut de lignes pour paragraphes, encadrement du plan). Rature propre (trait unique). Le canal \textbf{ne doit pas parasiter le message}.
    \item \textbf{Construire le Message (le texte) :} Structure \textbf{rigide} (Introduction $\rightarrow$ Développement équilibré $\rightarrow$ Conclusion). Argumentation \textbf{explicite} (thèses, arguments, exemples). Pas d'implicite culturel (contrairement au palabre).
    \item \textbf{Gérer le Référent (Sujet) :} Traiter \textbf{exactement} le sujet posé (mots-clés, instruction). Pas de hors-sujet (référent dévié). La problématique ancre le référent.
    \item \textbf{Fonction dominante visée :} \textbf{RÉFÉRENTIELLE} (démontrer, analyser) + \textbf{POÉTIQUE} (qualité de l'expression, élégance syntaxique) pour valoriser le message. \textbf{Pas de fonction Émotive} (pas de « Je pense que », « À mon avis », pathos personnel).
\end{enumerate}
\end{methode}

\begin{aretenir}[La « Situation Dissertation » en résumé]
\begin{center}
\begin{tabularx}{\linewidth}{|l|X|}\hline
\rowcolor{popgreenL}
\textbf{Facteur} & \textbf{Configuration « Probatoire »} \\ \hline
\textbf{Émetteur} & Candidat (anonyme, numéro de table) \\ \hline
\textbf{Récepteur} & Correcteur (expert, évaluateur, destinataire unique) \\ \hline
\textbf{Code} & Français standard soutenu, terminologie littéraire/grammaticale précise \\ \hline
\textbf{Canal} & Copie papier, manuscrite, lisible, structurée visuellement \\ \hline
\textbf{Référent} & Sujet du devoir (corpus + instruction / sujet d'invention) \\ \hline
\textbf{Message} & Dissertation structurée (Intro / Plan détaillé / Concl) \\ \hline
\textbf{Fonction dominante} & Référentielle (argumentation) + Poétique (qualité langue) \\ \hline
\textbf{Normes} & Respect du temps, de la consigne, de la présentation, de la correction \\ \hline
\end{tabularx}
\end{center}
\end{aretenir}

\begin{exercice}[Entraînement : Analyser une situation de communication]
\textbf{Consigne :} Pour chacune des situations ci-dessous, identifiez les 6 facteurs de Jakobson, la fonction dominante du langage et la norme de registre attendue.
\begin{enumerate}
    \item Un élève de Première Tech écrit un mail formel à son proviseur pour demander une autorisation d'absence.
    \item Un commentateur sportif radio décrit en direct l'action d'un match Canon vs Tonnerre.
    \item Un chef de quartier préside une réunion de médiation entre deux familles en conflit (palabre).
    \item Vous répondez à une question de grammaire à l'oral devant la classe.
\end{enumerate}
\end{exercice}

\begin{corrige}[Exercice n°1 : Entraînement]
\begin{enumerate}
    \item \textbf{Mail au proviseur :} Émetteur=Élève ; Récepteur=Proviseur (asymétrie forte) ; Code=Français administratif standard ; Canal=Écrit numérique (mail) ; Référent=Demande absence ; Message=Texte structuré (objet, formule politesse, motif, signature). \textbf{Fonction : Conative} (obtenir signature). \textbf{Norme : Formelle, écrite, polie, structurée.}
    \item \textbf{Commentaire radio :} Émetteur=Journaliste ; Récepteur=Auditeurs (masse) ; Code=Français sportif technique + excitatif ; Canal=Ondes radio (oral) ; Référent=Match en cours ; Message=Description instantanée. \textbf{Fonction : Référentielle} (informer) + \textbf{Phatique} (tenir l'antenne). \textbf{Norme : Orale, publique, immédiate, technique.}
    \item \textbf{Palabre médiation :} Émetteur=Chef / Porte-paroles ; Récepteur=Familles + Public ; Code=Langue locale + Proverbes ; Canal=Oral direct ; Référent=Conflit / Paix sociale ; Message=Paroles de sagesse, propositions. \textbf{Fonction : Conative} (apaiser, décider) + \textbf{Poétique} (proverbes). \textbf{Norme : Orale, publique, codifiée (tour de parole), respectueuse.}
    \item \textbf{Réponse orale en classe :} Émetteur=Élève ; Récepteur=Prof + Classe ; Code=Français standard scolaire ; Canal=Oral (air) ; Référent=Notion de grammaire ; Message=Définition / Exemple. \textbf{Fonction : Référentielle} (savoir) + \textbf{Métalinguistique} (parler de la langue). \textbf{Norme : Orale, formelle (levée de main), précise, debout.}
\end{enumerate}
\end{corrige}

\begin{attention}[Dernier piège : L'illusion de l'oralité à l'écrit]
Ne transposez \textbf{jamais} les normes de l'oral (hésitations, répétitions, syntaxe parataxique « et puis... et puis... », tics de langage « en fait », « du coup ») dans votre dissertation (écrit formel). La dissertation exige une \textbf{syntaxe hypotaxique} (subordination, articulation logique par connecteurs : \emph{donc, par conséquent, toutefois, en outre}). Le correcteur attend une \textbf{pensée achevée et architecturée}, non un flux de conscience.
\end{attention}

\coursec{Langue : la ponctuation forte et lecture méthodique du Lion et la Perle}

Dans la section précédente, nous avons vu que toute communication suppose un émetteur, un récepteur, un message, un code, un canal et une situation d'énonciation. Or, à l'écrit, le message ne peut pas s'appuyer sur la voix, les gestes ou les regards : c'est la \cle{ponctuation} qui remplace l'intonation et les silences. Parmi les signes de ponctuation, certains servent à \emph{séparer} à l'intérieur de la phrase (virgule, point-virgule, deux-points) : ce sont les signes de ponctuation faible. D'autres servent à \emph{clore} la phrase et à lui donner son sens définitif : ce sont les signes de \cle{ponctuation forte}, que nous étudions maintenant. Nous verrons ensuite comment ces signes travaillent dans un texte de théâtre, à travers une première lecture méthodique de \emph{Le Lion et la Perle} de Wole Soyinka.

\courssub{La ponctuation forte : définition et inventaire}

\begin{definition}[La ponctuation forte]
La \cle{ponctuation forte} est l'ensemble des signes qui marquent la \textbf{fin d'une unité de sens complète}, c'est-à-dire la clôture d'une phrase. Ces signes sont : le \cle{point} (.), le \cle{point d'interrogation} (?), le \cle{point d'exclamation} (!) et les \cle{points de suspension} (...). Ils indiquent au lecteur qu'une phrase est terminée, quelle est sa nature (énonciative, interrogative, exclamative, inachevée) et qu'une nouvelle phrase, commençant par une majuscule, va suivre.
\end{definition}

L'intuition est simple : quand vous parlez, vous vous arrêtez, vous respirez, vous changez de ton. À l'écrit, ces pauses et ces tons doivent être signalés. Lisez à voix haute la phrase : « Sidi est la perle du village ». Votre voix descend à la fin : c'est le point. Lisez maintenant : « Sidi est-elle la perle du village ? » Votre voix monte : c'est le point d'interrogation. La ponctuation forte est donc la \emph{partition musicale} de la phrase écrite.

\courssub{Le point (.)}

\begin{definition}[Le point]
Le \cle{point} marque la fin d'une \textbf{phrase énonciative} (qui constate, raconte ou décrit) ou d'une \textbf{phrase impérative calme}. Après le point, la phrase suivante commence obligatoirement par une \textbf{majuscule}.
\end{definition}

\begin{exemplebox}[Emplois du point]
\begin{itemize}
\item Phrase énonciative : \emph{Lakunle est l'instituteur du village d'Ilujinle.}
\item Ordre calme : \emph{Écoutez-moi attentivement.}
\item Reprise en majuscule : \emph{Sidi refuse. Lakunle insiste.}
\end{itemize}
\end{exemplebox}

\begin{attention}[Piège fréquent]
Après un point, la majuscule est \textbf{obligatoire}. Écrire « Sidi refuse. lakunle insiste. » est une faute. Inversement, on ne met \textbf{jamais} de point au milieu d'une phrase qui n'est pas finie : la virgule sert à cela.
\end{attention}

\courssub{Le point d'interrogation (?)}

\begin{definition}[Le point d'interrogation]
Le \cle{point d'interrogation} clôt une \textbf{phrase interrogative directe}, c'est-à-dire une question posée directement, avec inversion du sujet, mot interrogatif ou intonation montante.
\end{definition}

\begin{exemplebox}[Interrogation directe et indirecte]
\begin{itemize}
\item Directe : \emph{Accepteras-tu de devenir ma femme ?} (Lakunle à Sidi)
\item Indirecte : \emph{Lakunle demande à Sidi si elle acceptera de devenir sa femme.} (point simple, pas de point d'interrogation)
\end{itemize}
\end{exemplebox}

\begin{attention}[Erreur fréquente : l'interrogation indirecte]
On ne met \textbf{jamais} de point d'interrogation après une \cle{interrogation indirecte}, c'est-à-dire une question rapportée dans une subordonnée introduite par \emph{si}, \emph{ce que}, \emph{où}, \emph{quand}, \emph{pourquoi}. On écrit : « Je me demande s'il viendra\textbf{.} » et non « Je me demande s'il viendra\textbf{?} ». La phrase entière est énonciative : c'est elle qui commande le signe final.
\end{attention}

\courssub{Le point d'exclamation (!)}

\begin{definition}[Le point d'exclamation]
Le \cle{point d'exclamation} clôt une phrase qui exprime une \textbf{émotion vive} (joie, colère, admiration, surprise, douleur), un \textbf{ordre énergique} ou un \textbf{souhait}. Il donne à la phrase une valeur exclamative.
\end{definition}

\begin{exemplebox}[« La Perle ! »]
Dans \emph{Le Lion et la Perle}, le surnom de Sidi, « la Perle », revient souvent. Selon la \textbf{situation d'énonciation}, l'exclamation « La Perle ! » n'a pas le même sens :
\begin{itemize}
\item dans la bouche d'un villageois admiratif, elle exprime l'\textbf{admiration} pour la beauté de Sidi ;
\item dans la bouche de Lakunle, moqueur ou jaloux, elle peut devenir \textbf{ironique} ;
\item dans la bouche de Baroka, elle peut marquer le \textbf{désir} et la convoitise.
\end{itemize}
Le même signe produit donc des \textbf{effets de sens} différents selon qui parle et dans quelle situation.
\end{exemplebox}

\courssub{Les points de suspension (...)}

\begin{definition}[Les points de suspension]
Les \cle{points de suspension} (trois points, jamais deux ni quatre) indiquent que la phrase reste \textbf{inachevée}. Ils signalent une \textbf{énumération non close}, une \textbf{hésitation}, une \textbf{interruption} ou un \textbf{sous-entendu} que le lecteur doit compléter.
\end{definition}

\begin{exemplebox}[Emplois des points de suspension]
\begin{itemize}
\item Énumération inachevée : \emph{Le photographe a montré les champs, le marché, l'école...}
\item Hésitation : \emph{Je... je ne sais pas quoi répondre, Lakunle.}
\item Interruption au théâtre : \emph{Sidi, écoute...} (le personnage est coupé par un autre).
\end{itemize}
\end{exemplebox}

\courssub{Tableau récapitulatif}

\begin{popfigure}
\begin{center}
\begin{tikzpicture}
\fill[popblueL] (0,0) rectangle (14.5,1);
\node[anchor=west,text=popdark,font=\bfseries] at (0.2,0.5) {Signe};
\node[anchor=west,text=popdark,font=\bfseries] at (2.2,0.5) {Fonction principale};
\node[anchor=west,text=popdark,font=\bfseries] at (8.2,0.5) {Exemple inspiré du \emph{Lion et la Perle}};
\draw[popline] (0,0) -- (14.5,0);
\fill[popcream] (0,-1) rectangle (14.5,0);
\node[anchor=west] at (0.2,-0.5) {\textbf{.}};
\node[anchor=west,text width=5.6cm] at (2.2,-0.5) {Fin d'une phrase énonciative ; reprise en majuscule};
\node[anchor=west,text width=6cm] at (8.2,-0.5) {\emph{Sidi est la plus belle fille d'Ilujinle.}};
\fill[white] (0,-2) rectangle (14.5,-1);
\node[anchor=west] at (0.2,-1.5) {\textbf{?}};
\node[anchor=west,text width=5.6cm] at (2.2,-1.5) {Question directe (jamais dans l'interrogation indirecte)};
\node[anchor=west,text width=6cm] at (8.2,-1.5) {\emph{Veux-tu m'épouser sans dot ?}};
\fill[popcream] (0,-3) rectangle (14.5,-2);
\node[anchor=west] at (0.2,-2.5) {\textbf{!}};
\node[anchor=west,text width=5.6cm] at (2.2,-2.5) {Émotion, ordre, admiration ou ironie};
\node[anchor=west,text width=6cm] at (8.2,-2.5) {\emph{La Perle ! Quelle beauté !}};
\fill[white] (0,-4) rectangle (14.5,-3);
\node[anchor=west] at (0.2,-3.5) {\textbf{...}};
\node[anchor=west,text width=5.6cm] at (2.2,-3.5) {Phrase inachevée : hésitation, interruption, énumération ouverte};
\node[anchor=west,text width=6cm] at (8.2,-3.5) {\emph{Moi, épouser un vieillard... jamais !}};
\draw[popline] (0,-4) -- (14.5,-4);
\draw[popline] (0,1) -- (0,-4);
\draw[popline] (2,1) -- (2,-4);
\draw[popline] (8,1) -- (8,-4);
\draw[popline] (14.5,1) -- (14.5,-4);
\end{tikzpicture}
\end{center}
\legende{Tableau récapitulatif des quatre signes de ponctuation forte : signe, fonction et exemple tiré de l'univers du \emph{Lion et la Perle} de Wole Soyinka.}
\end{popfigure}

\begin{aretenir}[L'essentiel]
Quatre signes clôturent la phrase : le \textbf{point} (constat), le \textbf{point d'interrogation} (question directe uniquement), le \textbf{point d'exclamation} (émotion ou ordre), les \textbf{points de suspension} (inachèvement). Après chacun d'eux, la phrase suivante commence par une majuscule.
\end{aretenir}

\courssub{Lecture méthodique 1 du Lion et la Perle : la rencontre Sidi-Lakunle}

\textbf{Rappel de l'œuvre.} \emph{Le Lion et la Perle} (1959) est une pièce du dramaturge nigérian Wole Soyinka, prix Nobel de littérature en 1986. Dans le village d'Ilujinle, le vieux chef Baroka, « le Lion », et le jeune instituteur moderniste Lakunle se disputent Sidi, la plus belle fille du village, « la Perle ». L'extrait étudié ici est la première rencontre entre Sidi et Lakunle au début de la pièce.

\begin{methode}[Analyser un extrait de théâtre]
\begin{enumerate}
\item \textbf{Identifier les personnages} présents : qui parle ? qui est sur scène ?
\item \textbf{Situer la scène} : où et quand se déroule-t-elle ? Que s'est-il passé juste avant ?
\item \textbf{Dégager les enjeux} : que veut chaque personnage ? Quel conflit les oppose ?
\item \textbf{Observer les didascalies} (indications scéniques en italique) : gestes, tons, déplacements.
\item \textbf{Repérer la ponctuation forte} : relever les points d'exclamation, d'interrogation, de suspension, et expliquer leur \textbf{effet de sens} (colère, moquerie, hésitation, coupure de parole...).
\item \textbf{Conclure} : quelle vision des personnages et quel thème (tradition contre modernité) cette scène fait-elle naître ?
\end{enumerate}
\end{methode}

\textbf{Application à l'extrait.} Dans la scène d'ouverture, Sidi porte une cruche d'eau sur la tête ; Lakunle l'arrête et la sermonne. Le dialogue est traversé par la ponctuation forte :

\begin{itemize}
\item Les \textbf{points d'exclamation} de Lakunle (« Sidi ! ») traduisent son exaspération de maître d'école qui veut corriger la « sauvagerie » du village : ils construisent un personnage \textbf{exalté et autoritaire}.
\item Les \textbf{points d'interrogation} de Sidi montrent sa vivacité : elle retourne les arguments de Lakunle contre lui et le ridiculise. L'interrogation devient une \textbf{arme de réplique}.
\item Les \textbf{points de suspension} marquent les interruptions : les deux personnages se coupent la parole, signe d'un dialogue \textbf{vif et conflictuel}, presque un duel verbal.
\item Les \textbf{didascalies} (gestes, tons) complètent la ponctuation : elles indiquent comment la réplique doit être dite, ce que le texte seul ne peut montrer.
\end{itemize}

\begin{exoresolu}[Repérage de la ponctuation dans un dialogue]
Énoncé : dans le passage suivant, inspiré de la rencontre Sidi-Lakunle, identifie chaque signe de ponctuation forte et explique son effet de sens.

\emph{LAKUNLE : Sidi ! Reviens ! Une femme ne doit pas porter de charges sur la tête...\\
SIDI : Et pourquoi pas, l'instituteur ?\\
LAKUNLE : Parce que c'est une coutume arriérée ! Je me demande si tu comprendras un jour.}

\tcblower
Solution détaillée :
\begin{enumerate}
\item « Sidi ! » : point d'exclamation ; il exprime l'appel énergique et l'agacement de Lakunle.
\item « Reviens ! » : point d'exclamation après un impératif ; il marque l'\textbf{ordre}, le ton de maître que se donne Lakunle.
\item « ...sur la tête... » : points de suspension ; ils signalent une phrase qui se prolonge, une énumération de reproches ou une hésitation rhétorique.
\item « Et pourquoi pas, l'instituteur ? » : point d'interrogation ; question \textbf{directe} (intonation montante, apostrophe) ; Sidi conteste, l'interrogation est ici une réplique moqueuse.
\item « ...une coutume arriérée ! » : point d'exclamation ; jugement péremptoire, mépris de Lakunle pour la tradition.
\item « Je me demande si tu comprendras un jour. » : \textbf{point simple} et non point d'interrogation, car il s'agit d'une \textbf{interrogation indirecte} (introduite par \emph{si}) incluse dans une phrase énonciative.
\end{enumerate}
Conclusion : la ponctuation forte révèle le conflit — un Lakunle donneur de leçons (exclamations, ordres) face à une Sidi vive et ironique (interrogations).
\end{exoresolu}

\courssub{Contrôle de lecture, compte rendu et négociation du projet d'étude}

Après cette première lecture méthodique, la classe procède à trois activités :

\begin{itemize}
\item \textbf{Contrôle de lecture} : questions rapides sur l'œuvre (personnages, lieu, début de l'intrigue) pour vérifier que chacun a lu le passage. Exemples : Qui est Baroka ? Pourquoi Sidi est-elle surnommée « la Perle » ? Que reproche Lakunle aux coutumes du village ?
\item \textbf{Compte rendu} : chaque élève (ou groupe) formule à l'oral ou à l'écrit un résumé objectif de la scène, sans jugement personnel, en employant correctement la ponctuation forte dans ses propres phrases.
\item \textbf{Négociation du projet d'étude} : la classe et le professeur décident ensemble de la suite du travail sur l'œuvre (lecture par scènes, jeux de rôle, exposés sur le conflit tradition/modernité, préparation d'une dissertation). Négocier, c'est discuter pour aboutir à un accord : chacun propose, justifie, puis la classe tranche.
\end{itemize}

\courssub{Remédiation sur les acquis}

\begin{exercice}[Ponctuation forte]
Ponctue correctement les phrases suivantes (signe fort et majuscule) :
\begin{enumerate}
\item lakunle demande à sidi si elle veut l'épouser
\item sidi s'écrie jamais je n'épouserai un homme sans dot
\item le lion ruse il attend son heure
\item veux-tu vraiment moderniser ce village
\end{enumerate}
\end{exercice}

\begin{corrige}[Exercice : ponctuation forte]
\begin{enumerate}
\item \emph{Lakunle demande à Sidi si elle veut l'épouser.} — Point simple : interrogation \textbf{indirecte} (introduite par \emph{si}).
\item \emph{Sidi s'écrie : « Jamais ! Je n'épouserai un homme sans dot ! »} — Points d'exclamation : refus énergique, émotion.
\item \emph{Le Lion ruse... Il attend son heure.} — Points de suspension (sous-entendu, attente) puis point ; majuscule après chaque signe fort.
\item \emph{Veux-tu vraiment moderniser ce village ?} — Point d'interrogation : question \textbf{directe} avec inversion.
\end{enumerate}
\end{corrige}

\begin{aretenir}[Bilan de la section]
La ponctuation forte (., ?, !, ...) clôt la phrase et en fixe le sens : constat, question directe, émotion ou inachèvement. Au théâtre, elle est un indice précieux du ton et des rapports de force entre personnages. Dans \emph{Le Lion et la Perle}, le duel verbal entre Sidi et Lakunle se lit autant dans les signes de ponctuation que dans les mots : savoir les repérer et les interpréter est une compétence essentielle pour le commentaire et la dissertation au Probatoire et au Baccalauréat technique.
\end{aretenir}

\newpage

\coursec{Activité d'intégration}

\coursec{Leçon 17 : Produire un plan détaillé d'une dissertation}

\courssub{Objectifs de l'activité}

À l'issue de cette activité d'intégration, l'élève doit être capable de :

\begin{itemize}
\item analyser un sujet de dissertation (dégager le thème, définir les mots-clés, repérer la question posée) ;
\item formuler une \cle{problématique} claire, pertinente et directrice ;
\item conduire une \cle{recherche d'idées} organisée (brainstorming, carte mentale, tri et hiérarchisation) ;
\item choisir un \cle{type de plan} adapté à la nature du sujet et justifier ce choix ;
\item rédiger un \cle{plan détaillé} complet et équilibré (introduction, parties avec arguments et exemples précis, conclusion) ;
\item employer correctement la \cle{ponctuation forte} (point, deux-points, point-virgule, point d'interrogation, points de suspension) ;
\item évaluer la production d'autrui à l'aide d'une grille critériée et justifier son jugement.
\end{itemize}

\courssub{Situation}

Le journal du lycée de Nkongsamba, \emph{La Plume du Mungo}, organise son concours annuel de dissertation à l'occasion de la semaine de la culture. Le sujet retenu cette année par le jury est le suivant :

\begin{center}
\textbf{« Le développement de nos localités exige-t-il l'abandon de nos traditions ? »}
\end{center}

Le règlement du concours précise que chaque candidat doit remettre, avant la rédaction finale, un \textbf{plan détaillé} de sa dissertation. Ce plan doit comporter : une introduction (amorce, présentation du sujet, problématique, annonce du plan), deux ou trois parties équilibrées (chacune avec au moins deux arguments illustrés par un exemple précis), et une conclusion (bilan, réponse à la problématique, ouverture). Le jury accorde une attention particulière à la pertinence de la problématique, à l'équilibre des parties, à la qualité des exemples et à la correction de la langue, notamment l'emploi de la ponctuation forte.

Dans ta classe de Première technique, le professeur de français décide de préparer la classe à ce concours. Vous travaillez par groupes de quatre élèves. Vous disposez de votre étude de l'œuvre \emph{Le Lion et la Perle} de Wole Soyinka, qui met en scène le conflit entre Baroka, le chef traditionnel du village d'Ilujinle, et Lakunle, le jeune instituteur partisan de la modernité. Vous pouvez aussi vous appuyer sur la réalité camerounaise : chefferies traditionnelles de l'Ouest, modernisation des villages, électrification rurale, téléphonie mobile et TIC, fêtes traditionnelles comme le Ngondo ou le Nguon.

À toi de jouer : avec ton groupe, produis le meilleur plan détaillé possible pour remporter le concours.

\courssub{Consignes}

\begin{enumerate}
\item \textbf{Analyse du sujet.} Souligne les mots-clés du sujet (« développement », « localités », « exige », « abandon », « traditions »). Définis chacun d'eux en une phrase. Quel type de sujet est-ce : question ouverte ou fermée ? Quelle réponse le sujet semble-t-il attendre ?
\item \textbf{Problématique.} Formule la problématique sous la forme d'une ou deux questions qui traduisent le véritable enjeu du sujet : le progrès et la tradition sont-ils forcément ennemis ?
\item \textbf{Recherche des idées.} Pendant dix minutes, conduis un brainstorming en groupe : note toutes les idées qui viennent, sans les juger. Appuie-toi sur le conflit Baroka/Lakunle dans \emph{Le Lion et la Perle} et sur des exemples camerounais concrets (chefferies, modernisation des villages, TIC, fêtes traditionnelles).
\item \textbf{Tri et hiérarchisation.} Reporte tes idées sur une carte mentale : le thème au centre, les grandes orientations en branches principales, les exemples en branches secondaires. Élimine les idées hors sujet et regroupe les idées proches.
\item \textbf{Choix du plan.} Parmi les types de plans étudiés (dialectique, thématique, analytique), choisis celui qui convient le mieux à ce sujet. Justifie ton choix par écrit en trois ou quatre lignes.
\item \textbf{Rédaction du plan détaillé.} Rédige le plan complet : introduction (amorce, sujet, problématique, annonce du plan), deux ou trois parties avec arguments et exemples, conclusion. Respecte scrupuleusement la ponctuation forte.
\item \textbf{Présentation et évaluation.} Présente le plan de ton groupe au tableau en cinq minutes. Écoute les autres groupes et évalue leur production à l'aide de la grille critériée fournie par le professeur.
\end{enumerate}

\courssub{Ressources à mobiliser}

\begin{aretenir}[Les étapes de la dissertation]
\begin{itemize}
\item \textbf{1. Analyser le sujet} : définir les mots-clés, délimiter le thème, repérer la question et le type de sujet.
\item \textbf{2. Formuler la problématique} : transformer le sujet en question directrice qui ouvre la discussion.
\item \textbf{3. Rechercher les idées} : brainstorming (quantité), puis tri, regroupement et hiérarchisation (qualité).
\item \textbf{4. Choisir un plan} : dialectique (thèse / antithèse / synthèse), thématique (aspects d'une même question), analytique (constat / causes / conséquences ou causes / effets / solutions).
\item \textbf{5. Rédiger le plan détaillé} : introduction codifiée, parties équilibrées (argument + exemple + explication), conclusion structurée.
\end{itemize}
\end{aretenir}

\begin{aretenir}[La ponctuation forte -- Rappel]
\begin{itemize}
\item Le \textbf{point} (.) marque la fin d'une phrase énonciative affirmative ou négative.
\item Les \textbf{deux-points} (:) annoncent une explication, une énumération, une citation ou une conséquence.
\item Le \textbf{point-virgule} (;) sépare des propositions longues ou déjà ponctuées de virgules, liées par le sens (opposition, cause, conséquence, énumération).
\item Le \textbf{point d'interrogation} (?) termine une question directe ; c'est le signe obligatoire de la problématique.
\item Les \textbf{points de suspension} (...) marquent une hésitation, une omission ou une ouverture (utile en conclusion).
\end{itemize}
\end{aretenir}

\begin{methode}[Analyser un sujet de dissertation]
\begin{enumerate}
\item Recopier le sujet et souligner les \textbf{mots-clés} (noms, verbes, adjectifs qualificatifs).
\item \textbf{Définir} chaque mot-clé dans le contexte du sujet (dictionnaire + contexte).
\item Identifier le \textbf{thème} (sujet général) et la \textbf{question} posée (verbe + sujet).
\item Classer le sujet : \textbf{question fermée} (Oui/Non, Pour/Contre $\rightarrow$ plan dialectique) ou \textbf{question ouverte} (Comment/Pourquoi/Quels $\rightarrow$ plan thématique ou analytique).
\item Reformuler le sujet avec ses propres mots pour vérifier la compréhension.
\end{enumerate}
\end{methode}

\begin{methode}[Formuler une problématique]
\begin{enumerate}
\item Partir de l'opposition ou de la tension repérée dans le sujet (ex : développement vs traditions).
\item Rédiger une \textbf{question directe} (terminée par ?) qui ne soit pas une simple reformulation du sujet.
\item Vérifier que la problématique \textbf{annonce la structure du plan} (les grandes parties y répondront).
\item Éviter les questions trop générales (hors sujet) ou trop fermées (réponse Oui/Non immédiate).
\end{enumerate}
\end{methode}

\begin{methode}[Rechercher et hiérarchiser les idées]
\begin{enumerate}
\item \textbf{Brainstorming} (5-10 min) : noter \emph{toutes} les idées sans tri (mots, phrases, références littéraires, exemples vécus).
\item \textbf{Tri} : éliminer le hors-sujet, les doublons, les idées trop vagues.
\item \textbf{Regroupement} : rassembler les idées par affinités (causes, conséquences, arguments pour, arguments contre, exemples littéraires, exemples réels).
\item \textbf{Hiérarchisation} : identifier 2 ou 3 \textbf{grands axes} (futures parties) et ordonner les arguments à l'intérieur (du plus fort au plus faible, ou chronologique, ou logique).
\item \textbf{Cartographie} : tracer la carte mentale (thème au centre $\rightarrow$ axes $\rightarrow$ arguments $\rightarrow$ exemples).
\end{enumerate}
\end{methode}

\begin{methode}[Choisir et élaborer un plan détaillé]
\begin{enumerate}
\item \textbf{Choisir le type} selon le sujet : Dialectique (opposition), Thématique (facettes), Analytique (processus).
\item \textbf{Titrer les parties} : noms courts, parallèles, sans verbes conjugués (ex: \emph{Les freins traditionnels au développement} / \emph{La tradition, levier du développement}).
\item \textbf{Structurer chaque partie} : Idée directrice (phrase nominale) $\rightarrow$ Argument 1 + Exemple 1 $\rightarrow$ Argument 2 + Exemple 2.
\item \textbf{Rédiger l'introduction} : Amorce $\rightarrow$ Sujet posé $\rightarrow$ Problématique $\rightarrow$ Annonce du plan.
\item \textbf{Rédiger la conclusion} : Bilan synthétique $\rightarrow$ Réponse nette à la problématique $\rightarrow$ Ouverture pertinente.
\end{enumerate}
\end{methode}

\begin{attention}[Pièges fréquents à éviter]
\begin{itemize}
\item \textbf{Hors-sujet :} Ne pas définir les mots-clés ou traiter « développement » sans « traditions » (ou l'inverse).
\item \textbf{Problématique faible :} Simple reformulation du sujet (ex: « Faut-il abandonner nos traditions pour développer nos localités ? ») ou question trop large (ex: « Qu'est-ce que le développement ? »).
\item \textbf{Plan déséquilibré :} Une partie avec 3 arguments, l'autre avec 1 seul ; ou des exemples tous littéraires (ou tous réels) sans variété.
\item \textbf{Exemples « flottants » :} Citer \emph{Le Lion et la Perle} sans lien explicite avec l'argument (ex: « Baroka est polygame » sans expliquer en quoi cela freine ou aide le développement).
\item \textbf{Ponctuation défaillante :} Oublier le ? à la problématique ; confondre ; et , ; enchaîner des phrases sans point.
\end{itemize}
\end{attention}

\begin{savaistu}[Contexte culturel et littéraire : Le Cameroun et Wole Soyinka]
\begin{itemize}
\item \textbf{Wole Soyinka} (Nigéria, Prix Nobel 1986) : \emph{Le Lion et la Perle} (1963) met en scène le choc des cultures. \textbf{Baroka} (le Lion), chef rusé, incarne une tradition vivante, capable d'absorber la modernité (le timbre-poste, la machine à coudre) sans perdre son autorité. \textbf{Lakunle}, l'instituteur, incarne une modernité mimétique, rejetant la dot et les coutumes, mais déconnecté des réalités villageoises. \textbf{Sidi} (la Perle) choisit la tradition vivante (Baroka) plutôt que la modernité stérile (Lakunle).
\item \textbf{Réalité camerounaise} : Les chefferies traditionnelles (Bamiléké, Bamoun, Bassa, Douala, etc.) sont des acteurs majeurs du développement local (construction d'écoles, de centres de santé, gestion des conflits fonciers, organisation des fêtes culturelles). Le \textbf{Ngondo} (peuple Sawa) et le \textbf{Nguon} (Bamoun) sont des moteurs économiques (tourisme, artisanat) et identitaires. L'électrification rurale et la téléphonie mobile transforment les villages sans nécessairement effacer les coutumes (ex: tontines digitalisées, chefferies sur Facebook).
\end{itemize}
\end{savaistu}

\begin{exemplebox}[Grille critériée d'évaluation du plan détaillé]
\begin{center}
\begin{tabularx}{\linewidth}{|X|c|c|}
\hline
\rowcolor{popblueL} \textbf{Critère} & \textbf{Indicateurs de réussite} & \textbf{Note /5} \\
\hline
Analyse du sujet & Mots-clés définis, thème délimité, type de sujet identifié & \\
\hline
Pertinence de la problématique & Question directe, enjeu clair, annonce du plan & \\
\hline
Choix et justification du plan & Type adapté, justification écrite cohérente & \\
\hline
Équilibre et cohérence des parties & 2 ou 3 parties, titres parallèles, progression logique & \\
\hline
Qualité des arguments et exemples & Arguments distincts, exemples précis (littéraire + réel), lien explicite & \\
\hline
Structure de l'introduction & Amorce, sujet posé, problématique, annonce du plan & \\
\hline
Structure de la conclusion & Bilan, réponse à la problématique, ouverture & \\
\hline
Correction de la langue & Ponctuation forte maîtrisée, syntaxe correcte, vocabulaire précis & \\
\hline
Présentation orale & Clarté, respect du temps, posture, réponses aux questions & \\
\hline
\end{tabularx}
\end{center}
\end{exemplebox}

\courssub{Démarche et corrigé commenté}

\begin{exoresolu}[Du sujet au plan détaillé -- Sujet : « Le développement de nos localités exige-t-il l'abandon de nos traditions ? »]
Produire le plan détaillé d'une dissertation.
\tcblower
\textbf{Étape 1 : Analyse du sujet.}
\begin{itemize}
\item \textbf{Développement :} processus d'amélioration durable des conditions de vie (économiques, sociales, techniques, culturelles) d'une population.
\item \textbf{Localités :} villages, quartiers, arrondissements, entités territoriales de base.
\item \textbf{Exige :} verbe marquant la contrainte, la nécessité absolue, l'inconditionnalité.
\item \textbf{Abandon :} renoncement définitif, rupture totale, disparition.
\item \textbf{Traditions :} ensemble des coutumes, valeurs, croyances, institutions (chefferies, rites, langues, fêtes) transmises par les ancêtres.
\item \textbf{Type de sujet :} Question fermée (appelle Oui ou Non). Le verbe « exige » pose une condition absolue qui appelle la discussion (nuance).
\end{itemize}

\textbf{Étape 2 : Problématique.}
Le développement de nos localités est-il nécessairement incompatible avec la survie de nos traditions ? Faut-il opposer modernité et coutume, ou peut-on concevoir un progrès qui s'enracine dans l'héritage ancestral ?

\textbf{Étape 3 : Recherche des idées (Brainstorming -- 10 min).}
\textit{Idées littéraires (Le Lion et la Perle) :} Lakunle refuse la dot, veut transformer Ilujinle en ville moderne (routes, écoles) ; Baroka détourne le progrès (timbre, machine) pour asseoir son pouvoir ; Sidi choisit Baroka ; Sadiku et les femmes : gardiennes de la tradition ; le bale (fête) comme moment de cohésion.
\textit{Idées contextuelles (Cameroun) :} Chefferies finançant des infrastructures ; Ngondo/Nguon : attractivité touristique et cohésion ; Tontines traditionnelles $\rightarrow$ micro-crédits modernes ; Téléphonie mobile facilitant l'agriculture et la santé au village ; Exode rural des jeunes ; Perte des langues maternelles ; Élites construisant au village ; Conflits fonciers tradition vs droit moderne.

\textbf{Étape 4 : Tri et hiérarchisation (Carte mentale).}
\begin{center}
\textbf{THÈME : Développement vs Traditions}
\end{center}
\begin{itemize}
\item \textbf{Branche A : Tensions / Freins} (Le développement \emph{semble} exiger l'abandon)
  \begin{itemize}
  \item Argument : Certaines coutumes freinent l'innovation (ex: dot, refus de l'école pour les filles, gestion foncière coutumière bloquant les investissements).
  \item Argument : La modernisation impose de nouveaux modes de vie (ex: exode rural, individualisme, perte des langues).
  \end{itemize}
\item \textbf{Branche B : Complémentarité / Appui} (Les traditions \emph{servent} le développement)
  \begin{itemize}
  \item Argument : Valeurs de solidarité et d'autorité légitime (ex: chefferies moteurs de développement local, tontines).
  \item Argument : Patrimoine culturel comme ressource économique (ex: Ngondo, Nguon, artisanat, tourisme culturel).
  \end{itemize}
\item \textbf{Branche C : Dépassement / Synthèse} (Concilier les deux par une adaptation intelligente)
  \begin{itemize}
  \item Argument : La tradition n'est pas figée, elle évolue (ex: Baroka utilise la modernité ; tontines digitalisées).
  \item Argument : Un développement endogène durable s'appuie sur l'identité (ex: élites investissant au village, écoles bilingues français/langue locale).
  \end{itemize}
\end{itemize}
\textit{Idées écartées :} Description pure d'Ilujinle, biographie de Soyinka, liste des fêtes sans lien avec le développement.

\textbf{Étape 5 : Choix du plan.}
\textbf{Plan dialectique (Thèse / Antithèse / Synthèse).}
\textit{Justification :} Le sujet est une question fermée opposant deux termes forts (« développement » vs « traditions », « exige » vs « abandon »). La structure thèse/antithèse permet d'examiner tour à tour l'opposition apparente et la réalité du terrain, avant de les dépasser dans une synthèse qui répond « Non, mais... » à la question « Exige-t-il ? ».

\textbf{Étape 6 : Rédaction du plan détaillé.}

\medskip
\noindent\textbf{INTRODUCTION}
\begin{itemize}
\item \textbf{Amorce :} De nos villages de l'Ouest aux rives du Wouri, le visage de nos localités change à une vitesse inédite : bitume, pylônes, smartphones côtoient cases à toit de paille et cases de chefferie.
\item \textbf{Présentation du sujet :} Face à cette mutation, une interrogation s'impose : le progrès impose-t-il la table rase du passé ? Nombreux sont ceux qui pensent que moderniser, c'est « occidentaliser », et donc abandonner nos coutumes.
\item \textbf{Problématique :} Le développement de nos localités est-il incompatible avec la pérennité de nos traditions ?
\item \textbf{Annonce du plan :} Nous analyserons d'abord les tensions réelles entre modernisation et coutumes (I), avant de montrer que les traditions constituent un levier puissant pour le développement (II), pour enfin démontrer qu'un développement authentique suppose une réinvention créative de l'héritage (III).
\end{itemize}

\medskip
\noindent\textbf{PREMIÈRE PARTIE : L'ILLUSION D'UNE RUPTURE NÉCESSAIRE (THÈSE)}
\begin{itemize}
\item \textbf{Idée directrice :} Certaines pratiques traditionnelles apparaissent comme des freins objectifs à la modernisation technique et sociale.
\item \textbf{Argument 1 :} Le poids de certaines coutumes peut entraver l'innovation et l'égalité.
\item \textbf{Exemple 1 :} Dans \emph{Le Lion et la Perle}, Lakunle dénonce la dot (« prix de la mariée ») comme une marchandisation de la femme et refuse de la payer, voyant en elle un symbole d'archaïsme bloquant l'émancipation féminine et l'entrée dans la modernité juridique.
\item \textbf{Argument 2 :} L'arrivée des infrastructures modernes bouleverse les structures sociales traditionnelles.
\item \textbf{Exemple 2 :} Dans de nombreux villages de l'Ouest camerounais, l'ouverture de routes goudronnées et l'installation de la téléphonie mobile ont accéléré l'exode des jeunes vers les villes, fragilisant la transmission orale et l'autorité des notables.
\end{itemize}

\medskip
\noindent\textbf{DEUXIÈME PARTIE : LA TRADITION, FONDEMENT ET MOTEUR DU DÉVELOPPEMENT LOCAL (ANTITHÈSE)}
\begin{itemize}
\item \textbf{Idée directrice :} Loin d'être un obstacle, l'héritage ancestral fournit le capital social, identitaire et économique indispensable à tout développement durable.
\item \textbf{Argument 1 :} Les institutions traditionnelles (chefferies, sociétés secrètes, tontines) sont des acteurs majeurs de l'aménagement du territoire.
\item \textbf{Exemple 1 :} Au Cameroun, les chefferies de 1er et 2nd degré (Bamiléké, Bamoun, Sawa, etc.) pilotent la construction d'écoles, de centres de santé, de marchés et gèrent les conflits fonciers, assurant une gouvernance de proximité que l'État ne peut seul assurer.
\item \textbf{Argument 2 :} Le patrimoine culturel immatériel et matériel est une ressource économique directe (tourisme, artisanat, attractivité).
\item \textbf{Exemple 2 :} Les fêtes du \textbf{Ngondo} (Douala) et du \textbf{Nguon} (Foumban) drainent des milliers de visiteurs, génèrent des revenus pour les artisans, hôtels, transporteurs, et renforcent la fierté identitaire, condition du « bien-vivre ensemble ».
\end{itemize}

\medskip
\noindent\textbf{TROISIÈME PARTIE : VERS UN DÉVELOPPEMENT ENDOGÈNE : ADAPTER LA TRADITION AU LIEU DE L'ABANDONNER (SYNTHÈSE)}
\begin{itemize}
\item \textbf{Idée directrice :} Le vrai développement ne détruit pas l'identité ; il la met en mouvement. La tradition n'est pas un musée, c'est une ressource vivante.
\item \textbf{Argument 1 :} Les figures d'autorité traditionnelles savent intégrer la modernité sans perdre leur légitimité (hybridation réussie).
\item \textbf{Exemple 1 :} \textbf{Baroka} lui-même, dans la pièce de Soyinka, n'est pas un réactionnaire : il utilise la ruse et la modernité (le timbre-poste, la promesse de machine) pour conforter son pouvoir et développer son village, prouvant que tradition et modernité peuvent être complices.
\item \textbf{Argument 2 :} Les élites locales et la diaspora construisent un développement « par le bas » ancré dans les valeurs communautaires.
\item \textbf{Exemple 2 :} De nombreux cadres camerounais (originaires de l'Ouest, du Nord, du Sud) financent l'électrification solaire, les forages, les écoles bilingues (français/ghomala', bassa, fulfulde) dans leurs villages d'origine, gérés par les comités de développement villageois issus de la tradition.
\end{itemize}

\medskip
\noindent\textbf{CONCLUSION}
\begin{itemize}
\item \textbf{Bilan :} Au terme de cette analyse, il apparaît que le verbe « exiger » du sujet est trop fort. Le développement ne \emph{contraint} pas à l'abandon ; il \emph{provoque} une transformation.
\item \textbf{Réponse à la problématique :} Non, le développement de nos localités n'exige pas l'abandon de nos traditions. Il exige au contraire leur \textbf{réinvention} : les traditions doivent cesser d'être des chaînes pour devenir des racines.
\item \textbf{Ouverture :} Le défi majeur pour la jeunesse technique camerounaise est d'inventer les outils (numérique, agro-industrie, énergies renouvelables) qui permettront à nos chefferies, nos langues et nos rites de traverser le XXIe siècle non pas en survivants, mais en acteurs. Comme le suggère Soyinka, le Lion doit apprendre à rugir dans la langue de l'ordinateur.
\end{itemize}

\textbf{Étape 7 : Vérification finale (Auto-évaluation).}
\begin{itemize}
\item \checkmark Problématique terminée par un point d'interrogation (?).
\item \checkmark Deux-points (:) utilisés pour introduire les exemples et les explications.
\item \checkmark Point-virgule (;) utilisé pour lier les arguments au sein des parties (ex: « frein objectif ; le poids... »).
\item \checkmark Parties équilibrées (2 arguments + 2 exemples chacune).
\item \checkmark Exemples variés : 2 littéraires (Soyinka), 4 contextuels Cameroun (chefferies, Ngondo/Nguon, Baroka/ruse, élites/diaspora).
\item \checkmark Annonce du plan respectée dans le développement.
\item \checkmark Ouverture en lien direct avec la formation technique (jeunesse, outils, innovation).
\end{itemize}
\end{exoresolu}

\courssub{Bilan de l'activité}

Cette activité a permis de mobiliser, dans une tâche complète et authentique (concours du journal scolaire), l'ensemble des savoir-faire de la leçon. Elle a d'abord montré que l'\textbf{analyse rigoureuse du sujet} conditionne toute la dissertation : définir les mots-clés (« exige », « abandon ») et repérer le type de sujet (question fermée) évite le hors-sujet et oriente le choix du plan. Elle a ensuite mis en évidence que la \textbf{problématique} n'est pas une simple reformulation, mais la question directrice qui annonce la discussion et structure la pensée. Le \textbf{brainstorming} suivi de la \textbf{carte mentale} a prouvé son utilité opérationnelle : c'est le tri et le regroupement des idées qui font naître naturellement les grandes parties. Le choix du \textbf{plan dialectique} a été justifié par la nature même du sujet, opposition binaire qu'il faut discuter avant de la dépasser. Enfin, la rédaction du plan détaillé a permis de vérifier que chaque argument doit être soutenu par un exemple précis — tiré de \emph{Le Lion et la Perle} ou de la réalité camerounaise vécue — et que la \textbf{ponctuation forte} (?, :, ;) organise la clarté du propos et signale l'architecture de la pensée. L'évaluation croisée par grille critériée a en outre développé chez les élèves l'esprit critique et la capacité à justifier un jugement, compétences transversales essentielles pour l'épreuve de dissertation au Probatoire et au Baccalauréat technique.

\newpage

\coursec{Méthodes et exercices résolus}

\coursec{Produire un plan détaillé d'une dissertation : de l'analyse du sujet à la rédaction du plan}

\courssub{Réception du texte : Paratextes et lecture ouvroir de \emph{Le Lion et la Perle}}

\begin{methode}[Analyser les paratextes d'une œuvre théâtrale]
\begin{enumerate}
\item \textbf{Identifier les éléments paratextuels} : titre, sous-titre, nom de l'auteur, préface/avant-propos, dédicace, page de garde, distribution des personnages, didascalies initiales, notes de l'éditeur.
\item \textbf{Interroger le titre} : \emph{Le Lion et la Perle} $\rightarrow$ opposition métaphorique (force/puissance vs beauté/fragilité, tradition vs modernité), annonce du conflit central.
\item \textbf{Situer l'auteur et le contexte} : Wole \cle{Soyinka}, prix Nobel 1986, Nigéria, écriture en anglais (traduction française), théâtre postcolonial, confrontation cultures yoruba/occidentale.
\item \textbf{Exploiter la distribution et les didascalies} : repérer les archétypes (Baroka le Bale, Lakunle l'instituteur, Sidi la belle), les lieux (Ilujinle, la brousse, la salle de classe), le temps (une journée).
\item \textbf{Formuler des hypothèses de lecture} : à partir de ces indices, prédire les thèmes (mariage, dot, pouvoir, séduction, identité) et le type de comédie (satire, comédie de mœurs).
\end{enumerate}
\begin{aretenir}[Réflexe paratexte]
Les paratextes sont le \emph{seuil} de l'œuvre (Genette). Ils orientent la lecture avant même le premier vers. Ne les sautez jamais : ils donnent les clés du contrat de lecture.
\end{aretenir}
\end{methode}

\begin{exoresolu}[Lecture ouvroir : les paratextes de \emph{Le Lion et la Perle}]
\textbf{Énoncé.} À partir de la page de garde, de la distribution et des didascalies initiales, présentez les hypothèses de lecture pour l'étude de la pièce.
\tcblower
\textbf{Solution détaillée.}
\begin{enumerate}
\item \textbf{Titre} : \og Le Lion et la Perle \fg{} $\rightarrow$ allégorie. Le Lion = Baroka (chef traditionnel, virilité, ruse, pouvoir). La Perle = Sidi (beauté, jeunesse, objet de convoitise, mais aussi valeur). L'article défini \og Le \fg{} / \og La \fg{} singularise : un duel précis, pas une généralité.
\item \textbf{Auteur} : Wole Soyinka (né 1934), figure majeure des lettres africaines anglophones. Théâtre engagé, mélange rituel yoruba / forme occidentale.
\item \textbf{Distribution (personnages principaux)} :
\begin{itemize}
\item \textbf{Baroka} (le Bale) : 62 ans, rusé, polygame, incarne la tradition vivante, adaptative.
\item \textbf{Lakunle} : instituteur, ~23 ans, \cle{moderniste} dogmatique, refuse la dot, parle un anglais ampoulé, caricature de l'occidentalisé maladroit.
\item \textbf{Sidi} : la \og Perle \fg{}, villageoise belle, fière, consciente de son pouvoir de séduction (photos du magazine), pivot du conflit.
\item \textbf{Sadiku} : première épouse de Baroka, médiatrice, loyale mais manipulée.
\end{itemize}
\item \textbf{Didascalies initiales} : \og Le marché, le matin \fg{} (espace public, vie quotidienne) ; \og Près de l'école \fg{} (espace de la modernité imposée). Unité de lieu relative (village), unité de temps (une journée).
\item \textbf{Hypothèses de lecture} :
\begin{itemize}
\item Conflit générationnel et culturel : tradition vs modernité (mais nuancé : Baroka utilise la modernité — machine à timbres — ; Lakunle est ridicule).
\item La femme comme enjeu : Sidi n'est pas passive, elle choisit. La dot = symbole du respect coutumier vs achat perçu par Lakunle.
\item Comédie satirique : rire des excès des deux camps pour mieux questionner l'avenir de l'Afrique postcoloniale.
\end{itemize}
\end{enumerate}
\textbf{Conclusion.} La lecture ouvroir pose le décor : Ilujinle, microcosme nigérian où se joue, en une journée, le sort d'une identité culturelle.
\end{exoresolu}

\courssub{Analyse du sujet et formulation de la problématique}

\begin{methode}[Analyser le sujet en 4 étapes]
\begin{enumerate}
\item \textbf{Lire le sujet plusieurs fois} et souligner les \cle{mots-clés} (termes opératoires, concepts centraux).
\item \textbf{Définir chaque mot-clé} : sens littéral (dictionnaire) puis sens contextuel (dans le sujet précis).
\item \textbf{Dégager le(s) présupposé(s)} : affirmation implicite, postulat de départ que le sujet invite à discuter (ex. \og tradition et progrès sont incompatibles \fg{}).
\item \textbf{Repérer la tension dialectique} : opposition entre deux pôles (ex. conservation / changement, individu / société, liberté / sécurité) qui structure le débat.
\end{enumerate}
\begin{aretenir}[Réflexes à retenir]
Un sujet bien analysé = mots-clés définis + présupposé explicité + tension identifiée. \textbf{Interdiction de rédiger} avant cette étape de \og mise au clair \fg{}.
\end{aretenir}
\end{methode}

\begin{exoresolu}[Analyse d'un sujet type Probatoire technique]
\textbf{Énoncé.} Sujet : \og La tradition est un obstacle au progrès. \fg{} Analysez : mots-clés, définitions, présupposé, tension.
\tcblower
\textbf{Solution détaillée.}
\begin{enumerate}
\item \textbf{Mots-clés} : \emph{tradition}, \emph{obstacle}, \emph{progrès}.
\item \textbf{Définitions contextuelles} :
\begin{itemize}
\item \emph{Tradition} : ensemble des coutumes, croyances, savoirs, pratiques transmis intergénérationnellement ; socle identitaire.
\item \emph{Obstacle} : entrave, frein, barrage à un mouvement vers l'avant ; mais aussi \emph{épreuve} qui oblige à se surpasser.
\item \emph{Progrès} : amélioration mesurable des conditions de vie (technique, sanitaire, éducatif, démocratique) ; notion évaluative et historique.
\end{itemize}
\item \textbf{Présupposé} : le sujet pose comme acquis que tradition et progrès sont antinomiques, que le passé pèse sur l'avenir comme un boulet. C'est cette \emph{évidence} qu'il faut soumettre à l'épreuve de l'argumentation.
\item \textbf{Tension} : \textbf{Thèse} (tradition = frein : pesanteur coutumière, refus de l'innovation, inégalités figées) \textbf{vs} \textbf{Antithèse} (tradition = ressource : valeurs éthiques, savoirs endogènes, identité motrice du développement). Le sujet est \cle{discutable}.
\end{enumerate}
\textbf{Lien œuvre programme.} Dans \emph{Le Lion et la Perle}, Baroka (tradition) achète une machine à timbres (modernité) pour tromper Sidi ; Lakunle (modernité) refuse la dot (tradition) mais veut posséder Sidi. La tension est incarnée, non abstraite.
\end{exoresolu}

\begin{methode}[Construire une problématique]
\begin{enumerate}
\item \textbf{Transformer la tension en question ouverte} : la problématique \emph{est} la question qui guide tout le devoir.
\item \textbf{Privilégier les interrogatifs opératoires} : \og Dans quelle mesure... ? \fg{} (nuance, gradation), \og Comment... ? \fg{} (mécanismes), \og Faut-il... ? \fg{} (choix éthique), \og Peut-on... ? \fg{} (possibilité/conditions).
\item \textbf{Tester la discutabilité} : si on peut répondre \og Oui \fg{} ou \og Non \fg{} sans argumenter, la question est mauvaise.
\item \textbf{Préfigurer le plan} : la problématique principale peut s'accompagner de 2 questions secondaires (une par grande partie) qui annoncent la structure dialectique.
\end{enumerate}
\begin{attention}[Piège fréquent]
La problématique n'est \textbf{jamais} une reformulation du sujet (\og Est-ce que la tradition est un obstacle au progrès ? \fg{} $\rightarrow$ trop binaire). Elle doit \emph{problématiser}, c'est-à-dire faire surgir la complexité.
\end{attention}
\end{methode}

\begin{exoresolu}[Formulation d'une problématique complète]
\textbf{Énoncé.} Sujet : \og La tradition est un obstacle au progrès. \fg{} Proposez une problématique principale et deux questions secondaires.
\tcblower
\textbf{Solution détaillée.}
\begin{enumerate}
\item \textbf{Problématique principale} : \og \textbf{Dans quelle mesure la tradition constitue-t-elle un obstacle au progrès ?} \fg{} (La formule \og dans quelle mesure \fg{} impose la nuance : ni tout obstacle, ni nul obstacle, mais une relation complexe).
\item \textbf{Questions secondaires (annonce du plan)} :
\begin{itemize}
\item \og \textbf{Quelles formes prend l'opposition entre coutumes et modernité dans nos sociétés ?} \fg{} (Partie I : Thèse — constater les freins).
\item \og \textbf{La tradition ne peut-elle pas, paradoxalement, fonder un progrès authentique ?} \fg{} (Partie II : Antithèse — révéler les ressources).
\end{itemize}
\item \textbf{Synthèse attendue} : dépasser l'alternative stérile pour penser un \og progrès enraciné \fg{} (Partie III).
\end{enumerate}
\textbf{Conclusion.} Cette problématique valide le plan dialectique (Thèse / Antithèse / Synthèse) adapté aux sujets polémiques.
\end{exoresolu}

\courssub{La recherche des idées}

\begin{methode}[Rechercher, trier et valider les arguments]
\begin{enumerate}
\item \textbf{Brainstorming (remue-méninges)} : sur brouillon, lister \emph{tout} ce qui vient (idées, exemples, citations, souvenirs de cours, actualité) sans censure ni ordre.
\item \textbf{Mobiliser les réservoirs obligatoires} :
\begin{itemize}
\item L'œuvre au programme (\emph{Le Lion et la Perle} : scènes précises, répliques, gestes).
\item Les textes du manuel / corpus de l'année.
\item Le contexte camerounais et africain (coutumes locales, tontines, chefferies, politiques publiques, actualité récente).
\item La culture générale (proverbes, histoire, philosophie, sciences).
\end{itemize}
\item \textbf{Classer par pôles} : deux colonnes \og Pour la thèse \fg{} / \og Pour l'antithèse \fg{} (ou selon les axes du plan thématique/analytique).
\item \textbf{Sélectionner et calibrer} : retenir 2 à 3 arguments solides par partie. \textbf{Un argument valide = Idée directrice + Explication (raisonnement) + Exemple précis et daté si possible.} Éliminer le hors-sujet, les doublons, les généralités non étayées.
\end{enumerate}
\begin{aretenir}[Formule Bac]
\textbf{Argument = Idée + Explication + Exemple concret.} Un argument sans exemple est une assertion, pas une démonstration. Au Probatoire, l'exemple \emph{fait} la note.
\end{aretenir}
\end{methode}

\begin{exoresolu}[Recherche et classement : 4 arguments illustrés]
\textbf{Énoncé.} Pour le sujet \og La tradition est un obstacle au progrès \fg{}, proposez 4 idées classées (Thèse / Antithèse) avec exemples.
\tcblower
\textbf{Solution détaillée.}
\begin{center}
\begin{tabularx}{\linewidth}{|X|X|}
\hline
\rowcolor{popblueL} \textbf{Thèse : La tradition freine le progrès} & \textbf{Antithèse : La tradition guide le progrès} \\
\hline
\textbf{Arg. 1} : Certaines pratiques coutumières perpétuent des inégalités (ex. \emph{mariages précoces / forcés} de filles dans l'Extrême-Nord et le Nord, freinant la scolarisation et l'autonomie économique). & \textbf{Arg. 1} : Les valeurs traditionnelles de \emph{solidarité} (tontines, \emph{njangi}, entraide villageoise, palabre) assurent la cohésion sociale et financent l'entrepreneuriat local là où l'État faillit. \\
\hline
\textbf{Arg. 2} : Le \emph{refus du changement technique} au nom des ancêtres retarde la productivité (ex. attachement exclusif à la houe / daba vs motorisation agricole ; résistance à la vaccination dans certaines communautés). & \textbf{Arg. 2} : Les \emph{savoirs endogènes} sont des viviers d'innovation (ex. pharmacopée traditionnelle $\rightarrow$ recherche pharmaceutique moderne ; architecture en banco $\rightarrow$ isolation thermique écologique ; Baroka qui détourne la machine à timbres pour asseoir son pouvoir). \\
\hline
\end{tabularx}
\end{center}
\textbf{Vérification} : Équilibre des colonnes (2/2), exemples ancrés (Cameroun, œuvre), pas de hors-sujet. Matière prête pour le plan détaillé.
\end{exoresolu}

\courssub{Les types de plans et l'élaboration du plan détaillé}

\begin{methode}[Choisir et construire le plan détaillé]
\begin{enumerate}
\item \textbf{Identifier le type de plan} imposé par la nature du sujet :
\begin{itemize}
\item \textbf{Plan dialectique} (Thèse / Antithèse / Synthèse) : sujet \emph{polémique}, affirmation discutable (\og Dans quelle mesure... \fg{}, \og Est-il vrai que... \fg{}). \emph{Standard pour le Bac technique.}
\item \textbf{Plan thématique} : sujet \emph{d'inventaire} ou d'analyse transversale (\og Les formes de... \fg{}, \og Les fonctions de... \fg{}).
\item \textbf{Plan analytique} (Constat / Causes / Conséquences \emph{ou} Solutions) : sujet \emph{phénoménologique} (\og Pourquoi... \fg{}, \og Quelles solutions pour... \fg{}).
\end{itemize}
\item \textbf{Rédiger le squelette} : pour chaque grande partie (I, II, III), écrire une \textbf{phrase-idée directrice} (titre de la partie).
\item \textbf{Détailler les sous-parties} : 2 ou 3 arguments par partie, chacun formulé en \textbf{phrase complète} (sujet + verbe + complément), suivi de l'exemple précis entre parenthèses.
\item \textbf{Rédiger les transitions} : une phrase logique reliant la fin d'une partie à l'amorce de la suivante (reprise + ouverture).
\item \textbf{Préparer Introduction et Conclusion} : 
\begin{itemize}
\item \emph{Intro} : Amorce (accroche) $\rightarrow$ Présentation du sujet $\rightarrow$ Problématique $\rightarrow$ Annonce du plan.
\item \emph{Conclusion} : Bilan synthétique (réponse à la problématique) $\rightarrow$ Ouverture (autre sujet, actualité, perspective).
\end{itemize}
\end{enumerate}
\begin{attention}[Conditions d'application au Bac]
Le plan détaillé \textbf{n'est pas la dissertation}. C'est un document de travail : phrases-arguments, pas mots isolés. Chaque partie doit répondre \emph{directement} à la problématique. Vérifier la \textbf{cohérence interne} (arguments non contradictoires dans une même partie) et la \textbf{progressivité} (I $\rightarrow$ II $\rightarrow$ III).
\end{attention}
\end{methode}

\begin{exoresolu}[Plan détaillé complet : Dissertation dialectique]
\textbf{Énoncé (type Bac technique).} Sujet : \og La tradition est un obstacle au progrès. \fg{} Rédigez le plan détaillé.
\tcblower
\textbf{Solution détaillée.}

\textbf{INTRODUCTION}
\begin{itemize}
\item \textbf{Amorce} : \og La tradition n'est pas le culte des cendres, mais la conservation du feu \fg{} (G. Mahler) — nos sociétés africaines vivent au rythme de cette flamme.
\item \textbf{Présentation du sujet} : L'affirmation « la tradition est un obstacle au progrès » oppose deux forces vives de l'histoire.
\item \textbf{Problématique} : \og Dans quelle mesure la tradition constitue-t-elle un obstacle au progrès ? \fg{}
\item \textbf{Annonce du plan} : Nous examinerons d'abord les freins réels (I), puis les ressources insoupçonnées (II), pour enfin penser leur articulation féconde (III).
\end{itemize}

\textbf{I. Thèse : La tradition peut freiner le progrès.}
\begin{itemize}
\item \textbf{Arg. 1} : La pérennisation de coutumes discriminatoires entrave l'émancipation humaine (ex. mariages précoces des filles dans le Septentrion camerounais, exclusion de l'héritage foncier pour les femmes dans certaines chefferies Bamiléké/Bassa).
\item \textbf{Arg. 2} : L'attachement identitaire rigide peut bloquer l'innovation technique (ex. Lakunle dans \emph{Le Lion et la Perle} dénonce le refus des villageois d'adopter des méthodes agricoles modernes ; résistance réelle aux OGM ou à la mécanisation par peur de « trahir les ancêtres »).
\end{itemize}
\textbf{Transition} : \og Cependant, réduire la tradition à un musée de vieilleries, c'est ignorer sa capacité d'adaptation et sa force normative. \fg{}

\textbf{II. Antithèse : La tradition est une richesse et un levier pour le progrès.}
\begin{itemize}
\item \textbf{Arg. 1} : Les valeurs éthiques traditionnelles (respect des aînés, solidarité clanique, parole donnée, palabre) fondent le \emph{vivre-ensemble} et la gouvernance locale (ex. fonctionnement des tontines \emph{njangi} pour le crédit populaire ; rôle des chefs dans la médiation des conflits fonciers).
\item \textbf{Arg. 2} : Les savoirs ancestraux nourrissent l'innovation contemporaine (ex. pharmacopée traditionnelle $\rightarrow$ molécules brevetées ; agroforesterie paysanne $\rightarrow$ agriculture durable ; Baroka qui s'approprie la machine à timbres — modernité — pour servir son dessein traditionnel).
\end{itemize}
\textbf{Transition} : \og Il ne s'agit pas de choisir l'un contre l'autre, mais de les articuler dans une dynamique de \emph{progrès enraciné}. \fg{}

\textbf{III. Synthèse : Vers un progrès enraciné dans la tradition (dépassement dialectique).}
\begin{itemize}
\item \textbf{Arg. 1} : Le vrai progrès opère un \emph{triage critique} : garder les valeurs universelles (dignité, justice), réformer les pratiques néfastes (ex. loi camerounaise sur l'âge du mariage, code des personnes et de la famille).
\item \textbf{Arg. 2} : L'identité culturelle devient avantage comparatif dans la mondialisation (ex. Sidi choisit Baroka : la modernité ne s'impose pas \emph{contre} la culture, mais \emph{par} elle ; industrie culturelle, tourisme mémoriel, diplomatie culturelle).
\end{itemize}

\textbf{CONCLUSION}
\begin{itemize}
\item \textbf{Bilan} : La tradition n'est un obstacle que si elle se fige en traditionalisme ; vivante, elle est le socle d'un progrès authentique.
\item \textbf{Ouverture} : \og Quelle place pour la chefferie traditionnelle dans la décentralisation camerounaise de demain ? \fg{} (Lien avec l'Éducation à la Citoyenneté).
\end{itemize}
\end{exoresolu}

\courssub{Langue : Conditions et situations de communication}

\begin{methode}[Analyser une situation de communication (Modèle de Jakobson)]
\begin{enumerate}
\item \textbf{Identifier les 6 facteurs} (schéma ci-dessous) :
\begin{itemize}
\item \textbf{Émetteur} (qui parle/écrit) $\rightarrow$ \textbf{Récepteur} (à qui ?) : \cle{Fonction émotive} / \cle{Fonction conative}.
\item \textbf{Message} (quoi ?) $\rightarrow$ \textbf{Code} (langue, jargon, registre) : \cle{Fonction poétique} (focus sur le message).
\item \textbf{Canal} (support : air, papier, numérique) $\rightarrow$ \cle{Fonction phatique} (contact).
\item \textbf{Contexte / Référent} (où, quand, quoi ?) $\rightarrow$ \cle{Fonction référentielle} (information).
\item \textbf{Fonction métalinguistique} : vérifier/préciser le code (ex. \og Que voulez-vous dire par « progrès » ? \fg{}).
\end{itemize}
\item \textbf{Vérifier les conditions de réussite} : code partagé (langue commune), canal opérationnel, contexte connu, attention du récepteur, intention claire.
\item \textbf{Qualifier la situation} :
\begin{itemize}
\item \emph{Formelle} (institutionnelle, contrainte, codifiée) vs \emph{Informelle} (libre, familiale, marchande).
\item \emph{Unilatérale} (discours, conférence, notice) vs \emph{Interactive} (dialogue, entretien, négociation).
\item \emph{Médiatisée} (média, réseau) vs \emph{Présentielle} (face-à-face).
\end{itemize}
\end{enumerate}
\end{methode}

\begin{popfigure}
\begin{tikzpicture}[node distance=1.5cm and 2cm, >=Stealth, font=\small]
% Nodes
\node (emetteur) [draw, rounded rectangle, fill=popblueL, thick] {ÉMETTEUR};
\node (message) [draw, rounded rectangle, fill=popgreenL, thick, right=of emetteur] {MESSAGE};
\node (recepteur) [draw, rounded rectangle, fill=popblueL, thick, right=of message] {RÉCEPTEUR};
\node (code) [draw, ellipse, fill=poporangeL, thick, above=of message] {CODE};
\node (canal) [draw, ellipse, fill=poporangeL, thick, below=of message] {CANAL};
\node (contexte) [draw, ellipse, fill=poppurpleL, thick, above right=of recepteur] {CONTEXTE / RÉFÉRENT};

% Arrows
\draw[->, thick, popdark] (emetteur) -- node[above, font=\tiny] {encode} (message);
\draw[->, thick, popdark] (message) -- node[above, font=\tiny] {décode} (recepteur);
\draw[->, thick, popmuted] (code) -- node[left, font=\tiny] {partagé} (message);
\draw[->, thick, popmuted] (canal) -- node[left, font=\tiny] {support} (message);
\draw[->, thick, popmuted] (contexte) -- node[above, font=\tiny] {situe} (message);

% Functions labels
\node[align=center, font=\tiny, text=popblue] at ($(emetteur)!0.5!(message) + (0, -0.8)$) {Fonction \textbf{Émotive}\\(exprime)};
\node[align=center, font=\tiny, text=popblue] at ($(message)!0.5!(recepteur) + (0, -0.8)$) {Fonction \textbf{Conative}\\(incite)};
\node[align=center, font=\tiny, text=popgreen] at ($(code)!0.5!(message) + (0, 0.5)$) {Fonction \textbf{Poétique}\\(esthétique)};
\node[align=center, font=\tiny, text=poporange] at ($(canal)!0.5!(message) + (0, -0.5)$) {Fonction \textbf{Phatique}\\(contact)};
\node[align=center, font=\tiny, text=poppurple] at ($(contexte)!0.5!(message) + (-0.5, 0)$) {Fonction \textbf{Référentielle}\\(informe)};
\node[align=center, font=\tiny, text=popdark] at ($(emetteur)!0.5!(code) + (-0.5, 0.5)$) {Fonction \textbf{Métalinguistique}\\(vérifie code)};
\end{tikzpicture}
\legende{Schéma des 6 fonctions du langage (Jakobson) appliqué aux facteurs de la communication.}
\end{popfigure}

\begin{exoresolu}[Analyse d'une situation : Marché de Mokolo]
\textbf{Énoncé.} Au marché de Mokolo (Yaoundé), une vendeuse de \og bayam-sellam \fg{} annonce ses prix en \emph{ewondo} à une cliente qui marchande en \emph{camfranglais}. Identifiez les 6 facteurs, les conditions et qualifiez la situation.
\tcblower
\textbf{Solution détaillée.}
\begin{itemize}
\item \textbf{Émetteur} : Vendeuse (puis cliente, alternance) ; \textbf{Récepteur} : Cliente (puis vendeuse).
\item \textbf{Message} : Prix, qualité, négociation, salutations rituelles.
\item \textbf{Code} : \emph{Ewondo} (langue locale commune) + \emph{Camfranglais} (vernaculaire urbain) $\rightarrow$ \textbf{code hybride partagé} (condition remplie).
\item \textbf{Canal} : Air (oral, voix), bruits ambiants (bruit de fond $\rightarrow$ risque de parasitage, compensé par gestes, répétition).
\item \textbf{Contexte} : Marché de Mokolo, jour de marché, matinée, espace public marchand, relations commerciales récurrentes.
\item \textbf{Conditions} : Code partagé (oui), Canal disponible (oui, mais bruité), Contexte partagé (oui), Attention (oui, enjeu financier), Intention claire (achat/vente).
\item \textbf{Fonctions dominantes} : \textbf{Conative} (marchander = inciter à baisser le prix), \textbf{Référentielle} (prix, quantité), \textbf{Phatique} (\og Mami, ça va ? \fg{} pour maintenir le contact).
\item \textbf{Qualification} : Situation \textbf{informelle} (cadre non institutionnel), \textbf{interactive} (échange tour à tour, feedback immédiat), \textbf{présentielle} (co-présence physique), \textbf{multilingue} (pratique camerounaise typique).
\end{itemize}
\textbf{Conclusion.} La communication réussit grâce à la compétence plurilingue des acteurs et au rituel marchand partagé.
\end{exoresolu}

\courssub{Langue : La ponctuation forte et lecture méthodique}

\begin{methode}[Maîtriser la ponctuation forte (marques de fin de phrase et structurantes)]
\begin{enumerate}
\item \textbf{Le point (.)} : Clôture une phrase \emph{déclarative} complète. La suivante commence par une majuscule. \emph{Attention :} point après abréviation (M., Dr, etc.) $\neq$ fin de phrase.
\item \textbf{Le point d'interrogation (?)} : Clôture une \emph{question directe} (ordre des mots inversé ou \og est-ce que \fg{}). \textbf{Interdit après question indirecte} (subordonnée complétive).
\item \textbf{Le point d'exclamation (!)} : Marque l'ordre (impératif), l'exclamation (émotion, surprise), l'interjection. À utiliser avec parcimonie en dissertation.
\item \textbf{Les deux-points (:)} : Annoncent une \emph{conséquence}, une \emph{explication}, une \emph{énumération} ou une \emph{citation}. La suite commence en minuscule (sauf citation complète ou nom propre).
\item \textbf{Les points de suspension (...)} : 3 points exactement. Marquent l'\emph{ellipse} (liste ouverte), l'\emph{hésitation}, la \emph{rupture} (dialogue), la \emph{suspense}. Jamais plus de 3.
\end{enumerate}
\begin{attention}[Erreurs sanctionnées au Bac]
\begin{enumerate}
\item \textbf{Point d'interrogation après question indirecte} : \og Il demande si tu viens? \fg{} $\rightarrow$ \textbf{Faute}. Correct : \og Il demande si tu viens. \fg{}
\item \textbf{Majuscule après deux-points} (hors citation) : \og Il a dit : Les élèves sont là. \fg{} $\rightarrow$ \textbf{Faute} (sauf citation). Correct : \og Il a dit : les élèves sont là. \fg{}
\item \textbf{Oubli des deux-points avant énumération/citation} : \og Trois arguments : l'école, la santé, l'emploi. \fg{} (Obligatoire).
\item \textbf{Points de suspension « décoratifs »} en fin de phrase affirmative : \og Il est parti... \fg{} (inadapté en rédaction formelle).
\end{enumerate}
\end{attention}
\end{methode}

\begin{exoresolu}[Ponctuation forte : Exercice d'application]
\textbf{Énoncé.} Ponctuez correctement le texte ci-dessous (ponctuation forte uniquement) :
\og le professeur demanda pourquoi la tradition freinerait le progrès ( ) les élèves citèrent trois exemples ( ) les mariages précoces ( ) le refus des techniques nouvelles ( ) et le rejet de l'école ( ) quel débat ( ) \fg{}
\tcblower
\textbf{Solution détaillée.}
\begin{enumerate}
\item \og ...freinerait le progrès\textbf{.} \fg{} (Question \emph{indirecte} = subordonnée complétive de \og demanda \fg{} $\rightarrow$ point final).
\item \og ...citèrent trois exemples\textbf{:} \fg{} (Annonce d'énumération $\rightarrow$ deux-points).
\item \og ...les mariages précoces\textbf{,} le refus des techniques nouvelles\textbf{,} et le rejet de l'école\textbf{.} \fg{} (Virgules = ponctuation faible interne ; point final = fin de phrase déclarative).
\item \og \textbf{Quel débat!} \fg{} (Exclamation admirative $\rightarrow$ point d'exclamation).
\end{enumerate}
\textbf{Texte corrigé} : \og Le professeur demanda pourquoi la tradition freinerait le progrès. Les élèves citèrent trois exemples : les mariages précoces, le refus des techniques nouvelles et le rejet de l'école. Quel débat! \fg{}
\end{exoresolu}

\begin{methode}[Lecture méthodique 1 : La fiche de lecture analytique (œuvre intégrale)]
\begin{enumerate}
\item \textbf{Fiche d'identité} : Titre, Auteur, Date, Genre, Mouvement, Contexte historique/culturel.
\item \textbf{Résumé structuré} : Par acte/scène ou séquence narrative (début, péripéties, dénouement). Neutre, au présent.
\item \textbf{Personnages principaux} : Portrait physique/moral, fonction dramatique (protagoniste, antagoniste, confident), évolution.
\item \textbf{Thèmes majeurs} : 3 à 5 thèmes avec citations courtes (référence acte/scène).
\item \textbf{Procédés d'écriture / Style} : Registres (comique, tragique, satirique), figures de style, langue (niveau, oralité), structure dramatique.
\item \textbf{Apport personnel / Jugement critique} : Ce que l'œuvre m'apporte, lien avec le programme (tradition/modernité), actualité.
\end{enumerate}
\begin{aretenir}[Contrôle de lecture]
Le contrôle porte sur la \textbf{fiche} (connaissance de l'intrigue, personnages, thèmes) et la \textbf{capacité à mobiliser} l'œuvre dans une argumentation (dissertation, commentaire). La \og négociation du projet d'étude \fg{} valide le choix des axes d'approfondissement (ex. : \og La figure du trickster chez Baroka \fg{}, \og La satire de l'occidentalisé chez Lakunle \fg{}).
\end{aretenir}
\end{methode}

\courssub{Contrôle de lecture, compte rendu et négociation du projet d'étude}

\begin{exoresolu}[Modèle de fiche de lecture synthétique pour \emph{Le Lion et la Perle}]
\textbf{Énoncé.} Remplissez la fiche de lecture type pour le contrôle de l'œuvre intégrale.
\tcblower
\textbf{Solution détaillée.}
\begin{center}
\begin{tabularx}{\linewidth}{|l|X|}
\hline
\rowcolor{popblueL} \textbf{Champ} & \textbf{Contenu attendu (exemples)} \\
\hline
\textbf{Titre / Auteur / Date} & \emph{Le Lion et la Perle}, Wole Soyinka, 1963 (création), 1964 (publication). \\
\hline
\textbf{Genre / Mouvement} & Théâtre / Comédie satirique postcoloniale / Théâtre yoruba modernisé. \\
\hline
\textbf{Contexte} & Nigéria pré-indépendance (1960), tension tradition (coutumes yoruba) / modernité (école britannique, administration). \\
\hline
\textbf{Résumé (3 actes)} & \textbf{Matin} : Lakunle courtise Sidi, refuse la dot. Photographe \og stranger \fg{} publie photos de Sidi. \textbf{Midi} : Sadiku annonce veuveté feinte de Baroka. Sidi va moquer Baroka. \textbf{Soir} : Baroka séduit Sidi (machine à timbres). Lakunle arrive trop tard. Sidi épouse Baroka. \\
\hline
\textbf{Personnages clés} & \textbf{Baroka} : Bale, 62 ans, rusé, vital, adaptatif (tradition vivante). \textbf{Lakunle} : Instituteur, dogmatique, ridicule, modernité superficielle. \textbf{Sidi} : Belle, fière, pragmatique, enjeu puis actrice. \textbf{Sadiku} : Épouse senior, médiatrice loyale mais naïve. \\
\hline
\textbf{Thèmes majeurs} & 1. Tradition vs Modernité (nuancée). 2. Pouvoir de la séduction / Image (photos). 3. Condition féminine / Dot / Mariage. 4. Ruse / Trickster (Baroka). 5. Langue / Communication (proverbes, anglais, pidgin). \\
\hline
\textbf{Procédés marquants} & Proverbes yoruba traduits (oralité), chansons/danses (rituel), satire (caricature Lakunle), ironie dramatique (Sidi ne sait pas que Baroka ment), didascalies scéniques précises. \\
\hline
\textbf{Projets d'étude négociés} & \textit{Axe 1} : Baroka, figure du \emph{trickster} moderne ? \textit{Axe 2} : La dot, obstacle ou fondement du lien social ? \textit{Axe 3} : Le corps de la femme comme espace de guerre culturelle. \\
\hline
\end{tabularx}
\end{center}
\textbf{Conclusion.} Cette fiche, apprise par cœur, est la \og boîte à outils \fg{} pour toute dissertation ou commentaire portant sur l'œuvre.
\end{exoresolu}

\begin{attention}[Les 5 erreurs qui coûtent des points au Probatoire / Bac Technique]
\begin{enumerate}
\item \textbf{Recopier le sujet comme problématique} : La problématique doit être une \textbf{vraie question ouverte} (\og Dans quelle mesure... ? \fg{}, \og Comment... ? \fg{}), jamais une paraphrase de l'énoncé.
\item \textbf{Un plan sans exemples concrets} : Chaque argument = Idée + Explication + \textbf{Exemple précis} (œuvre, Cameroun, actualité datée). Argument nu = note faible.
\item \textbf{Le hors-sujet masqué} : Des idées justes mais qui ne répondent \emph{pas} à la problématique posée. \textbf{Vérifier systématiquement} : \og Cet argument répond-il à la question ? \fg{}.
\item \textbf{Confondre les types de plans} : Sujet polémique $\rightarrow$ \textbf{Plan dialectique} (Thèse/Antithèse/Synthèse). Imposer un plan thématique à un sujet de débat = démonstration affaiblie, note de structure basse.
\item \textbf{Négliger la langue (ponctuation, transitions, syntaxe)} : Point d'interrogation après question indirecte, absence de deux-points avant énumération, parties juxtaposées sans phrases de liaison, fautes d'accord majeures : \textbf{sanctionnées sur la note de langue ET de cohérence}.
\end{enumerate}
\end{attention}

\newpage

\coursec{Exercices}

\courssub{Je vérifie mes connaissances}

\begin{exercice}[QCM : la dissertation]
Pour chaque question, choisis la bonne réponse.
\begin{enumerate}
\item L'analyse d'un sujet de dissertation commence par :
\begin{enumerate}
\item la rédaction de l'introduction ;
\item le dégagement des mots-clés et du thème ;
\item la rédaction de la conclusion.
\end{enumerate}
\item La problématique est :
\begin{enumerate}
\item une question directrice née de l'analyse du sujet ;
\item le titre de l'œuvre étudiée ;
\item une citation d'auteur.
\end{enumerate}
\item Le plan dialectique comporte :
\begin{enumerate}
\item deux parties ;
\item trois parties (thèse, antithèse, synthèse) ;
\item cinq parties.
\end{enumerate}
\item Un plan détaillé contient :
\begin{enumerate}
\item uniquement les titres des parties ;
\item les parties, les sous-parties, les idées et les exemples ;
\item le texte intégral de la dissertation.
\end{enumerate}
\end{enumerate}
\end{exercice}

\begin{exercice}[Vrai ou faux]
Réponds par vrai ou faux, puis justifie ta réponse en une phrase.
\begin{enumerate}
\item La recherche des idées se fait avant l'analyse du sujet.
\item Le plan thématique convient à un sujet qui appelle un débat d'idées avec confrontation de deux opinions opposées.
\item Les paratextes d'un livre comprennent le titre, le nom de l'auteur et la quatrième de couverture.
\item Le point d'exclamation peut exprimer l'ordre, la surprise ou l'indignation dans un dialogue.
\end{enumerate}
\end{exercice}

\begin{exercice}[Définitions]
Définis en une ou deux phrases chacun des termes suivants :
\begin{enumerate}
\item la problématique ;
\item le plan détaillé ;
\item le paratexte ;
\item la situation de communication.
\end{enumerate}
\end{exercice}

\begin{exercice}[Identifier les mots-clés]
On donne le sujet suivant : \emph{Le travail est-il pour l'homme une contrainte ou une source d'épanouissement ?}
\begin{enumerate}
\item Relève les mots-clés du sujet.
\item Indique le domaine (thème) auquel appartient ce sujet.
\item Repère la forme du sujet (question directe, affirmation à discuter, citation) et dis ce que cela implique pour la problématique.
\end{enumerate}
\end{exercice}

\courssub{Je m'entraîne}

\begin{exercice}[Formuler une problématique]
Pour chacun des sujets suivants, dégage les mots-clés, précise le thème et formule une problématique sous forme d'une question directrice.
\begin{enumerate}
\item \emph{Faut-il interdire les téléphones portables à l'école ?}
\item \emph{La tradition est-elle un obstacle au progrès ?}
\end{enumerate}
\end{exercice}

\begin{exercice}[Trier et regrouper des idées]
Sujet : \emph{Le sport est-il utile à la formation du jeune ?} Voici huit idées issues d'une recherche : (a) le sport fatigue les élèves ; (b) le sport développe l'esprit d'équipe ; (c) le sport prend du temps sur les études ; (d) le sport entretient la santé ; (e) le sport apprend le respect des règles ; (f) le sport peut provoquer des blessures ; (g) le sport procure du plaisir ; (h) le sport peut ouvrir une carrière professionnelle.
\begin{enumerate}
\item Classe ces idées en deux groupes : avantages et inconvénients.
\item Regroupe les idées proches et donne un titre à chaque sous-groupe.
\item Propose une troisième partie possible pour un plan dialectique.
\end{enumerate}
\end{exercice}

\begin{exercice}[Ponctuation forte]
Ponctue correctement le passage suivant (points, points d'interrogation, points d'exclamation, points de suspension, deux-points), puis justifie trois de tes choix par leurs effets de sens :

\emph{Lakunle s'approcha de Sidi et s'écria tu refuses encore le mariage sans dot mais réfléchis un peu veux-tu rester toute ta vie attachée aux vieilles coutumes le village entier se moquera de nous écoute-moi je t'en supplie}

\begin{enumerate}
\item Ponctue le texte.
\item Justifie l'emploi d'un point d'exclamation, d'un point d'interrogation et d'un point final.
\end{enumerate}
\end{exercice}

\begin{exercice}[Conditions de communication]
Dans \emph{Le Lion et la Perle} de Wole Soyinka, Lakunle, l'instituteur, tente de convaincre Sidi de renoncer au prix de la dot.
\begin{enumerate}
\item Identifie l'émetteur, le récepteur, le message et le canal de cette communication.
\item Précise le lieu, le moment et le rapport entre les personnages (situation de communication).
\item Dis si la communication réussit et explique pourquoi.
\end{enumerate}
\end{exercice}

\courssub{Je me prépare au Bac}

\begin{exercice}[Type Bac : analyse complète d'un sujet]
Sujet : \emph{L'école est-elle la seule voie de la réussite ?}
\begin{enumerate}
\item Relève les mots-clés du sujet et dégage le thème.
\item Formule la problématique sous forme d'une question directrice.
\item Note au brouillon au moins six idées (arguments et exemples tirés de la vie camerounaise) susceptibles d'alimenter la réflexion.
\item Propose deux types de plans possibles pour traiter ce sujet (par exemple plan dialectique et plan thématique), en indiquant pour chacun les titres des grandes parties.
\item Justifie le choix du plan que tu retiens et rédige le plan détaillé correspondant (parties, sous-parties, idées et exemples).
\end{enumerate}
\end{exercice}

\begin{exercice}[Type Probatoire : construire un plan dialectique détaillé]
Sujet : \emph{Les réseaux sociaux : progrès ou danger pour la jeunesse camerounaise ?}
Voici huit idées mélangées issues de la recherche :
\begin{enumerate}
\item[a.] les réseaux sociaux facilitent la communication rapide entre les jeunes ;
\item[b.] certains jeunes y passent des heures au détriment des études ;
\item[c.] ils permettent de s'informer et de suivre des cours en ligne ;
\item[d.] ils exposent les jeunes aux fausses informations et aux mauvaises rencontres ;
\item[e.] ils aident les jeunes entrepreneurs à faire connaître leurs produits ;
\item[f.] ils peuvent provoquer la moquerie et le harcèlement entre élèves ;
\item[g.] ils créent des liens d'amitié au-delà des frontières ;
\item[h.] un usage encadré et modéré peut concilier avantages et prudence.
\end{enumerate}
Questions :
\begin{enumerate}
\item Trie ces huit idées en deux groupes (progrès / danger).
\item Regroupe les idées proches et formule une sous-partie pour chaque groupe.
\item Rédige un plan dialectique détaillé (thèse, antithèse, synthèse) avec sous-parties, idées et exemples.
\item Rédige la problématique et une phrase d'annonce de plan pour l'introduction.
\end{enumerate}
\end{exercice}

\begin{exercice}[Type Bac : réception et langue autour du \emph{Lion et la Perle}]
On donne les éléments de paratexte suivants : titre \emph{Le Lion et la Perle} ; auteur Wole Soyinka, dramaturge nigérian, prix Nobel de littérature ; quatrième de couverture annonçant une pièce de théâtre où le chef Baroka (le Lion) et le jeune instituteur Lakunle se disputent la belle Sidi (la Perle) dans un village nigérian partagé entre tradition et modernité.
\begin{enumerate}
\item Dégage trois informations fournies par ce paratexte (titre, auteur, quatrième de couverture).
\item Formule deux hypothèses de lecture sur le sens du titre et sur le conflit annoncé.
\item Confronte ces hypothèses à l'extrait étudié en classe (scène où Lakunle tente de convaincre Sidi de renoncer à la dot) : les hypothèses sont-elles confirmées ou infirmées ? Justifie.
\item Dans cette scène, identifie les conditions de communication (émetteur, récepteur, message, canal) et décris la situation de communication (lieu, moment, rapport entre les personnages).
\item Releve dans l'extrait deux signes de ponctuation forte et explique leur effet de sens sur le ton de la scène.
\end{enumerate}
\end{exercice}

\newpage

\coursec{Corrigés des exercices}

\coursec{Corrigés des exercices — Leçon 17 : Produire un plan détaillé d'une dissertation}

\begin{corrige}[Exercice 1 — QCM : la dissertation]
\begin{enumerate}
\item Réponse \textbf{b} : le dégagement des mots-clés et du thème. L'analyse du sujet commence toujours par la lecture attentive de l'énoncé : on repère les mots importants (mots-clés) et on dégage le domaine général (le thème). La rédaction de l'introduction et de la conclusion ne vient qu'après cette étape.
\item Réponse \textbf{a} : une question directrice née de l'analyse du sujet. La problématique est la question centrale que le sujet pose et à laquelle toute la dissertation devra répondre.
\item Réponse \textbf{b} : trois parties (thèse, antithèse, synthèse). Le plan dialectique confronte d'abord une opinion (thèse), puis l'opinion opposée (antithèse), avant de dépasser l'opposition dans une synthèse.
\item Réponse \textbf{b} : les parties, les sous-parties, les idées et les exemples. Le plan détaillé est le squelette complet du devoir : il annonce chaque partie, chaque sous-partie, chaque argument et chaque exemple, sans rédiger le texte intégral.
\end{enumerate}
\end{corrige}

\begin{corrige}[Exercice 2 — Vrai ou faux]
\begin{enumerate}
\item \textbf{Faux.} La recherche des idées se fait \emph{après} l'analyse du sujet et la formulation de la problématique, car ce sont elles qui orientent et guident la recherche.
\item \textbf{Faux.} C'est le \emph{plan dialectique} (thèse, antithèse, synthèse) qui convient à un sujet appelant la confrontation de deux opinions opposées ; le plan thématique convient plutôt à l'étude des différents aspects d'une notion.
\item \textbf{Vrai.} Les paratextes sont tous les éléments qui entourent le texte : titre, nom de l'auteur, couverture, quatrième de couverture, préface, illustrations.
\item \textbf{Vrai.} Le point d'exclamation traduit un sentiment fort : l'ordre, la surprise, la colère, l'indignation ou la joie, notamment dans les dialogues.
\end{enumerate}
\end{corrige}

\begin{corrige}[Exercice 3 — Définitions]
\begin{enumerate}
\item \textbf{La problématique} : c'est la question directrice qui naît de l'analyse du sujet et qui résume le problème à résoudre ; elle guide toute la réflexion de la dissertation.
\item \textbf{Le plan détaillé} : c'est l'organisation complète et écrite du devoir, présentant les grandes parties, les sous-parties, les idées (arguments) et les exemples qui serviront à la rédaction.
\item \textbf{Le paratexte} : c'est l'ensemble des éléments qui accompagnent et entourent un texte (titre, nom de l'auteur, couverture, quatrième de couverture, préface) et qui préparent sa lecture.
\item \textbf{La situation de communication} : c'est l'ensemble des circonstances dans lesquelles s'échange un message : le lieu, le moment, les personnes présentes et le rapport qu'elles entretiennent.
\end{enumerate}
\end{corrige}

\begin{corrige}[Exercice 4 — Identifier les mots-clés]
\begin{enumerate}
\item Les mots-clés sont : \cle{travail}, \cle{homme}, \cle{contrainte}, \cle{épanouissement}.
\item Le sujet appartient au domaine (thème) du \textbf{travail et de la condition humaine}.
\item Il s'agit d'une \textbf{question directe} à double terme (alternative « contrainte \emph{ou} épanouissement »). Cela implique que la problématique doit reprendre cette alternative sous forme d'interrogation et que le devoir devra examiner les deux points de vue avant de conclure. Problématique possible : \emph{Le travail est-il seulement une contrainte subie par l'homme, ou peut-il être aussi une source d'épanouissement ?}
\end{enumerate}
\end{corrige}

\begin{corrige}[Exercice 5 — Formuler une problématique]
\begin{enumerate}
\item \emph{Faut-il interdire les téléphones portables à l'école ?}
\begin{itemize}
\item Mots-clés : \cle{interdire}, \cle{téléphones portables}, \cle{école}.
\item Thème : l'école et les nouvelles technologies.
\item Problématique : \emph{Le téléphone portable est-il un outil utile à l'élève ou une nuisance qu'il faut bannir de l'école ?}
\end{itemize}
\item \emph{La tradition est-elle un obstacle au progrès ?}
\begin{itemize}
\item Mots-clés : \cle{tradition}, \cle{obstacle}, \cle{progrès}.
\item Thème : tradition et modernité (thème central de \emph{Le Lion et la Perle}).
\item Problématique : \emph{La tradition freine-t-elle le progrès, ou peut-elle au contraire l'accompagner et l'enrichir ?}
\end{itemize}
\end{enumerate}
\end{corrige}

\begin{corrige}[Exercice 6 — Trier et regrouper des idées]
\begin{enumerate}
\item Classement en deux groupes :
\begin{itemize}
\item \textbf{Avantages} : (b) esprit d'équipe ; (d) santé ; (e) respect des règles ; (g) plaisir ; (h) carrière professionnelle.
\item \textbf{Inconvénients} : (a) fatigue ; (c) temps pris sur les études ; (f) blessures.
\end{itemize}
\item Regroupement et titres des sous-groupes :
\begin{itemize}
\item Avantages : \emph{le sport forme le caractère} (b, e) ; \emph{le sport entretient le corps et procure du plaisir} (d, g) ; \emph{le sport ouvre des perspectives d'avenir} (h).
\item Inconvénients : \emph{le sport nuit à la scolarité} (a, c) ; \emph{le sport comporte des risques physiques} (f).
\end{itemize}
\item Troisième partie possible (synthèse) : \emph{Un sport pratiqué avec modération et bien encadré concilie formation du jeune et réussite scolaire.}
\end{enumerate}
\end{corrige}

\begin{corrige}[Exercice 7 — Ponctuation forte]
\begin{enumerate}
\item Texte ponctué :

\emph{Lakunle s'approcha de Sidi et s'écria : « Tu refuses encore le mariage sans dot ! Mais réfléchis un peu : veux-tu rester toute ta vie attachée aux vieilles coutumes ? Le village entier se moquera de nous\ldots\ Écoute-moi, je t'en supplie ! »}
\item Justifications :
\begin{itemize}
\item \textbf{Point d'exclamation} après « sans dot ! » : il exprime l'indignation et l'exaspération de Lakunle devant le refus de Sidi ; il donne à la réplique un ton véhément.
\item \textbf{Point d'interrogation} après « vieilles coutumes ? » : il signale une question directe par laquelle Lakunle tente de faire réfléchir Sidi ; c'est un procédé de persuasion.
\item \textbf{Points de suspension} après « se moquera de nous\ldots » : ils laissent planer la menace et donnent à Lakunle le temps de laisser son argument produire son effet ; le point final, s'il était utilisé, marquerait une affirmation catégorique, une prédiction présentée comme certaine.
\end{itemize}
\end{enumerate}
\end{corrige}

\begin{corrige}[Exercice 8 — Conditions de communication]
\begin{enumerate}
\item Éléments de la communication (schéma de Jakobson) :
\begin{itemize}
\item \textbf{Émetteur} : Lakunle, le jeune instituteur ;
\item \textbf{Récepteur} : Sidi, la belle jeune fille du village (la Perle) ;
\item \textbf{Message} : il faut renoncer au prix de la dot, coutume que Lakunle juge dépassée ;
\item \textbf{Canal} : la parole, dans un échange en face à face (voie orale) ;
\item \textbf{Code} : la langue parlée (anglais dans l'original, français en traduction) ;
\item \textbf{Contexte (référent)} : la coutume de la dot, la place du village, le regard de la communauté.
\end{itemize}
\item Situation de communication : la scène se déroule dans le \textbf{village} d'Ilujinle, en plein jour, près de la place du marché ; le rapport entre les personnages est celui d'un prétendant moderniste face à une jeune fille attachée aux coutumes : rapport de séduction et de confrontation, marqué par le mépris de Lakunle pour la tradition.
\item La communication \textbf{échoue} : Sidi refuse les arguments de Lakunle, car elle tient à la dot, qui fait sa fierté et sa valeur aux yeux du village. L'échec vient de ce que l'émetteur méprise les valeurs de son récepteur au lieu de les respecter : un message ne passe pas quand il blesse celui qui le reçoit.
\end{enumerate}
\end{corrige}

\begin{corrige}[Exercice 9 — Type Bac : analyse complète d'un sujet]
\begin{enumerate}
\item \textbf{Mots-clés} : \cle{école}, \cle{seule voie}, \cle{réussite}. \textbf{Thème} : l'éducation et la réussite sociale.
\item \textbf{Problématique} : \emph{L'école est-elle l'unique chemin qui mène à la réussite, ou existe-t-il d'autres voies légitimes ?}
\item \textbf{Idées au brouillon} (au moins six) :
\begin{itemize}
\item l'école donne des connaissances et des diplômes exigés par les employeurs ;
\item l'école forme l'esprit critique et ouvre l'accès aux fonctions publiques camerounaises (concours administratifs) ;
\item beaucoup de diplômés restent au chômage ;
\item l'artisanat, l'agriculture et le commerce permettent de réussir sans longues études (exemple : les commerçants prospères de nos marchés) ;
\item le sport et la musique offrent des carrières (exemple : les footballeurs camerounais) ;
\item la formation professionnelle et technique est une voie de réussite reconnue ;
\item la réussite dépend aussi du travail, de la persévérance et de l'honnêteté.
\end{itemize}
\item \textbf{Deux plans possibles} :
\begin{itemize}
\item \emph{Plan dialectique} : I. L'école est la voie principale de la réussite (thèse) ; II. Mais d'autres voies existent hors de l'école (antithèse) ; III. L'école reste précieuse, mais la réussite dépend surtout de l'effort et des talents de chacun (synthèse).
\item \emph{Plan thématique} : I. Ce que l'école apporte à la réussite ; II. Les autres voies de la réussite ; III. Les conditions personnelles de la réussite.
\end{itemize}
\item \textbf{Choix et plan détaillé} : on retient le plan dialectique, car le sujet oppose deux opinions (« seule voie » ou non) et appelle un débat.

\textbf{I. L'école est une voie essentielle de la réussite.}
\begin{itemize}
\item A. Elle transmet les connaissances et les diplômes (ex. : les concours de la fonction publique exigent des diplômes).
\item B. Elle forme le jugement et ouvre l'esprit (ex. : médecins, ingénieurs, enseignants formés à l'école).
\end{itemize}
\textbf{II. Cependant, l'école n'est pas la seule voie.}
\begin{itemize}
\item A. Des diplômés restent sans emploi (ex. : les jeunes diplômés au chômage dans nos villes).
\item B. L'artisanat, le commerce, l'agriculture, le sport et les arts mènent aussi à la réussite (ex. : artisans et commerçants prospères, footballeurs camerounais).
\end{itemize}
\textbf{III. La réussite dépend avant tout de l'effort et de l'usage des talents.}
\begin{itemize}
\item A. L'école est un atout précieux qu'il faut valoriser.
\item B. Mais le travail, la persévérance et l'honnêteté décident de la réussite, à l'école comme hors de l'école.
\end{itemize}
\end{enumerate}

\textbf{Barème indicatif (sur 20)} : mots-clés et thème (3 pts) ; problématique bien formulée (3 pts) ; recherche des idées, au moins six idées pertinentes (4 pts) ; deux types de plans correctement présentés (4 pts) ; choix justifié et plan détaillé cohérent (6 pts).

\textbf{Points de vigilance} : la problématique doit être une \emph{question}, pas une affirmation ; les exemples doivent être concrets et tirés de la vie camerounaise ; le plan détaillé annonce les idées sans les rédiger ; la synthèse doit dépasser l'opposition, non la répéter.
\end{corrige}

\begin{corrige}[Exercice 10 — Type Probatoire : construire un plan dialectique détaillé]
\begin{enumerate}
\item \textbf{Tri des idées} :
\begin{itemize}
\item \emph{Progrès} : a (communication rapide), c (information et cours en ligne), e (promotion des produits des jeunes entrepreneurs), g (amitiés au-delà des frontières).
\item \emph{Danger} : b (temps perdu au détriment des études), d (fausses informations et mauvaises rencontres), f (moqueries et harcèlement).
\item L'idée h (usage encadré et modéré) prépare la synthèse.
\end{itemize}
\item \textbf{Regroupement et sous-parties} :
\begin{itemize}
\item Progrès : \emph{les réseaux sociaux facilitent la communication et les liens} (a, g) ; \emph{ils sont un outil d'apprentissage et d'activité économique} (c, e).
\item Danger : \emph{ils nuisent aux études} (b) ; \emph{ils exposent aux mauvaises influences et au harcèlement} (d, f).
\end{itemize}
\item \textbf{Plan dialectique détaillé} :

\textbf{I. Les réseaux sociaux sont un progrès pour la jeunesse camerounaise (thèse).}
\begin{itemize}
\item A. Ils facilitent la communication et créent des liens (idées a, g ; ex. : échanges rapides entre élèves de Yaoundé et de Garoua, amitiés avec des jeunes d'autres pays).
\item B. Ils servent l'apprentissage et l'initiative économique (idées c, e ; ex. : cours en ligne, jeune vendeuse qui fait connaître ses produits sur les réseaux).
\end{itemize}
\textbf{II. Mais ils présentent aussi de réels dangers (antithèse).}
\begin{itemize}
\item A. Ils nuisent aux études par le temps qu'ils absorbent (idée b ; ex. : élève qui passe ses nuits en ligne et échoue en classe).
\item B. Ils exposent aux fausses informations, aux mauvaises rencontres et au harcèlement (idées d, f ; ex. : rumeurs, moqueries entre élèves).
\end{itemize}
\textbf{III. Un usage encadré et modéré permet de concilier avantages et prudence (synthèse).}
\begin{itemize}
\item A. Les jeunes doivent limiter le temps passé en ligne et vérifier les informations (idée h).
\item B. Parents et école doivent éduquer à un usage responsable (ex. : règles d'utilisation fixées en famille et au lycée).
\end{itemize}
\item \textbf{Problématique} : \emph{Les réseaux sociaux représentent-ils un progrès au service de la jeunesse camerounaise ou un danger qui la menace ?}

\textbf{Annonce de plan} : \emph{Nous verrons d'abord que les réseaux sociaux constituent un progrès pour les jeunes, puis qu'ils comportent de véritables dangers, avant de montrer qu'un usage encadré permet d'en tirer profit sans en subir les méfaits.}
\end{enumerate}

\textbf{Barème indicatif (sur 20)} : tri correct des idées (4 pts) ; regroupement et sous-parties bien formulées (4 pts) ; plan dialectique détaillé complet et cohérent (8 pts) ; problématique et annonce de plan (4 pts).

\textbf{Points de vigilance} : respecter l'ordre thèse / antithèse / synthèse ; ne pas mélanger dans une même partie des idées favorables et défavorables ; l'idée h appartient à la synthèse, pas à la thèse ; chaque sous-partie doit contenir une idée \emph{et} un exemple.
\end{corrige}

\begin{corrige}[Exercice 11 — Type Bac : réception et langue autour de \emph{Le Lion et la Perle}]
\begin{enumerate}
\item \textbf{Trois informations du paratexte} :
\begin{itemize}
\item le \emph{titre} annonce deux personnages symboliques : un « Lion » (force, pouvoir) et une « Perle » (beauté, objet précieux) ;
\item l'\emph{auteur} est Wole Soyinka, dramaturge nigérian et prix Nobel de littérature, ce qui promet une œuvre de valeur reconnue ;
\item la \emph{quatrième de couverture} indique le genre (pièce de théâtre), le cadre (un village nigérian) et le conflit : Baroka et Lakunle se disputent Sidi, sur fond d'opposition entre tradition et modernité.
\end{itemize}
\item \textbf{Hypothèses de lecture} :
\begin{itemize}
\item le « Lion » désigne probablement le chef Baroka, homme âgé et puissant, et la « Perle » la belle Sidi, convoitée par tous ;
\item le conflit annoncé opposera vraisemblablement la tradition (incarnée par Baroka et la dot) à la modernité (incarnée par Lakunle, l'instituteur).
\end{itemize}
\item \textbf{Confrontation à l'extrait} : les hypothèses sont \textbf{confirmées}. Dans la scène étudiée, Lakunle tente bien de convaincre Sidi de renoncer à la dot, qu'il juge arriérée : il incarne la modernité, tandis que Sidi défend la coutume, fière de sa valeur. La rivalité autour de la « Perle » et l'opposition tradition / modernité sont bien au cœur de la pièce.
\item \textbf{Conditions et situation de communication} : l'émetteur est Lakunle, le récepteur Sidi, le message est le refus de la dot, le canal est la parole en face à face. La scène se déroule dans le village, près de la place du marché, en plein jour ; le rapport entre les personnages est un rapport de séduction conflictuelle : Lakunle veut imposer ses idées modernes à une jeune fille attachée aux traditions, ce qui provoque le rejet du message.
\item \textbf{Ponctuation forte et effets de sens} :
\begin{itemize}
\item le \textbf{point d'exclamation} (par exemple dans les reproches de Lakunle) traduit l'exaspération et la véhémence de l'instituteur : le ton est passionné, presque agressif ;
\item le \textbf{point d'interrogation} (questions répétées de Lakunle ou réponses de Sidi) marque l'affrontement : chacun interroge et met l'autre en difficulté, ce qui donne à la scène un ton vif et combatif.
\end{itemize}
\end{enumerate}

\textbf{Barème indicatif (sur 20)} : informations du paratexte (3 pts) ; hypothèses de lecture (4 pts) ; confrontation justifiée à l'extrait (4 pts) ; conditions et situation de communication (5 pts) ; relevé et explication de la ponctuation forte (4 pts).

\textbf{Points de vigilance} : distinguer clairement \emph{paratexte} (ce qui entoure le texte) et \emph{texte} ; une hypothèse de lecture doit être formulée au conditionnel ou comme une supposition ; pour la communication, citer les quatre éléments (émetteur, récepteur, message, canal) \emph{et} la situation (lieu, moment, rapport) ; tout effet de ponctuation doit être relié au ton de la scène, pas seulement nommé.
\end{corrige}

\newpage

\coursec{Fiche bilan}

\begin{popfigure}
\begin{tikzpicture}[mindmap, grow cyclic, every node/.style={concept}, root concept/.append style={concept color=popblue, font=\small\bfseries, minimum size=2.6cm, text=white}, level 1 concept/.append style={sibling angle=60, level distance=3.6cm, font=\scriptsize, minimum size=1.9cm, text=white}]
\node[root concept]{Produire un plan détaillé de dissertation}
  child[concept color=poporange]{ node {Paratextes et lecture ouvroir : titre, couverture, 4\textsuperscript{e} de couverture} }
  child[concept color=popgreen]{ node {Analyse du sujet : mots-clés, problématique} }
  child[concept color=poppurple]{ node {Recherche des idées : remue-méninges, classement} }
  child[concept color=poppink]{ node {Types de plans : dialectique, thématique, analytique} }
  child[concept color=popgold]{ node {Plan détaillé : thèse, arguments, exemples} }
  child[concept color=popblue!70]{ node {Communication : émetteur, récepteur, message, canal, code, contexte} };
\end{tikzpicture}
\legende{Carte mentale de la leçon : les six savoirs essentiels autour du plan détaillé de dissertation.}
\end{popfigure}

\begin{bilan}
\textbf{Définitions clés}
\begin{itemize}
  \item \cle{Paratexte} : ensemble des éléments qui entourent le texte (titre, nom de l'auteur, couverture, quatrième de couverture, préface).
  \item \cle{Lecture ouvroir} : lecture du début de l'œuvre (l'« ouvroir »), première approche qui ouvre le texte.
  \item \cle{Problématique} : question centrale, issue de l'analyse du sujet, qui oriente toute la réflexion.
  \item \cle{Plan détaillé} : plan où chaque partie est développée en sous-parties avec thèse, arguments et exemples précis.
  \item \cle{Ponctuation forte} : signes qui marquent une pause nette ou une fin de phrase : le point (.), le point d'interrogation (?), le point d'exclamation (!), les points de suspension (...).
\end{itemize}

\textbf{Les trois types de plans à connaître par cœur}
\begin{center}
\begin{tabularx}{\linewidth}{|l|X|X|}
\hline
\rowcolor{popblueL} \textbf{Type de plan} & \textbf{Structure} & \textbf{Usage} \\
\hline
Dialectique & Thèse / Antithèse / Synthèse & Sujets polémiques, débats d'idées \\
\hline
Thématique & Plusieurs aspects d'un même thème & Sujets descriptifs ou explicatifs \\
\hline
Analytique & Constat / Causes / Conséquences (ou solutions) & Sujets de réflexion sur un problème \\
\hline
\end{tabularx}
\end{center}

\textbf{Les six éléments de la communication}
\begin{center}
\begin{tabularx}{\linewidth}{|l|X|}
\hline
\rowcolor{popgreenL} \textbf{Élément} & \textbf{Rôle} \\
\hline
Émetteur & celui qui envoie le message \\
\hline
Récepteur & celui qui reçoit le message \\
\hline
Message & contenu transmis \\
\hline
Canal & moyen matériel de transmission (voix, écrit, téléphone) \\
\hline
Code & langue ou système de signes partagé \\
\hline
Contexte & situation dans laquelle on communique \\
\hline
\end{tabularx}
\end{center}

\textbf{Méthode express : construire un plan détaillé}
\begin{enumerate}
  \item Analyser le sujet : souligner les mots-clés, définir les termes, repérer les limites.
  \item Formuler la problématique sous forme d'une question.
  \item Faire un remue-méninges : noter toutes les idées, arguments et exemples.
  \item Classer les idées et choisir le type de plan adapté.
  \item Rédiger le plan : 2 ou 3 grandes parties, chacune avec sous-parties, arguments et exemples (tirés notamment du \emph{Lion et la Perle} de Wole Soyinka).
\end{enumerate}
\end{bilan}

\begin{aretenir}[Checklist avant le Bac]
Avant de produire ton plan détaillé, vérifie :
\begin{itemize}
  \item $\square$ J'ai lu et compris le sujet en entier, sans rien ajouter ni oublier.
  \item $\square$ J'ai souligné et défini les mots-clés du sujet.
  \item $\square$ Ma problématique est formulée sous forme de question claire.
  \item $\square$ J'ai noté assez d'idées, d'arguments et d'exemples avant de planifier.
  \item $\square$ J'ai choisi le type de plan adapté au sujet (dialectique, thématique ou analytique).
  \item $\square$ Mon plan comporte une introduction, 2 ou 3 parties équilibrées et une conclusion.
  \item $\square$ Chaque partie contient une thèse, des arguments et au moins un exemple précis.
  \item $\square$ Mes exemples sont pertinents (œuvres lues, vie courante, société camerounaise).
  \item $\square$ J'emploie correctement la ponctuation forte : point, ?, !, points de suspension.
  \item $\square$ Je respecte les conditions de communication : message clair, code adapté au correcteur.
  \item $\square$ Je me relis pour corriger les fautes et vérifier la logique du plan.
\end{itemize}
\end{aretenir}

\findecours

\end{document}
